\documentclass[11pt,a4paper]{article}
\usepackage[a4paper,margin=20mm]{geometry}
\usepackage{fontspec}
\setmainfont{Noto Serif}
\setsansfont{Noto Sans}
\usepackage{amsmath,amssymb}
\usepackage{booktabs,tabularx,array}
\usepackage[table]{xcolor}
\usepackage{hyperref}
\definecolor{primary}{RGB}{27,79,114}
\definecolor{accent}{RGB}{40,116,166}
\definecolor{rowtint}{RGB}{235,245,251}
\hypersetup{colorlinks=true,linkcolor=primary,urlcolor=primary}
\setlength{\parskip}{5pt plus 2pt minus 1pt}
\setlength{\parindent}{0pt}
\begin{document}
\begin{center}
  {\LARGE\bfseries\color{primary} PHIẾU HỌC TẬP TRẮC NGHIỆM TOÁN 10}\\[4pt]
  {\normalsize\itshape\color{accent} Bộ sách: Kết nối tri thức với cuộc sống (Tập 1 \& Tập 2 – 9 Chương)}\\[4pt]
  {\small Cấu trúc mỗi chương: 20 câu Dễ (Câu 1--20) $\bullet$ 20 câu Trung bình (Câu 21--40) $\bullet$ 10 câu Khó (Câu 41--50) -- Tổng 50 câu/chương}
\end{center}
\vspace{6pt}\hrule\vspace{10pt}
\section*{\color{primary} CHƯƠNG I: MỆNH ĐỀ VÀ TẬP HỢP}
\textit{Trọng tâm kiến thức SGK: Mệnh đề, mệnh đề chứa biến, mệnh đề phủ định, mệnh đề kéo theo, mệnh đề tương đương, tập hợp và các phép toán trên tập hợp (giao, hợp, hiệu, phần bù).}\par\vspace{4pt}
\subsection*{\color{accent} Phần I. Mức độ Dễ (Nhận biết -- Câu 1 đến Câu 20)}
\textbf{\color{primary} Câu 1.} Mệnh đề phủ định của mệnh đề P: '$\forall$x $\in$ $\mathbb{R}$, x$^2$ + 1 > 0' là:\par
\hspace*{1.5em}\textbf{A.} $\bar{P}$: '$\forall$x $\in$ $\mathbb{R}$, x$^2$ + 1 $\le$ 0'\quad\quad \textbf{B.} $\bar{P}$: '$\exists$x $\in$ $\mathbb{R}$, x$^2$ + 1 < 0'\quad\quad \textbf{C.} $\bar{P}$: '$\forall$x $\in$ $\mathbb{R}$, x$^2$ + 1 < 0'\quad\quad \textbf{D.} $\bar{P}$: '$\exists$x $\in$ $\mathbb{R}$, x$^2$ + 1 $\le$ 0'\par\vspace{2pt}
\textbf{\color{primary} Câu 2.} Thực hiện yêu cầu nhận biết sau: Cho hai tập hợp A = \{1; 2; 3; 5\} và B = \{2; 5; 7\}. Tập hợp giao A $\cap$ B là.\par
\hspace*{1.5em}\textbf{A.} \{7\}\quad\quad \textbf{B.} \{2; 5\}\quad\quad \textbf{C.} \{1; 2; 3; 5; 7\}\quad\quad \textbf{D.} \{1; 3\}\par\vspace{2pt}
\textbf{\color{primary} Câu 3.} Trong bài học thuộc Chương I, xét nội dung: Một tập hợp A có 5 phần tử. Số tập hợp con của tập hợp A là. Đáp án đúng là:\par
\hspace*{1.5em}\textbf{A.} 32 tập con\quad\quad \textbf{B.} 31 tập con\quad\quad \textbf{C.} 10 tập con\quad\quad \textbf{D.} 25 tập con\par\vspace{2pt}
\textbf{\color{primary} Câu 4.} Dựa vào kiến thức chuẩn SGK của Chương I (Mệnh đề và tập hợp), mệnh đề phủ định của mệnh đề P: '$\forall$x $\in$ $\mathbb{R}$, x$^2$ + 4 > 0' là:\par
\hspace*{1.5em}\textbf{A.} $\bar{P}$: '$\exists$x $\in$ $\mathbb{R}$, x$^2$ + 4 $\le$ 0'\quad\quad \textbf{B.} $\bar{P}$: '$\forall$x $\in$ $\mathbb{R}$, x$^2$ + 4 $\le$ 0'\quad\quad \textbf{C.} $\bar{P}$: '$\forall$x $\in$ $\mathbb{R}$, x$^2$ + 4 < 0'\quad\quad \textbf{D.} $\bar{P}$: '$\exists$x $\in$ $\mathbb{R}$, x$^2$ + 4 < 0'\par\vspace{2pt}
\textbf{\color{primary} Câu 5.} Chọn kết quả đúng nhất khi xét: Cho hai tập hợp A = \{1; 2; 3; 8\} và B = \{2; 8; 10\}. Tập hợp giao A $\cap$ B là:\par
\hspace*{1.5em}\textbf{A.} \{1; 3\}\quad\quad \textbf{B.} \{1; 2; 3; 8; 10\}\quad\quad \textbf{C.} \{2; 8\}\quad\quad \textbf{D.} \{10\}\par\vspace{2pt}
\textbf{\color{primary} Câu 6.} Khi làm bài tập cơ bản về mệnh đề và tập hợp, một học sinh gặp câu hỏi: Một tập hợp A có 3 phần tử. Số tập hợp con của tập hợp A là. Kết quả nào dưới đây đúng?\par
\hspace*{1.5em}\textbf{A.} 8 tập con\quad\quad \textbf{B.} 6 tập con\quad\quad \textbf{C.} 7 tập con\quad\quad \textbf{D.} 9 tập con\par\vspace{2pt}
\textbf{\color{primary} Câu 7.} Xác định giá trị hoặc khẳng định đúng cho bài toán sau: Mệnh đề phủ định của mệnh đề P: '$\forall$x $\in$ $\mathbb{R}$, x$^2$ + 7 > 0' là:\par
\hspace*{1.5em}\textbf{A.} $\bar{P}$: '$\exists$x $\in$ $\mathbb{R}$, x$^2$ + 7 $\le$ 0'\quad\quad \textbf{B.} $\bar{P}$: '$\forall$x $\in$ $\mathbb{R}$, x$^2$ + 7 < 0'\quad\quad \textbf{C.} $\bar{P}$: '$\exists$x $\in$ $\mathbb{R}$, x$^2$ + 7 < 0'\quad\quad \textbf{D.} $\bar{P}$: '$\forall$x $\in$ $\mathbb{R}$, x$^2$ + 7 $\le$ 0'\par\vspace{2pt}
\textbf{\color{primary} Câu 8.} Áp dụng trực tiếp định nghĩa và quy tắc trong Chương I: Cho hai tập hợp A = \{1; 2; 3; 11\} và B = \{2; 11; 13\}. Tập hợp giao A $\cap$ B là:\par
\hspace*{1.5em}\textbf{A.} \{1; 3\}\quad\quad \textbf{B.} \{13\}\quad\quad \textbf{C.} \{2; 11\}\quad\quad \textbf{D.} \{1; 2; 3; 11; 13\}\par\vspace{2pt}
\textbf{\color{primary} Câu 9.} Kiểm tra nhanh kiến thức nền tảng: Một tập hợp A có 6 phần tử. Số tập hợp con của tập hợp A là:\par
\hspace*{1.5em}\textbf{A.} 36 tập con\quad\quad \textbf{B.} 12 tập con\quad\quad \textbf{C.} 63 tập con\quad\quad \textbf{D.} 64 tập con\par\vspace{2pt}
\textbf{\color{primary} Câu 10.} Cho dữ kiện nhận biết sau đây: Mệnh đề phủ định của mệnh đề P: '$\forall$x $\in$ $\mathbb{R}$, x$^2$ + 10 > 0' là. Phương án chính xác là:\par
\begin{itemize}\setlength{\itemsep}{1pt}\setlength{\parskip}{0pt}
  \item[] \textbf{A.} $\bar{P}$: '$\forall$x $\in$ $\mathbb{R}$, x$^2$ + 10 < 0'
  \item[] \textbf{B.} $\bar{P}$: '$\exists$x $\in$ $\mathbb{R}$, x$^2$ + 10 < 0'
  \item[] \textbf{C.} $\bar{P}$: '$\forall$x $\in$ $\mathbb{R}$, x$^2$ + 10 $\le$ 0'
  \item[] \textbf{D.} $\bar{P}$: '$\exists$x $\in$ $\mathbb{R}$, x$^2$ + 10 $\le$ 0'
\end{itemize}\vspace{2pt}
\textbf{\color{primary} Câu 11.} Hãy chỉ ra kết quả đúng trong bốn lựa chọn A, B, C, D đối với yêu cầu: Cho hai tập hợp A = \{1; 2; 3; 14\} và B = \{2; 14; 16\}. Tập hợp giao A $\cap$ B là:\par
\hspace*{1.5em}\textbf{A.} \{1; 2; 3; 14; 16\}\quad\quad \textbf{B.} \{1; 3\}\quad\quad \textbf{C.} \{16\}\quad\quad \textbf{D.} \{2; 14\}\par\vspace{2pt}
\textbf{\color{primary} Câu 12.} Câu hỏi ôn tập lý thuyết và kỹ năng cơ bản (Chương I): Một tập hợp A có 4 phần tử. Số tập hợp con của tập hợp A là:\par
\hspace*{1.5em}\textbf{A.} 16 tập con\quad\quad \textbf{B.} 8 tập con\quad\quad \textbf{C.} 16 tập con (1)\quad\quad \textbf{D.} 15 tập con\par\vspace{2pt}
\textbf{\color{primary} Câu 13.} Theo nội dung sách giáo khoa Kết nối tri thức với cuộc sống (Chương I), mệnh đề phủ định của mệnh đề P: '$\forall$x $\in$ $\mathbb{R}$, x$^2$ + 13 > 0' là:\par
\begin{itemize}\setlength{\itemsep}{1pt}\setlength{\parskip}{0pt}
  \item[] \textbf{A.} $\bar{P}$: '$\exists$x $\in$ $\mathbb{R}$, x$^2$ + 13 < 0'
  \item[] \textbf{B.} $\bar{P}$: '$\exists$x $\in$ $\mathbb{R}$, x$^2$ + 13 $\le$ 0'
  \item[] \textbf{C.} $\bar{P}$: '$\forall$x $\in$ $\mathbb{R}$, x$^2$ + 13 $\le$ 0'
  \item[] \textbf{D.} $\bar{P}$: '$\forall$x $\in$ $\mathbb{R}$, x$^2$ + 13 < 0'
\end{itemize}\vspace{2pt}
\textbf{\color{primary} Câu 14.} Tìm đáp số đúng của bài toán cơ bản sau: Cho hai tập hợp A = \{1; 2; 3; 17\} và B = \{2; 17; 19\}. Tập hợp giao A $\cap$ B là:\par
\hspace*{1.5em}\textbf{A.} \{2; 17\}\quad\quad \textbf{B.} \{19\}\quad\quad \textbf{C.} \{1; 2; 3; 17; 19\}\quad\quad \textbf{D.} \{1; 3\}\par\vspace{2pt}
\textbf{\color{primary} Câu 15.} Xét đối tượng toán học trong Chương I: Một tập hợp A có 7 phần tử. Số tập hợp con của tập hợp A là. Giá trị hoặc tính chất tương ứng là:\par
\hspace*{1.5em}\textbf{A.} 127 tập con\quad\quad \textbf{B.} 14 tập con\quad\quad \textbf{C.} 49 tập con\quad\quad \textbf{D.} 128 tập con\par\vspace{2pt}
\textbf{\color{primary} Câu 16.} Hoàn thành phép xác định cơ bản sau đây: Mệnh đề phủ định của mệnh đề P: '$\forall$x $\in$ $\mathbb{R}$, x$^2$ + 16 > 0' là:\par
\begin{itemize}\setlength{\itemsep}{1pt}\setlength{\parskip}{0pt}
  \item[] \textbf{A.} $\bar{P}$: '$\exists$x $\in$ $\mathbb{R}$, x$^2$ + 16 $\le$ 0'
  \item[] \textbf{B.} $\bar{P}$: '$\forall$x $\in$ $\mathbb{R}$, x$^2$ + 16 < 0'
  \item[] \textbf{C.} $\bar{P}$: '$\exists$x $\in$ $\mathbb{R}$, x$^2$ + 16 < 0'
  \item[] \textbf{D.} $\bar{P}$: '$\forall$x $\in$ $\mathbb{R}$, x$^2$ + 16 $\le$ 0'
\end{itemize}\vspace{2pt}
\textbf{\color{primary} Câu 17.} Trong các khẳng định hoặc giá trị dưới đây, đâu là đáp án đúng cho câu hỏi: Cho hai tập hợp A = \{1; 2; 3; 20\} và B = \{2; 20; 22\}. Tập hợp giao A $\cap$ B là?\par
\hspace*{1.5em}\textbf{A.} \{1; 2; 3; 20; 22\}\quad\quad \textbf{B.} \{22\}\quad\quad \textbf{C.} \{2; 20\}\quad\quad \textbf{D.} \{1; 3\}\par\vspace{2pt}
\textbf{\color{primary} Câu 18.} Bài tập khởi động Chương I: Một tập hợp A có 5 phần tử. Số tập hợp con của tập hợp A là:\par
\hspace*{1.5em}\textbf{A.} 25 tập con\quad\quad \textbf{B.} 10 tập con\quad\quad \textbf{C.} 32 tập con\quad\quad \textbf{D.} 31 tập con\par\vspace{2pt}
\textbf{\color{primary} Câu 19.} Quan sát dữ kiện đề bài: Mệnh đề phủ định của mệnh đề P: '$\forall$x $\in$ $\mathbb{R}$, x$^2$ + 19 > 0' là. Hãy chọn phương án đúng:\par
\begin{itemize}\setlength{\itemsep}{1pt}\setlength{\parskip}{0pt}
  \item[] \textbf{A.} $\bar{P}$: '$\forall$x $\in$ $\mathbb{R}$, x$^2$ + 19 $\le$ 0'
  \item[] \textbf{B.} $\bar{P}$: '$\forall$x $\in$ $\mathbb{R}$, x$^2$ + 19 < 0'
  \item[] \textbf{C.} $\bar{P}$: '$\exists$x $\in$ $\mathbb{R}$, x$^2$ + 19 < 0'
  \item[] \textbf{D.} $\bar{P}$: '$\exists$x $\in$ $\mathbb{R}$, x$^2$ + 19 $\le$ 0'
\end{itemize}\vspace{2pt}
\textbf{\color{primary} Câu 20.} Sử dụng công thức hoặc tính chất đầu bài của Chương I để giải quyết yêu cầu: Cho hai tập hợp A = \{1; 2; 3; 23\} và B = \{2; 23; 25\}. Tập hợp giao A $\cap$ B là:\par
\hspace*{1.5em}\textbf{A.} \{25\}\quad\quad \textbf{B.} \{1; 3\}\quad\quad \textbf{C.} \{1; 2; 3; 23; 25\}\quad\quad \textbf{D.} \{2; 23\}\par\vspace{2pt}
\subsection*{\color{accent} Phần II. Mức độ Trung bình (Thông hiểu -- Câu 21 đến Câu 40)}
\textbf{\color{primary} Câu 21.} Cho hai tập hợp A = (-3; 4] và B = (1; 7). Tập hợp A $\cap$ B bằng:\par
\hspace*{1.5em}\textbf{A.} (-3; 7)\quad\quad \textbf{B.} (1; 4]\quad\quad \textbf{C.} [-3; 1]\quad\quad \textbf{D.} (4; 7)\par\vspace{2pt}
\textbf{\color{primary} Câu 22.} Vận dụng quy tắc trong Chương I (Mệnh đề và tập hợp) để giải bài toán: Cho hai tập hợp A = [-2; 4) và B = [0; 7]. Tập hợp hiệu A \textbackslash\{\} B bằng:\par
\hspace*{1.5em}\textbf{A.} (-2; 4)\quad\quad \textbf{B.} [4; 7]\quad\quad \textbf{C.} [-2; 0)\quad\quad \textbf{D.} [-2; 0]\par\vspace{2pt}
\textbf{\color{primary} Câu 23.} Biến đổi và tính toán theo kiến thức Chương I: Liệt kê các phần tử của tập hợp M = \{x $\in$ $\mathbb{Z}$ | x$^2$ - 5x + 6 = 0\} ta được. Kết quả thu được là:\par
\hspace*{1.5em}\textbf{A.} \{1; 6\}\quad\quad \textbf{B.} \{2; 3\}\quad\quad \textbf{C.} \{-2; -3\}\quad\quad \textbf{D.} $\emptyset$\par\vspace{2pt}
\textbf{\color{primary} Câu 24.} Xét bài toán mức độ thông hiểu sau: Cho hai tập hợp A = (-3; 7] và B = (1; 10). Tập hợp A $\cap$ B bằng. Đáp án đúng là:\par
\hspace*{1.5em}\textbf{A.} (7; 10)\quad\quad \textbf{B.} (1; 7]\quad\quad \textbf{C.} [-3; 1]\quad\quad \textbf{D.} (-3; 10)\par\vspace{2pt}
\textbf{\color{primary} Câu 25.} Để hoàn thành phiếu học tập Chương I, hãy xác định: Cho hai tập hợp A = [-2; 7) và B = [0; 10]. Tập hợp hiệu A \textbackslash\{\} B bằng:\par
\hspace*{1.5em}\textbf{A.} [7; 10]\quad\quad \textbf{B.} [-2; 0)\quad\quad \textbf{C.} (-2; 7)\quad\quad \textbf{D.} [-2; 0]\par\vspace{2pt}
\textbf{\color{primary} Câu 26.} Phân tích dữ kiện bài toán: Liệt kê các phần tử của tập hợp M = \{x $\in$ $\mathbb{Z}$ | x$^2$ - 8x + 12 = 0\} ta được. Giá trị cần tìm bằng:\par
\hspace*{1.5em}\textbf{A.} $\emptyset$\quad\quad \textbf{B.} \{1; 12\}\quad\quad \textbf{C.} \{-2; -6\}\quad\quad \textbf{D.} \{2; 6\}\par\vspace{2pt}
\textbf{\color{primary} Câu 27.} Áp dụng tính chất đã học ở Chương I vào bài tập sau: Cho hai tập hợp A = (-3; 10] và B = (1; 13). Tập hợp A $\cap$ B bằng:\par
\hspace*{1.5em}\textbf{A.} (1; 10]\quad\quad \textbf{B.} (-3; 13)\quad\quad \textbf{C.} (10; 13)\quad\quad \textbf{D.} [-3; 1]\par\vspace{2pt}
\textbf{\color{primary} Câu 28.} Thực hiện các bước suy luận hợp lý để giải quyết yêu cầu: Cho hai tập hợp A = [-2; 10) và B = [0; 13]. Tập hợp hiệu A \textbackslash\{\} B bằng:\par
\hspace*{1.5em}\textbf{A.} [10; 13]\quad\quad \textbf{B.} [-2; 0)\quad\quad \textbf{C.} (-2; 10)\quad\quad \textbf{D.} [-2; 0]\par\vspace{2pt}
\textbf{\color{primary} Câu 29.} Cho bài toán tự luyện (Mệnh đề và tập hợp): Liệt kê các phần tử của tập hợp M = \{x $\in$ $\mathbb{Z}$ | x$^2$ - 11x + 18 = 0\} ta được. Phương án nào sau đây chính xác?\par
\hspace*{1.5em}\textbf{A.} \{-2; -9\}\quad\quad \textbf{B.} \{2; 9\}\quad\quad \textbf{C.} \{1; 18\}\quad\quad \textbf{D.} $\emptyset$\par\vspace{2pt}
\textbf{\color{primary} Câu 30.} Kết hợp các phép biến đổi trong Chương I, hãy tìm kết quả của bài toán: Cho hai tập hợp A = (-3; 13] và B = (1; 16). Tập hợp A $\cap$ B bằng:\par
\hspace*{1.5em}\textbf{A.} [-3; 1]\quad\quad \textbf{B.} (1; 13]\quad\quad \textbf{C.} (13; 16)\quad\quad \textbf{D.} (-3; 16)\par\vspace{2pt}
\textbf{\color{primary} Câu 31.} Một bài kiểm tra định kỳ nội dung Chương I nêu câu hỏi: Cho hai tập hợp A = [-2; 13) và B = [0; 16]. Tập hợp hiệu A \textbackslash\{\} B bằng. Hãy chọn đáp số đúng:\par
\hspace*{1.5em}\textbf{A.} [13; 16]\quad\quad \textbf{B.} [-2; 0]\quad\quad \textbf{C.} (-2; 13)\quad\quad \textbf{D.} [-2; 0)\par\vspace{2pt}
\textbf{\color{primary} Câu 32.} Tìm lời giải đúng cho bài toán thông hiểu: Liệt kê các phần tử của tập hợp M = \{x $\in$ $\mathbb{Z}$ | x$^2$ - 14x + 24 = 0\} ta được:\par
\hspace*{1.5em}\textbf{A.} \{-2; -12\}\quad\quad \textbf{B.} $\emptyset$\quad\quad \textbf{C.} \{2; 12\}\quad\quad \textbf{D.} \{1; 24\}\par\vspace{2pt}
\textbf{\color{primary} Câu 33.} Sử dụng mối liên hệ giữa các đại lượng trong Chương I: Cho hai tập hợp A = (-3; 16] và B = (1; 19). Tập hợp A $\cap$ B bằng:\par
\hspace*{1.5em}\textbf{A.} (1; 16]\quad\quad \textbf{B.} (16; 19)\quad\quad \textbf{C.} [-3; 1]\quad\quad \textbf{D.} (-3; 19)\par\vspace{2pt}
\textbf{\color{primary} Câu 34.} Hãy tính toán và chọn phương án đúng đối với bài tập: Cho hai tập hợp A = [-2; 16) và B = [0; 19]. Tập hợp hiệu A \textbackslash\{\} B bằng:\par
\hspace*{1.5em}\textbf{A.} (-2; 16)\quad\quad \textbf{B.} [-2; 0)\quad\quad \textbf{C.} [-2; 0]\quad\quad \textbf{D.} [16; 19]\par\vspace{2pt}
\textbf{\color{primary} Câu 35.} Dựa trên các tính chất trọng tâm của Chương I, xét bài toán: Liệt kê các phần tử của tập hợp M = \{x $\in$ $\mathbb{Z}$ | x$^2$ - 17x + 30 = 0\} ta được:\par
\hspace*{1.5em}\textbf{A.} \{2; 15\}\quad\quad \textbf{B.} \{-2; -15\}\quad\quad \textbf{C.} $\emptyset$\quad\quad \textbf{D.} \{1; 30\}\par\vspace{2pt}
\textbf{\color{primary} Câu 36.} Thực hành giải toán (Mệnh đề và tập hợp): Cho hai tập hợp A = (-3; 19] và B = (1; 22). Tập hợp A $\cap$ B bằng. Đáp án nào dưới đây thỏa mãn?\par
\hspace*{1.5em}\textbf{A.} (1; 19]\quad\quad \textbf{B.} (19; 22)\quad\quad \textbf{C.} [-3; 1]\quad\quad \textbf{D.} (-3; 22)\par\vspace{2pt}
\textbf{\color{primary} Câu 37.} Kiểm tra kỹ năng biến đổi và suy luận qua bài toán: Cho hai tập hợp A = [-2; 19) và B = [0; 22]. Tập hợp hiệu A \textbackslash\{\} B bằng:\par
\hspace*{1.5em}\textbf{A.} [-2; 0)\quad\quad \textbf{B.} [19; 22]\quad\quad \textbf{C.} (-2; 19)\quad\quad \textbf{D.} [-2; 0]\par\vspace{2pt}
\textbf{\color{primary} Câu 38.} Cho giả thiết thuộc phạm vi Chương I: Liệt kê các phần tử của tập hợp M = \{x $\in$ $\mathbb{Z}$ | x$^2$ - 20x + 36 = 0\} ta được. Kết luận đúng là:\par
\hspace*{1.5em}\textbf{A.} \{2; 18\}\quad\quad \textbf{B.} \{-2; -18\}\quad\quad \textbf{C.} $\emptyset$\quad\quad \textbf{D.} \{1; 36\}\par\vspace{2pt}
\textbf{\color{primary} Câu 39.} Rút gọn hoặc tính giá trị theo yêu cầu sau của Chương I: Cho hai tập hợp A = (-3; 22] và B = (1; 25). Tập hợp A $\cap$ B bằng:\par
\hspace*{1.5em}\textbf{A.} (-3; 25)\quad\quad \textbf{B.} [-3; 1]\quad\quad \textbf{C.} (1; 22]\quad\quad \textbf{D.} (22; 25)\par\vspace{2pt}
\textbf{\color{primary} Câu 40.} Xác định phương án đúng sau khi thực hiện phép tính / suy luận: Cho hai tập hợp A = [-2; 22) và B = [0; 25]. Tập hợp hiệu A \textbackslash\{\} B bằng:\par
\hspace*{1.5em}\textbf{A.} (-2; 22)\quad\quad \textbf{B.} [-2; 0]\quad\quad \textbf{C.} [-2; 0)\quad\quad \textbf{D.} [22; 25]\par\vspace{2pt}
\subsection*{\color{accent} Phần III. Mức độ Khó (Vận dụng \& Vận dụng cao -- Câu 41 đến Câu 50)}
\textbf{\color{primary} Câu 41.} Lớp 10A có 40 học sinh, trong đó có 25 bạn thích môn Toán, 20 bạn thích môn Anh và 10 bạn thích cả hai môn Toán và Anh. Số học sinh không thích cả hai môn đó là:\par
\hspace*{1.5em}\textbf{A.} 5 học sinh\quad\quad \textbf{B.} 10 học sinh\quad\quad \textbf{C.} 0 học sinh\quad\quad \textbf{D.} 15 học sinh\par\vspace{2pt}
\textbf{\color{primary} Câu 42.} Bài toán vận dụng (Chương I – Mệnh đề và tập hợp): Tìm điều kiện của tham số m để hai khoảng A = (-$\infty$; m) và B = (3; +$\infty$) có giao khác rỗng (A $\cap$ B $\ne$ $\emptyset$). Kết quả đúng là:\par
\hspace*{1.5em}\textbf{A.} m $\le$ 3\quad\quad \textbf{B.} m > 3\quad\quad \textbf{C.} m $\ge$ 3\quad\quad \textbf{D.} m < 3\par\vspace{2pt}
\textbf{\color{primary} Câu 43.} Vận dụng kiến thức nâng cao trong phạm vi Chương I để giải bài toán: Lớp 10A có 42 học sinh, trong đó có 27 bạn thích môn Toán, 20 bạn thích môn Anh và 10 bạn thích cả hai môn Toán và Anh. Số học sinh không thích cả hai môn đó là:\par
\hspace*{1.5em}\textbf{A.} 10 học sinh\quad\quad \textbf{B.} 15 học sinh\quad\quad \textbf{C.} 5 học sinh\quad\quad \textbf{D.} 0 học sinh\par\vspace{2pt}
\textbf{\color{primary} Câu 44.} Phân tích và giải bài toán tổng hợp sau của Chương I: Tìm điều kiện của tham số m để hai khoảng A = (-$\infty$; m) và B = (5; +$\infty$) có giao khác rỗng (A $\cap$ B $\ne$ $\emptyset$):\par
\hspace*{1.5em}\textbf{A.} m > 5\quad\quad \textbf{B.} m $\le$ 5\quad\quad \textbf{C.} m $\ge$ 5\quad\quad \textbf{D.} m < 5\par\vspace{2pt}
\textbf{\color{primary} Câu 45.} Thử thách tư duy trong phạm vi kiến thức Chương I: Lớp 10A có 44 học sinh, trong đó có 29 bạn thích môn Toán, 20 bạn thích môn Anh và 10 bạn thích cả hai môn Toán và Anh. Số học sinh không thích cả hai môn đó là. Đáp số chính xác là:\par
\hspace*{1.5em}\textbf{A.} 5 học sinh\quad\quad \textbf{B.} 10 học sinh\quad\quad \textbf{C.} 15 học sinh\quad\quad \textbf{D.} 0 học sinh\par\vspace{2pt}
\textbf{\color{primary} Câu 46.} Sử dụng phương pháp giải toán của Chương I để tìm đáp án cho bài toán: Tìm điều kiện của tham số m để hai khoảng A = (-$\infty$; m) và B = (7; +$\infty$) có giao khác rỗng (A $\cap$ B $\ne$ $\emptyset$):\par
\hspace*{1.5em}\textbf{A.} m $\le$ 7\quad\quad \textbf{B.} m $\ge$ 7\quad\quad \textbf{C.} m < 7\quad\quad \textbf{D.} m > 7\par\vspace{2pt}
\textbf{\color{primary} Câu 47.} Xét bài toán vận dụng thực tiễn / nâng cao thuộc Chương I: Lớp 10A có 46 học sinh, trong đó có 31 bạn thích môn Toán, 20 bạn thích môn Anh và 10 bạn thích cả hai môn Toán và Anh. Số học sinh không thích cả hai môn đó là:\par
\hspace*{1.5em}\textbf{A.} 5 học sinh\quad\quad \textbf{B.} 0 học sinh\quad\quad \textbf{C.} 15 học sinh\quad\quad \textbf{D.} 10 học sinh\par\vspace{2pt}
\textbf{\color{primary} Câu 48.} Tìm kết quả cuối cùng của bài toán phân loại học sinh khá giỏi (Chương I): Tìm điều kiện của tham số m để hai khoảng A = (-$\infty$; m) và B = (9; +$\infty$) có giao khác rỗng (A $\cap$ B $\ne$ $\emptyset$):\par
\hspace*{1.5em}\textbf{A.} m $\le$ 9\quad\quad \textbf{B.} m > 9\quad\quad \textbf{C.} m $\ge$ 9\quad\quad \textbf{D.} m < 9\par\vspace{2pt}
\textbf{\color{primary} Câu 49.} Kết hợp các tính chất đã học trong Chương I, hãy giải quyết bài toán: Lớp 10A có 48 học sinh, trong đó có 33 bạn thích môn Toán, 20 bạn thích môn Anh và 10 bạn thích cả hai môn Toán và Anh. Số học sinh không thích cả hai môn đó là:\par
\hspace*{1.5em}\textbf{A.} 0 học sinh\quad\quad \textbf{B.} 15 học sinh\quad\quad \textbf{C.} 10 học sinh\quad\quad \textbf{D.} 5 học sinh\par\vspace{2pt}
\textbf{\color{primary} Câu 50.} Câu hỏi vận dụng cao cuối Chương I: Tìm điều kiện của tham số m để hai khoảng A = (-$\infty$; m) và B = (11; +$\infty$) có giao khác rỗng (A $\cap$ B $\ne$ $\emptyset$). Phương án đúng là:\par
\hspace*{1.5em}\textbf{A.} m > 11\quad\quad \textbf{B.} m < 11\quad\quad \textbf{C.} m $\ge$ 11\quad\quad \textbf{D.} m $\le$ 11\par\vspace{2pt}
\subsection*{\color{primary} Bảng đáp án Chương I: Mệnh đề và tập hợp}
\begin{center}
\renewcommand{\arraystretch}{1.25}
\begin{tabularx}{\textwidth}{*{10}{>{\centering\arraybackslash}X}}
\toprule
\textbf{\color{primary} Cột 1} & \textbf{\color{primary} Cột 2} & \textbf{\color{primary} Cột 3} & \textbf{\color{primary} Cột 4} & \textbf{\color{primary} Cột 5} & \textbf{\color{primary} Cột 6} & \textbf{\color{primary} Cột 7} & \textbf{\color{primary} Cột 8} & \textbf{\color{primary} Cột 9} & \textbf{\color{primary} Cột 10} \\
\midrule
\textbf{1.D} & \textbf{2.B} & \textbf{3.A} & \textbf{4.A} & \textbf{5.C} & \textbf{6.A} & \textbf{7.A} & \textbf{8.C} & \textbf{9.D} & \textbf{10.D} \\
\rowcolor{rowtint} \textbf{11.D} & \textbf{12.A} & \textbf{13.B} & \textbf{14.A} & \textbf{15.D} & \textbf{16.A} & \textbf{17.C} & \textbf{18.C} & \textbf{19.D} & \textbf{20.D} \\
\textbf{21.B} & \textbf{22.C} & \textbf{23.B} & \textbf{24.B} & \textbf{25.B} & \textbf{26.D} & \textbf{27.A} & \textbf{28.B} & \textbf{29.B} & \textbf{30.B} \\
\rowcolor{rowtint} \textbf{31.D} & \textbf{32.C} & \textbf{33.A} & \textbf{34.B} & \textbf{35.A} & \textbf{36.A} & \textbf{37.A} & \textbf{38.A} & \textbf{39.C} & \textbf{40.C} \\
\textbf{41.A} & \textbf{42.B} & \textbf{43.C} & \textbf{44.A} & \textbf{45.A} & \textbf{46.D} & \textbf{47.A} & \textbf{48.B} & \textbf{49.D} & \textbf{50.A} \\
\bottomrule
\end{tabularx}
\end{center}\vspace{10pt}
\section*{\color{primary} CHƯƠNG II: BẤT PHƯƠNG TRÌNH VÀ HỆ BẤT PHƯƠNG TRÌNH BẬC NHẤT HAI ẨN}
\textit{Trọng tâm kiến thức SGK: Bất phương trình bậc nhất hai ẩn, miền nghiệm của bất phương trình và hệ bất phương trình bậc nhất hai ẩn, bài toán quy hoạch tuyến tính hai biến.}\par\vspace{4pt}
\subsection*{\color{accent} Phần I. Mức độ Dễ (Nhận biết -- Câu 1 đến Câu 20)}
\textbf{\color{primary} Câu 1.} Cặp số (x; y) nào dưới đây thuộc miền nghiệm của bất phương trình 2x - y + 2 > 0?\par
\hspace*{1.5em}\textbf{A.} (0; 5)\quad\quad \textbf{B.} (0; 3)\quad\quad \textbf{C.} (-2; 1)\quad\quad \textbf{D.} (1; 0)\par\vspace{2pt}
\textbf{\color{primary} Câu 2.} Hãy trả lời câu hỏi cơ bản sau: Bất phương trình nào sau đây là bất phương trình bậc nhất hai ẩn x, y?\par
\hspace*{1.5em}\textbf{A.} x$^2$ + 2y > 0\quad\quad \textbf{B.} 3x - 3y $\le$ 7\quad\quad \textbf{C.} x + y + z < 5\quad\quad \textbf{D.} 2x - 3xy $\ge$ 4\par\vspace{2pt}
\textbf{\color{primary} Câu 3.} Trong bài học thuộc Chương II, xét nội dung: Điểm O(0; 0) thuộc miền nghiệm của bất phương trình nào sau đây. Đáp án đúng là:\par
\hspace*{1.5em}\textbf{A.} 2x + 3y - 6 < 0\quad\quad \textbf{B.} 2x + 3y + 6 < 0\quad\quad \textbf{C.} 3x + y $\ge$ 3\quad\quad \textbf{D.} x - y > 4\par\vspace{2pt}
\textbf{\color{primary} Câu 4.} Dựa vào kiến thức chuẩn SGK của Chương II (Bất phương trình và hệ bất phương trình bậc nhất hai ẩn), cặp số (x; y) nào dưới đây thuộc miền nghiệm của bất phương trình 2x - y + 5 > 0:\par
\hspace*{1.5em}\textbf{A.} (-5; 1)\quad\quad \textbf{B.} (1; 0)\quad\quad \textbf{C.} (0; 6)\quad\quad \textbf{D.} (0; 8)\par\vspace{2pt}
\textbf{\color{primary} Câu 5.} Chọn phương án trả lời đúng nhất: Bất phương trình nào sau đây là bất phương trình bậc nhất hai ẩn x, y?\par
\hspace*{1.5em}\textbf{A.} x$^2$ + 5y > 0\quad\quad \textbf{B.} 2x - 3xy $\ge$ 4\quad\quad \textbf{C.} 3x - 6y $\le$ 7\quad\quad \textbf{D.} x + y + z < 5\par\vspace{2pt}
\textbf{\color{primary} Câu 6.} Khi làm bài tập cơ bản về bất phương trình và hệ bất phương trình bậc nhất hai ẩn, một học sinh gặp câu hỏi: Điểm O(0; 0) thuộc miền nghiệm của bất phương trình nào sau đây. Kết quả nào dưới đây đúng?\par
\hspace*{1.5em}\textbf{A.} x - y > 7\quad\quad \textbf{B.} 2x + 3y - 9 < 0\quad\quad \textbf{C.} 2x + 3y + 9 < 0\quad\quad \textbf{D.} 3x + y $\ge$ 6\par\vspace{2pt}
\textbf{\color{primary} Câu 7.} Xác định giá trị hoặc khẳng định đúng cho bài toán sau: Cặp số (x; y) nào dưới đây thuộc miền nghiệm của bất phương trình 2x - y + 8 > 0:\par
\hspace*{1.5em}\textbf{A.} (0; 11)\quad\quad \textbf{B.} (-8; 1)\quad\quad \textbf{C.} (1; 0)\quad\quad \textbf{D.} (0; 9)\par\vspace{2pt}
\textbf{\color{primary} Câu 8.} Áp dụng trực tiếp định nghĩa và quy tắc trong Chương II: Bất phương trình nào sau đây là bất phương trình bậc nhất hai ẩn x, y:\par
\hspace*{1.5em}\textbf{A.} 3x - 9y $\le$ 7\quad\quad \textbf{B.} 2x - 3xy $\ge$ 4\quad\quad \textbf{C.} x + y + z < 5\quad\quad \textbf{D.} x$^2$ + 8y > 0\par\vspace{2pt}
\textbf{\color{primary} Câu 9.} Kiểm tra nhanh kiến thức nền tảng: Điểm O(0; 0) thuộc miền nghiệm của bất phương trình nào sau đây?\par
\hspace*{1.5em}\textbf{A.} 2x + 3y + 12 < 0\quad\quad \textbf{B.} 3x + y $\ge$ 9\quad\quad \textbf{C.} 2x + 3y - 12 < 0\quad\quad \textbf{D.} x - y > 10\par\vspace{2pt}
\textbf{\color{primary} Câu 10.} Cho dữ kiện nhận biết sau đây: Cặp số (x; y) nào dưới đây thuộc miền nghiệm của bất phương trình 2x - y + 11 > 0. Phương án chính xác là:\par
\hspace*{1.5em}\textbf{A.} (0; 12)\quad\quad \textbf{B.} (1; 0)\quad\quad \textbf{C.} (-11; 1)\quad\quad \textbf{D.} (0; 14)\par\vspace{2pt}
\textbf{\color{primary} Câu 11.} Hãy chỉ ra kết quả đúng trong bốn lựa chọn A, B, C, D đối với yêu cầu: Bất phương trình nào sau đây là bất phương trình bậc nhất hai ẩn x, y:\par
\hspace*{1.5em}\textbf{A.} 2x - 3xy $\ge$ 4\quad\quad \textbf{B.} x$^2$ + 11y > 0\quad\quad \textbf{C.} x + y + z < 5\quad\quad \textbf{D.} 3x - 12y $\le$ 7\par\vspace{2pt}
\textbf{\color{primary} Câu 12.} Câu hỏi ôn tập lý thuyết và kỹ năng cơ bản (Chương II): Điểm O(0; 0) thuộc miền nghiệm của bất phương trình nào sau đây:\par
\hspace*{1.5em}\textbf{A.} x - y > 13\quad\quad \textbf{B.} 3x + y $\ge$ 12\quad\quad \textbf{C.} 2x + 3y - 15 < 0\quad\quad \textbf{D.} 2x + 3y + 15 < 0\par\vspace{2pt}
\textbf{\color{primary} Câu 13.} Theo nội dung sách giáo khoa Kết nối tri thức với cuộc sống (Chương II), cặp số (x; y) nào dưới đây thuộc miền nghiệm của bất phương trình 2x - y + 14 > 0:\par
\hspace*{1.5em}\textbf{A.} (0; 17)\quad\quad \textbf{B.} (1; 0)\quad\quad \textbf{C.} (-14; 1)\quad\quad \textbf{D.} (0; 15)\par\vspace{2pt}
\textbf{\color{primary} Câu 14.} Tìm đáp số đúng của bài toán cơ bản sau: Bất phương trình nào sau đây là bất phương trình bậc nhất hai ẩn x, y:\par
\hspace*{1.5em}\textbf{A.} 3x - 15y $\le$ 7\quad\quad \textbf{B.} x + y + z < 5\quad\quad \textbf{C.} x$^2$ + 14y > 0\quad\quad \textbf{D.} 2x - 3xy $\ge$ 4\par\vspace{2pt}
\textbf{\color{primary} Câu 15.} Xét đối tượng toán học trong Chương II: Điểm O(0; 0) thuộc miền nghiệm của bất phương trình nào sau đây. Giá trị hoặc tính chất tương ứng là:\par
\hspace*{1.5em}\textbf{A.} 2x + 3y + 18 < 0\quad\quad \textbf{B.} 2x + 3y - 18 < 0\quad\quad \textbf{C.} x - y > 16\quad\quad \textbf{D.} 3x + y $\ge$ 15\par\vspace{2pt}
\textbf{\color{primary} Câu 16.} Hoàn thành phép xác định cơ bản sau đây: Cặp số (x; y) nào dưới đây thuộc miền nghiệm của bất phương trình 2x - y + 17 > 0:\par
\hspace*{1.5em}\textbf{A.} (1; 0)\quad\quad \textbf{B.} (-17; 1)\quad\quad \textbf{C.} (0; 18)\quad\quad \textbf{D.} (0; 20)\par\vspace{2pt}
\textbf{\color{primary} Câu 17.} Trong các khẳng định hoặc giá trị dưới đây, đâu là đáp án đúng cho câu hỏi: Bất phương trình nào sau đây là bất phương trình bậc nhất hai ẩn x, y?\par
\hspace*{1.5em}\textbf{A.} x + y + z < 5\quad\quad \textbf{B.} 2x - 3xy $\ge$ 4\quad\quad \textbf{C.} 3x - 18y $\le$ 7\quad\quad \textbf{D.} x$^2$ + 17y > 0\par\vspace{2pt}
\textbf{\color{primary} Câu 18.} Bài tập khởi động Chương II: Điểm O(0; 0) thuộc miền nghiệm của bất phương trình nào sau đây:\par
\hspace*{1.5em}\textbf{A.} x - y > 19\quad\quad \textbf{B.} 2x + 3y - 21 < 0\quad\quad \textbf{C.} 3x + y $\ge$ 18\quad\quad \textbf{D.} 2x + 3y + 21 < 0\par\vspace{2pt}
\textbf{\color{primary} Câu 19.} Quan sát dữ kiện đề bài: Cặp số (x; y) nào dưới đây thuộc miền nghiệm của bất phương trình 2x - y + 20 > 0. Hãy chọn phương án đúng:\par
\hspace*{1.5em}\textbf{A.} (0; 21)\quad\quad \textbf{B.} (-20; 1)\quad\quad \textbf{C.} (0; 23)\quad\quad \textbf{D.} (1; 0)\par\vspace{2pt}
\textbf{\color{primary} Câu 20.} Sử dụng công thức hoặc tính chất đầu bài của Chương II để giải quyết yêu cầu: Bất phương trình nào sau đây là bất phương trình bậc nhất hai ẩn x, y:\par
\hspace*{1.5em}\textbf{A.} x + y + z < 5\quad\quad \textbf{B.} 3x - 21y $\le$ 7\quad\quad \textbf{C.} x$^2$ + 20y > 0\quad\quad \textbf{D.} 2x - 3xy $\ge$ 4\par\vspace{2pt}
\subsection*{\color{accent} Phần II. Mức độ Trung bình (Thông hiểu -- Câu 21 đến Câu 40)}
\textbf{\color{primary} Câu 21.} Cho hệ bất phương trình: x $\ge$ 0, y $\ge$ 0 và x + y $\le$ 5. Điểm nào sau đây thuộc miền nghiệm của hệ?\par
\hspace*{1.5em}\textbf{A.} (4; 3)\quad\quad \textbf{B.} (1; 2)\quad\quad \textbf{C.} (-1; 2)\quad\quad \textbf{D.} (2; -1)\par\vspace{2pt}
\textbf{\color{primary} Câu 22.} Vận dụng quy tắc trong Chương II (Bất phương trình và hệ bất phương trình bậc nhất hai ẩn) để giải bài toán: Tìm tất cả các giá trị của tham số m để điểm M(2; -1) thuộc miền nghiệm của bất phương trình mx - 3y > 6:\par
\hspace*{1.5em}\textbf{A.} m > 4.5\quad\quad \textbf{B.} m > 1.5\quad\quad \textbf{C.} m < 1.5\quad\quad \textbf{D.} m $\le$ 1.5\par\vspace{2pt}
\textbf{\color{primary} Câu 23.} Biến đổi và tính toán theo kiến thức Chương II: Miền nghiệm của hệ bất phương trình x $\ge$ 0, y $\ge$ 0, x/5 + y/4 $\le$ 1 là một miền tam giác vuông OAB. Diện tích miền tam giác OAB đó bằng. Kết quả thu được là:\par
\hspace*{1.5em}\textbf{A.} 20 (đơn vị diện tích)\quad\quad \textbf{B.} 10 (đơn vị diện tích)\quad\quad \textbf{C.} 9 (đơn vị diện tích)\quad\quad \textbf{D.} 5 (đơn vị diện tích)\par\vspace{2pt}
\textbf{\color{primary} Câu 24.} Xét bài toán mức độ thông hiểu sau: Cho hệ bất phương trình: x $\ge$ 0, y $\ge$ 0 và x + y $\le$ 8. Điểm nào sau đây thuộc miền nghiệm của hệ. Đáp án đúng là:\par
\hspace*{1.5em}\textbf{A.} (2; -1)\quad\quad \textbf{B.} (7; 3)\quad\quad \textbf{C.} (-1; 5)\quad\quad \textbf{D.} (1; 5)\par\vspace{2pt}
\textbf{\color{primary} Câu 25.} Để hoàn thành phiếu học tập Chương II, hãy xác định: Tìm tất cả các giá trị của tham số m để điểm M(2; -1) thuộc miền nghiệm của bất phương trình mx - 3y > 9:\par
\hspace*{1.5em}\textbf{A.} m > 3\quad\quad \textbf{B.} m > 6.0\quad\quad \textbf{C.} m < 3.0\quad\quad \textbf{D.} m $\le$ 3.0\par\vspace{2pt}
\textbf{\color{primary} Câu 26.} Phân tích dữ kiện bài toán: Miền nghiệm của hệ bất phương trình x $\ge$ 0, y $\ge$ 0, x/8 + y/4 $\le$ 1 là một miền tam giác vuông OAB. Diện tích miền tam giác OAB đó bằng. Giá trị cần tìm bằng:\par
\hspace*{1.5em}\textbf{A.} 32 (đơn vị diện tích)\quad\quad \textbf{B.} 8 (đơn vị diện tích)\quad\quad \textbf{C.} 16 (đơn vị diện tích)\quad\quad \textbf{D.} 12 (đơn vị diện tích)\par\vspace{2pt}
\textbf{\color{primary} Câu 27.} Áp dụng tính chất đã học ở Chương II vào bài tập sau: Cho hệ bất phương trình: x $\ge$ 0, y $\ge$ 0 và x + y $\le$ 11. Điểm nào sau đây thuộc miền nghiệm của hệ:\par
\hspace*{1.5em}\textbf{A.} (-1; 8)\quad\quad \textbf{B.} (2; -1)\quad\quad \textbf{C.} (1; 8)\quad\quad \textbf{D.} (10; 3)\par\vspace{2pt}
\textbf{\color{primary} Câu 28.} Thực hiện các bước suy luận hợp lý để giải quyết yêu cầu: Tìm tất cả các giá trị của tham số m để điểm M(2; -1) thuộc miền nghiệm của bất phương trình mx - 3y > 12:\par
\hspace*{1.5em}\textbf{A.} m < 4.5\quad\quad \textbf{B.} m $\le$ 4.5\quad\quad \textbf{C.} m > 4.5\quad\quad \textbf{D.} m > 7.5\par\vspace{2pt}
\textbf{\color{primary} Câu 29.} Cho bài toán tự luyện (Bất phương trình và hệ bất phương trình bậc nhất hai ẩn): Miền nghiệm của hệ bất phương trình x $\ge$ 0, y $\ge$ 0, x/11 + y/4 $\le$ 1 là một miền tam giác vuông OAB. Diện tích miền tam giác OAB đó bằng. Phương án nào sau đây chính xác?\par
\hspace*{1.5em}\textbf{A.} 22 (đơn vị diện tích)\quad\quad \textbf{B.} 11 (đơn vị diện tích)\quad\quad \textbf{C.} 44 (đơn vị diện tích)\quad\quad \textbf{D.} 15 (đơn vị diện tích)\par\vspace{2pt}
\textbf{\color{primary} Câu 30.} Kết hợp các phép biến đổi trong Chương II, hãy tìm kết quả của bài toán: Cho hệ bất phương trình: x $\ge$ 0, y $\ge$ 0 và x + y $\le$ 14. Điểm nào sau đây thuộc miền nghiệm của hệ:\par
\hspace*{1.5em}\textbf{A.} (13; 3)\quad\quad \textbf{B.} (2; -1)\quad\quad \textbf{C.} (1; 11)\quad\quad \textbf{D.} (-1; 11)\par\vspace{2pt}
\textbf{\color{primary} Câu 31.} Một bài kiểm tra định kỳ nội dung Chương II nêu câu hỏi: Tìm tất cả các giá trị của tham số m để điểm M(2; -1) thuộc miền nghiệm của bất phương trình mx - 3y > 15. Hãy chọn đáp số đúng:\par
\hspace*{1.5em}\textbf{A.} m $\le$ 6.0\quad\quad \textbf{B.} m < 6.0\quad\quad \textbf{C.} m > 9.0\quad\quad \textbf{D.} m > 6\par\vspace{2pt}
\textbf{\color{primary} Câu 32.} Tìm lời giải đúng cho bài toán thông hiểu: Miền nghiệm của hệ bất phương trình x $\ge$ 0, y $\ge$ 0, x/14 + y/4 $\le$ 1 là một miền tam giác vuông OAB. Diện tích miền tam giác OAB đó bằng:\par
\hspace*{1.5em}\textbf{A.} 28 (đơn vị diện tích)\quad\quad \textbf{B.} 18 (đơn vị diện tích)\quad\quad \textbf{C.} 56 (đơn vị diện tích)\quad\quad \textbf{D.} 14 (đơn vị diện tích)\par\vspace{2pt}
\textbf{\color{primary} Câu 33.} Sử dụng mối liên hệ giữa các đại lượng trong Chương II: Cho hệ bất phương trình: x $\ge$ 0, y $\ge$ 0 và x + y $\le$ 17. Điểm nào sau đây thuộc miền nghiệm của hệ:\par
\hspace*{1.5em}\textbf{A.} (16; 3)\quad\quad \textbf{B.} (2; -1)\quad\quad \textbf{C.} (1; 14)\quad\quad \textbf{D.} (-1; 14)\par\vspace{2pt}
\textbf{\color{primary} Câu 34.} Hãy tính toán và chọn phương án đúng đối với bài tập: Tìm tất cả các giá trị của tham số m để điểm M(2; -1) thuộc miền nghiệm của bất phương trình mx - 3y > 18:\par
\hspace*{1.5em}\textbf{A.} m > 10.5\quad\quad \textbf{B.} m < 7.5\quad\quad \textbf{C.} m > 7.5\quad\quad \textbf{D.} m $\le$ 7.5\par\vspace{2pt}
\textbf{\color{primary} Câu 35.} Dựa trên các tính chất trọng tâm của Chương II, xét bài toán: Miền nghiệm của hệ bất phương trình x $\ge$ 0, y $\ge$ 0, x/17 + y/4 $\le$ 1 là một miền tam giác vuông OAB. Diện tích miền tam giác OAB đó bằng:\par
\hspace*{1.5em}\textbf{A.} 68 (đơn vị diện tích)\quad\quad \textbf{B.} 34 (đơn vị diện tích)\quad\quad \textbf{C.} 17 (đơn vị diện tích)\quad\quad \textbf{D.} 21 (đơn vị diện tích)\par\vspace{2pt}
\textbf{\color{primary} Câu 36.} Thực hành giải toán (Bất phương trình và hệ bất phương trình bậc nhất hai ẩn): Cho hệ bất phương trình: x $\ge$ 0, y $\ge$ 0 và x + y $\le$ 20. Điểm nào sau đây thuộc miền nghiệm của hệ. Đáp án nào dưới đây thỏa mãn?\par
\hspace*{1.5em}\textbf{A.} (-1; 17)\quad\quad \textbf{B.} (19; 3)\quad\quad \textbf{C.} (1; 17)\quad\quad \textbf{D.} (2; -1)\par\vspace{2pt}
\textbf{\color{primary} Câu 37.} Kiểm tra kỹ năng biến đổi và suy luận qua bài toán: Tìm tất cả các giá trị của tham số m để điểm M(2; -1) thuộc miền nghiệm của bất phương trình mx - 3y > 21:\par
\hspace*{1.5em}\textbf{A.} m > 12.0\quad\quad \textbf{B.} m > 9\quad\quad \textbf{C.} m $\le$ 9.0\quad\quad \textbf{D.} m < 9.0\par\vspace{2pt}
\textbf{\color{primary} Câu 38.} Cho giả thiết thuộc phạm vi Chương II: Miền nghiệm của hệ bất phương trình x $\ge$ 0, y $\ge$ 0, x/20 + y/4 $\le$ 1 là một miền tam giác vuông OAB. Diện tích miền tam giác OAB đó bằng. Kết luận đúng là:\par
\hspace*{1.5em}\textbf{A.} 20 (đơn vị diện tích)\quad\quad \textbf{B.} 24 (đơn vị diện tích)\quad\quad \textbf{C.} 40 (đơn vị diện tích)\quad\quad \textbf{D.} 80 (đơn vị diện tích)\par\vspace{2pt}
\textbf{\color{primary} Câu 39.} Rút gọn hoặc tính giá trị theo yêu cầu sau của Chương II: Cho hệ bất phương trình: x $\ge$ 0, y $\ge$ 0 và x + y $\le$ 23. Điểm nào sau đây thuộc miền nghiệm của hệ:\par
\hspace*{1.5em}\textbf{A.} (-1; 20)\quad\quad \textbf{B.} (1; 20)\quad\quad \textbf{C.} (2; -1)\quad\quad \textbf{D.} (22; 3)\par\vspace{2pt}
\textbf{\color{primary} Câu 40.} Xác định phương án đúng sau khi thực hiện phép tính / suy luận: Tìm tất cả các giá trị của tham số m để điểm M(2; -1) thuộc miền nghiệm của bất phương trình mx - 3y > 24:\par
\hspace*{1.5em}\textbf{A.} m > 10.5\quad\quad \textbf{B.} m > 13.5\quad\quad \textbf{C.} m $\le$ 10.5\quad\quad \textbf{D.} m < 10.5\par\vspace{2pt}
\subsection*{\color{accent} Phần III. Mức độ Khó (Vận dụng \& Vận dụng cao -- Câu 41 đến Câu 50)}
\textbf{\color{primary} Câu 41.} Giá trị lớn nhất của biểu thức F(x; y) = 3x + 2y trên miền đa giác có các đỉnh O(0; 0), A(0; 4), B(2; 3), C(3; 0) bằng:\par
\hspace*{1.5em}\textbf{A.} 12\quad\quad \textbf{B.} 8\quad\quad \textbf{C.} 0\quad\quad \textbf{D.} 9\par\vspace{2pt}
\textbf{\color{primary} Câu 42.} Bài toán vận dụng (Chương II – Bất phương trình và hệ bất phương trình bậc nhất hai ẩn): Có bao nhiêu cặp số nguyên không âm (x; y) thỏa mãn đồng thời x $\ge$ 0, y $\ge$ 0 và x + y $\le$ 4. Kết quả đúng là:\par
\hspace*{1.5em}\textbf{A.} 15 cặp\quad\quad \textbf{B.} 25 cặp\quad\quad \textbf{C.} 10 cặp\quad\quad \textbf{D.} 16 cặp\par\vspace{2pt}
\textbf{\color{primary} Câu 43.} Vận dụng kiến thức nâng cao trong phạm vi Chương II để giải bài toán: Giá trị lớn nhất của biểu thức F(x; y) = 3x + 2y trên miền đa giác có các đỉnh O(0; 0), A(0; 6), B(2; 5), C(3; 0) bằng:\par
\hspace*{1.5em}\textbf{A.} 0\quad\quad \textbf{B.} 12\quad\quad \textbf{C.} 16\quad\quad \textbf{D.} 9\par\vspace{2pt}
\textbf{\color{primary} Câu 44.} Phân tích và giải bài toán tổng hợp sau của Chương II: Có bao nhiêu cặp số nguyên không âm (x; y) thỏa mãn đồng thời x $\ge$ 0, y $\ge$ 0 và x + y $\le$ 6:\par
\hspace*{1.5em}\textbf{A.} 21 cặp\quad\quad \textbf{B.} 36 cặp\quad\quad \textbf{C.} 28 cặp\quad\quad \textbf{D.} 49 cặp\par\vspace{2pt}
\textbf{\color{primary} Câu 45.} Thử thách tư duy trong phạm vi kiến thức Chương II: Giá trị lớn nhất của biểu thức F(x; y) = 3x + 2y trên miền đa giác có các đỉnh O(0; 0), A(0; 8), B(2; 7), C(3; 0) bằng. Đáp số chính xác là:\par
\hspace*{1.5em}\textbf{A.} 9\quad\quad \textbf{B.} 0\quad\quad \textbf{C.} 20\quad\quad \textbf{D.} 16\par\vspace{2pt}
\textbf{\color{primary} Câu 46.} Sử dụng phương pháp giải toán của Chương II để tìm đáp án cho bài toán: Có bao nhiêu cặp số nguyên không âm (x; y) thỏa mãn đồng thời x $\ge$ 0, y $\ge$ 0 và x + y $\le$ 8:\par
\hspace*{1.5em}\textbf{A.} 45 cặp\quad\quad \textbf{B.} 64 cặp\quad\quad \textbf{C.} 36 cặp\quad\quad \textbf{D.} 81 cặp\par\vspace{2pt}
\textbf{\color{primary} Câu 47.} Xét bài toán vận dụng thực tiễn / nâng cao thuộc Chương II: Giá trị lớn nhất của biểu thức F(x; y) = 3x + 2y trên miền đa giác có các đỉnh O(0; 0), A(0; 10), B(2; 9), C(3; 0) bằng:\par
\hspace*{1.5em}\textbf{A.} 24\quad\quad \textbf{B.} 0\quad\quad \textbf{C.} 20\quad\quad \textbf{D.} 9\par\vspace{2pt}
\textbf{\color{primary} Câu 48.} Tìm kết quả cuối cùng của bài toán phân loại học sinh khá giỏi (Chương II): Có bao nhiêu cặp số nguyên không âm (x; y) thỏa mãn đồng thời x $\ge$ 0, y $\ge$ 0 và x + y $\le$ 10:\par
\hspace*{1.5em}\textbf{A.} 100 cặp\quad\quad \textbf{B.} 121 cặp\quad\quad \textbf{C.} 66 cặp\quad\quad \textbf{D.} 55 cặp\par\vspace{2pt}
\textbf{\color{primary} Câu 49.} Kết hợp các tính chất đã học trong Chương II, hãy giải quyết bài toán: Giá trị lớn nhất của biểu thức F(x; y) = 3x + 2y trên miền đa giác có các đỉnh O(0; 0), A(0; 12), B(2; 11), C(3; 0) bằng:\par
\hspace*{1.5em}\textbf{A.} 0\quad\quad \textbf{B.} 24\quad\quad \textbf{C.} 9\quad\quad \textbf{D.} 28\par\vspace{2pt}
\textbf{\color{primary} Câu 50.} Câu hỏi vận dụng cao cuối Chương II: Có bao nhiêu cặp số nguyên không âm (x; y) thỏa mãn đồng thời x $\ge$ 0, y $\ge$ 0 và x + y $\le$ 12. Phương án đúng là:\par
\hspace*{1.5em}\textbf{A.} 78 cặp\quad\quad \textbf{B.} 169 cặp\quad\quad \textbf{C.} 91 cặp\quad\quad \textbf{D.} 144 cặp\par\vspace{2pt}
\subsection*{\color{primary} Bảng đáp án Chương II: Bất phương trình và hệ bất phương trình bậc nhất hai ẩn}
\begin{center}
\renewcommand{\arraystretch}{1.25}
\begin{tabularx}{\textwidth}{*{10}{>{\centering\arraybackslash}X}}
\toprule
\textbf{\color{primary} Cột 1} & \textbf{\color{primary} Cột 2} & \textbf{\color{primary} Cột 3} & \textbf{\color{primary} Cột 4} & \textbf{\color{primary} Cột 5} & \textbf{\color{primary} Cột 6} & \textbf{\color{primary} Cột 7} & \textbf{\color{primary} Cột 8} & \textbf{\color{primary} Cột 9} & \textbf{\color{primary} Cột 10} \\
\midrule
\textbf{1.D} & \textbf{2.B} & \textbf{3.A} & \textbf{4.B} & \textbf{5.C} & \textbf{6.B} & \textbf{7.C} & \textbf{8.A} & \textbf{9.C} & \textbf{10.B} \\
\rowcolor{rowtint} \textbf{11.D} & \textbf{12.C} & \textbf{13.B} & \textbf{14.A} & \textbf{15.B} & \textbf{16.A} & \textbf{17.C} & \textbf{18.B} & \textbf{19.D} & \textbf{20.B} \\
\textbf{21.B} & \textbf{22.B} & \textbf{23.B} & \textbf{24.D} & \textbf{25.A} & \textbf{26.C} & \textbf{27.C} & \textbf{28.C} & \textbf{29.A} & \textbf{30.C} \\
\rowcolor{rowtint} \textbf{31.D} & \textbf{32.A} & \textbf{33.C} & \textbf{34.C} & \textbf{35.B} & \textbf{36.C} & \textbf{37.B} & \textbf{38.C} & \textbf{39.B} & \textbf{40.A} \\
\textbf{41.A} & \textbf{42.A} & \textbf{43.C} & \textbf{44.C} & \textbf{45.C} & \textbf{46.A} & \textbf{47.A} & \textbf{48.C} & \textbf{49.D} & \textbf{50.C} \\
\bottomrule
\end{tabularx}
\end{center}\vspace{10pt}
\section*{\color{primary} CHƯƠNG III: HỆ THỨC LƯỢNG TRONG TAM GIÁC}
\textit{Trọng tâm kiến thức SGK: Giá trị lượng giác của một góc từ 0° đến 180°, định lí côsin, định lí sin, các công thức tính diện tích tam giác và giải tam giác.}\par\vspace{4pt}
\subsection*{\color{accent} Phần I. Mức độ Dễ (Nhận biết -- Câu 1 đến Câu 20)}
\textbf{\color{primary} Câu 1.} Với góc $\alpha$ (0° $\le$ $\alpha$ $\le$ 180°), khẳng định nào sau đây là ĐÚNG về hai góc bù nhau?\par
\begin{itemize}\setlength{\itemsep}{1pt}\setlength{\parskip}{0pt}
  \item[] \textbf{A.} tan(180° - $\alpha$) = tan $\alpha$
  \item[] \textbf{B.} sin(180° - $\alpha$) = -sin $\alpha$
  \item[] \textbf{C.} sin(180° - $\alpha$) = sin $\alpha$ và cos(180° - $\alpha$) = -cos $\alpha$
  \item[] \textbf{D.} cos(180° - $\alpha$) = cos $\alpha$
\end{itemize}\vspace{2pt}
\textbf{\color{primary} Câu 2.} Thực hiện yêu cầu nhận biết sau: Trong tam giác ABC có BC = a, AC = b, AB = c và bán kính đường tròn ngoại tiếp R, định lí sin được phát biểu là.\par
\begin{itemize}\setlength{\itemsep}{1pt}\setlength{\parskip}{0pt}
  \item[] \textbf{A.} a / sin A = b / sin B = c / sin C = 2R
  \item[] \textbf{B.} a / sin A = b / sin B = c / sin C = R
  \item[] \textbf{C.} a$^2$ = b$^2$ + c$^2$ - 2bc$\cdot$cos A
  \item[] \textbf{D.} a$\cdot$sin A = b$\cdot$sin B = c$\cdot$sin C = 2R
\end{itemize}\vspace{2pt}
\textbf{\color{primary} Câu 3.} Trong bài học thuộc Chương III, xét nội dung: Cho tam giác ABC có hai cạnh b = 6, c = 6 và góc A = 30°. Diện tích S của tam giác ABC (S = 1/2 $\cdot$ bc$\cdot$sin A) bằng. Đáp án đúng là:\par
\hspace*{1.5em}\textbf{A.} 18\quad\quad \textbf{B.} 36\quad\quad \textbf{C.} 6$\sqrt{3}$\quad\quad \textbf{D.} 9\par\vspace{2pt}
\textbf{\color{primary} Câu 4.} Dựa vào kiến thức chuẩn SGK của Chương III (Hệ thức lượng trong tam giác), với góc $\alpha$ (0° $\le$ $\alpha$ $\le$ 180°), khẳng định nào sau đây là ĐÚNG về hai góc bù nhau:\par
\begin{itemize}\setlength{\itemsep}{1pt}\setlength{\parskip}{0pt}
  \item[] \textbf{A.} cos(180° - $\alpha$) = cos $\alpha$
  \item[] \textbf{B.} tan(180° - $\alpha$) = tan $\alpha$
  \item[] \textbf{C.} sin(180° - $\alpha$) = -sin $\alpha$
  \item[] \textbf{D.} sin(180° - $\alpha$) = sin $\alpha$ và cos(180° - $\alpha$) = -cos $\alpha$
\end{itemize}\vspace{2pt}
\textbf{\color{primary} Câu 5.} Chọn kết quả đúng nhất khi xét: Trong tam giác ABC có BC = a, AC = b, AB = c và bán kính đường tròn ngoại tiếp R, định lí sin được phát biểu là:\par
\begin{itemize}\setlength{\itemsep}{1pt}\setlength{\parskip}{0pt}
  \item[] \textbf{A.} a / sin A = b / sin B = c / sin C = 2R
  \item[] \textbf{B.} a$\cdot$sin A = b$\cdot$sin B = c$\cdot$sin C = 2R
  \item[] \textbf{C.} a$^2$ = b$^2$ + c$^2$ - 2bc$\cdot$cos A
  \item[] \textbf{D.} a / sin A = b / sin B = c / sin C = R
\end{itemize}\vspace{2pt}
\textbf{\color{primary} Câu 6.} Khi làm bài tập cơ bản về hệ thức lượng trong tam giác, một học sinh gặp câu hỏi: Cho tam giác ABC có hai cạnh b = 9, c = 6 và góc A = 30°. Diện tích S của tam giác ABC (S = 1/2 $\cdot$ bc$\cdot$sin A) bằng. Kết quả nào dưới đây đúng?\par
\hspace*{1.5em}\textbf{A.} 9$\sqrt{3}$\quad\quad \textbf{B.} 13.5\quad\quad \textbf{C.} 54\quad\quad \textbf{D.} 27\par\vspace{2pt}
\textbf{\color{primary} Câu 7.} Xác định giá trị hoặc khẳng định đúng cho bài toán sau: Với góc $\alpha$ (0° $\le$ $\alpha$ $\le$ 180°), khẳng định nào sau đây là ĐÚNG về hai góc bù nhau:\par
\begin{itemize}\setlength{\itemsep}{1pt}\setlength{\parskip}{0pt}
  \item[] \textbf{A.} cos(180° - $\alpha$) = cos $\alpha$
  \item[] \textbf{B.} sin(180° - $\alpha$) = -sin $\alpha$
  \item[] \textbf{C.} sin(180° - $\alpha$) = sin $\alpha$ và cos(180° - $\alpha$) = -cos $\alpha$
  \item[] \textbf{D.} tan(180° - $\alpha$) = tan $\alpha$
\end{itemize}\vspace{2pt}
\textbf{\color{primary} Câu 8.} Áp dụng trực tiếp định nghĩa và quy tắc trong Chương III: Trong tam giác ABC có BC = a, AC = b, AB = c và bán kính đường tròn ngoại tiếp R, định lí sin được phát biểu là:\par
\begin{itemize}\setlength{\itemsep}{1pt}\setlength{\parskip}{0pt}
  \item[] \textbf{A.} a$\cdot$sin A = b$\cdot$sin B = c$\cdot$sin C = 2R
  \item[] \textbf{B.} a / sin A = b / sin B = c / sin C = R
  \item[] \textbf{C.} a$^2$ = b$^2$ + c$^2$ - 2bc$\cdot$cos A
  \item[] \textbf{D.} a / sin A = b / sin B = c / sin C = 2R
\end{itemize}\vspace{2pt}
\textbf{\color{primary} Câu 9.} Kiểm tra nhanh kiến thức nền tảng: Cho tam giác ABC có hai cạnh b = 12, c = 6 và góc A = 30°. Diện tích S của tam giác ABC (S = 1/2 $\cdot$ bc$\cdot$sin A) bằng:\par
\hspace*{1.5em}\textbf{A.} 36\quad\quad \textbf{B.} 72\quad\quad \textbf{C.} 12$\sqrt{3}$\quad\quad \textbf{D.} 18\par\vspace{2pt}
\textbf{\color{primary} Câu 10.} Cho dữ kiện nhận biết sau đây: Với góc $\alpha$ (0° $\le$ $\alpha$ $\le$ 180°), khẳng định nào sau đây là ĐÚNG về hai góc bù nhau. Phương án chính xác là:\par
\begin{itemize}\setlength{\itemsep}{1pt}\setlength{\parskip}{0pt}
  \item[] \textbf{A.} tan(180° - $\alpha$) = tan $\alpha$
  \item[] \textbf{B.} sin(180° - $\alpha$) = sin $\alpha$ và cos(180° - $\alpha$) = -cos $\alpha$
  \item[] \textbf{C.} cos(180° - $\alpha$) = cos $\alpha$
  \item[] \textbf{D.} sin(180° - $\alpha$) = -sin $\alpha$
\end{itemize}\vspace{2pt}
\textbf{\color{primary} Câu 11.} Hãy chỉ ra kết quả đúng trong bốn lựa chọn A, B, C, D đối với yêu cầu: Trong tam giác ABC có BC = a, AC = b, AB = c và bán kính đường tròn ngoại tiếp R, định lí sin được phát biểu là:\par
\begin{itemize}\setlength{\itemsep}{1pt}\setlength{\parskip}{0pt}
  \item[] \textbf{A.} a$^2$ = b$^2$ + c$^2$ - 2bc$\cdot$cos A
  \item[] \textbf{B.} a / sin A = b / sin B = c / sin C = R
  \item[] \textbf{C.} a / sin A = b / sin B = c / sin C = 2R
  \item[] \textbf{D.} a$\cdot$sin A = b$\cdot$sin B = c$\cdot$sin C = 2R
\end{itemize}\vspace{2pt}
\textbf{\color{primary} Câu 12.} Câu hỏi ôn tập lý thuyết và kỹ năng cơ bản (Chương III): Cho tam giác ABC có hai cạnh b = 15, c = 6 và góc A = 30°. Diện tích S của tam giác ABC (S = 1/2 $\cdot$ bc$\cdot$sin A) bằng:\par
\hspace*{1.5em}\textbf{A.} 45\quad\quad \textbf{B.} 15$\sqrt{3}$\quad\quad \textbf{C.} 90\quad\quad \textbf{D.} 22.5\par\vspace{2pt}
\textbf{\color{primary} Câu 13.} Theo nội dung sách giáo khoa Kết nối tri thức với cuộc sống (Chương III), với góc $\alpha$ (0° $\le$ $\alpha$ $\le$ 180°), khẳng định nào sau đây là ĐÚNG về hai góc bù nhau:\par
\begin{itemize}\setlength{\itemsep}{1pt}\setlength{\parskip}{0pt}
  \item[] \textbf{A.} cos(180° - $\alpha$) = cos $\alpha$
  \item[] \textbf{B.} sin(180° - $\alpha$) = -sin $\alpha$
  \item[] \textbf{C.} tan(180° - $\alpha$) = tan $\alpha$
  \item[] \textbf{D.} sin(180° - $\alpha$) = sin $\alpha$ và cos(180° - $\alpha$) = -cos $\alpha$
\end{itemize}\vspace{2pt}
\textbf{\color{primary} Câu 14.} Tìm đáp số đúng của bài toán cơ bản sau: Trong tam giác ABC có BC = a, AC = b, AB = c và bán kính đường tròn ngoại tiếp R, định lí sin được phát biểu là:\par
\begin{itemize}\setlength{\itemsep}{1pt}\setlength{\parskip}{0pt}
  \item[] \textbf{A.} a$^2$ = b$^2$ + c$^2$ - 2bc$\cdot$cos A
  \item[] \textbf{B.} a / sin A = b / sin B = c / sin C = 2R
  \item[] \textbf{C.} a$\cdot$sin A = b$\cdot$sin B = c$\cdot$sin C = 2R
  \item[] \textbf{D.} a / sin A = b / sin B = c / sin C = R
\end{itemize}\vspace{2pt}
\textbf{\color{primary} Câu 15.} Xét đối tượng toán học trong Chương III: Cho tam giác ABC có hai cạnh b = 18, c = 6 và góc A = 30°. Diện tích S của tam giác ABC (S = 1/2 $\cdot$ bc$\cdot$sin A) bằng. Giá trị hoặc tính chất tương ứng là:\par
\hspace*{1.5em}\textbf{A.} 54\quad\quad \textbf{B.} 27\quad\quad \textbf{C.} 18$\sqrt{3}$\quad\quad \textbf{D.} 108\par\vspace{2pt}
\textbf{\color{primary} Câu 16.} Hoàn thành phép xác định cơ bản sau đây: Với góc $\alpha$ (0° $\le$ $\alpha$ $\le$ 180°), khẳng định nào sau đây là ĐÚNG về hai góc bù nhau:\par
\begin{itemize}\setlength{\itemsep}{1pt}\setlength{\parskip}{0pt}
  \item[] \textbf{A.} tan(180° - $\alpha$) = tan $\alpha$
  \item[] \textbf{B.} sin(180° - $\alpha$) = -sin $\alpha$
  \item[] \textbf{C.} cos(180° - $\alpha$) = cos $\alpha$
  \item[] \textbf{D.} sin(180° - $\alpha$) = sin $\alpha$ và cos(180° - $\alpha$) = -cos $\alpha$
\end{itemize}\vspace{2pt}
\textbf{\color{primary} Câu 17.} Trong các khẳng định hoặc giá trị dưới đây, đâu là đáp án đúng cho câu hỏi: Trong tam giác ABC có BC = a, AC = b, AB = c và bán kính đường tròn ngoại tiếp R, định lí sin được phát biểu là?\par
\begin{itemize}\setlength{\itemsep}{1pt}\setlength{\parskip}{0pt}
  \item[] \textbf{A.} a / sin A = b / sin B = c / sin C = 2R
  \item[] \textbf{B.} a$\cdot$sin A = b$\cdot$sin B = c$\cdot$sin C = 2R
  \item[] \textbf{C.} a / sin A = b / sin B = c / sin C = R
  \item[] \textbf{D.} a$^2$ = b$^2$ + c$^2$ - 2bc$\cdot$cos A
\end{itemize}\vspace{2pt}
\textbf{\color{primary} Câu 18.} Bài tập khởi động Chương III: Cho tam giác ABC có hai cạnh b = 21, c = 6 và góc A = 30°. Diện tích S của tam giác ABC (S = 1/2 $\cdot$ bc$\cdot$sin A) bằng:\par
\hspace*{1.5em}\textbf{A.} 126\quad\quad \textbf{B.} 63\quad\quad \textbf{C.} 21$\sqrt{3}$\quad\quad \textbf{D.} 31.5\par\vspace{2pt}
\textbf{\color{primary} Câu 19.} Quan sát dữ kiện đề bài: Với góc $\alpha$ (0° $\le$ $\alpha$ $\le$ 180°), khẳng định nào sau đây là ĐÚNG về hai góc bù nhau. Hãy chọn phương án đúng:\par
\begin{itemize}\setlength{\itemsep}{1pt}\setlength{\parskip}{0pt}
  \item[] \textbf{A.} sin(180° - $\alpha$) = sin $\alpha$ và cos(180° - $\alpha$) = -cos $\alpha$
  \item[] \textbf{B.} tan(180° - $\alpha$) = tan $\alpha$
  \item[] \textbf{C.} cos(180° - $\alpha$) = cos $\alpha$
  \item[] \textbf{D.} sin(180° - $\alpha$) = -sin $\alpha$
\end{itemize}\vspace{2pt}
\textbf{\color{primary} Câu 20.} Sử dụng công thức hoặc tính chất đầu bài của Chương III để giải quyết yêu cầu: Trong tam giác ABC có BC = a, AC = b, AB = c và bán kính đường tròn ngoại tiếp R, định lí sin được phát biểu là:\par
\begin{itemize}\setlength{\itemsep}{1pt}\setlength{\parskip}{0pt}
  \item[] \textbf{A.} a / sin A = b / sin B = c / sin C = 2R
  \item[] \textbf{B.} a / sin A = b / sin B = c / sin C = R
  \item[] \textbf{C.} a$^2$ = b$^2$ + c$^2$ - 2bc$\cdot$cos A
  \item[] \textbf{D.} a$\cdot$sin A = b$\cdot$sin B = c$\cdot$sin C = 2R
\end{itemize}\vspace{2pt}
\subsection*{\color{accent} Phần II. Mức độ Trung bình (Thông hiểu -- Câu 21 đến Câu 40)}
\textbf{\color{primary} Câu 21.} Cho tam giác ABC có AB = 3, AC = 5 và góc A = 120° (biết cos 120° = -1/2). Theo định lí côsin, bình phương cạnh BC$^2$ bằng:\par
\hspace*{1.5em}\textbf{A.} 49\quad\quad \textbf{B.} 19\quad\quad \textbf{C.} 64\quad\quad \textbf{D.} 34\par\vspace{2pt}
\textbf{\color{primary} Câu 22.} Vận dụng quy tắc trong Chương III (Hệ thức lượng trong tam giác) để giải bài toán: Cho tam giác ABC có cạnh BC = 7 và góc A = 30°. Bán kính R của đường tròn ngoại tiếp tam giác ABC bằng:\par
\hspace*{1.5em}\textbf{A.} R = 3.5\quad\quad \textbf{B.} R = 7$\sqrt{3}$\quad\quad \textbf{C.} R = 7\quad\quad \textbf{D.} R = 14\par\vspace{2pt}
\textbf{\color{primary} Câu 23.} Biến đổi và tính toán theo kiến thức Chương III: Cho tam giác ABC có ba cạnh a = 5, b = 6, c = 7 (nửa chu vi p = 9). Theo công thức Heron, diện tích S của tam giác ABC bằng. Kết quả thu được là:\par
\hspace*{1.5em}\textbf{A.} 18\quad\quad \textbf{B.} 12$\sqrt{6}$\quad\quad \textbf{C.} 6$\sqrt{3}$\quad\quad \textbf{D.} 6$\sqrt{6}$\par\vspace{2pt}
\textbf{\color{primary} Câu 24.} Xét bài toán mức độ thông hiểu sau: Cho tam giác ABC có AB = 6, AC = 8 và góc A = 120° (biết cos 120° = -1/2). Theo định lí côsin, bình phương cạnh BC$^2$ bằng. Đáp án đúng là:\par
\hspace*{1.5em}\textbf{A.} 52\quad\quad \textbf{B.} 148\quad\quad \textbf{C.} 100\quad\quad \textbf{D.} 196\par\vspace{2pt}
\textbf{\color{primary} Câu 25.} Để hoàn thành phiếu học tập Chương III, hãy xác định: Cho tam giác ABC có cạnh BC = 10 và góc A = 30°. Bán kính R của đường tròn ngoại tiếp tam giác ABC bằng:\par
\hspace*{1.5em}\textbf{A.} R = 10$\sqrt{3}$\quad\quad \textbf{B.} R = 10\quad\quad \textbf{C.} R = 5.0\quad\quad \textbf{D.} R = 20\par\vspace{2pt}
\textbf{\color{primary} Câu 26.} Phân tích dữ kiện bài toán: Cho tam giác ABC có ba cạnh a = 5, b = 6, c = 7 (nửa chu vi p = 9). Theo công thức Heron, diện tích S của tam giác ABC bằng. Giá trị cần tìm bằng:\par
\hspace*{1.5em}\textbf{A.} 6$\sqrt{3}$\quad\quad \textbf{B.} 18\quad\quad \textbf{C.} 12$\sqrt{6}$\quad\quad \textbf{D.} 6$\sqrt{6}$\par\vspace{2pt}
\textbf{\color{primary} Câu 27.} Áp dụng tính chất đã học ở Chương III vào bài tập sau: Cho tam giác ABC có AB = 9, AC = 11 và góc A = 120° (biết cos 120° = -1/2). Theo định lí côsin, bình phương cạnh BC$^2$ bằng:\par
\hspace*{1.5em}\textbf{A.} 400\quad\quad \textbf{B.} 301\quad\quad \textbf{C.} 202\quad\quad \textbf{D.} 103\par\vspace{2pt}
\textbf{\color{primary} Câu 28.} Thực hiện các bước suy luận hợp lý để giải quyết yêu cầu: Cho tam giác ABC có cạnh BC = 13 và góc A = 30°. Bán kính R của đường tròn ngoại tiếp tam giác ABC bằng:\par
\hspace*{1.5em}\textbf{A.} R = 6.5\quad\quad \textbf{B.} R = 26\quad\quad \textbf{C.} R = 13\quad\quad \textbf{D.} R = 13$\sqrt{3}$\par\vspace{2pt}
\textbf{\color{primary} Câu 29.} Cho bài toán tự luyện (Hệ thức lượng trong tam giác): Cho tam giác ABC có ba cạnh a = 5, b = 6, c = 7 (nửa chu vi p = 9). Theo công thức Heron, diện tích S của tam giác ABC bằng. Phương án nào sau đây chính xác?\par
\hspace*{1.5em}\textbf{A.} 12$\sqrt{6}$\quad\quad \textbf{B.} 6$\sqrt{3}$\quad\quad \textbf{C.} 6$\sqrt{6}$\quad\quad \textbf{D.} 18\par\vspace{2pt}
\textbf{\color{primary} Câu 30.} Kết hợp các phép biến đổi trong Chương III, hãy tìm kết quả của bài toán: Cho tam giác ABC có AB = 12, AC = 14 và góc A = 120° (biết cos 120° = -1/2). Theo định lí côsin, bình phương cạnh BC$^2$ bằng:\par
\hspace*{1.5em}\textbf{A.} 172\quad\quad \textbf{B.} 676\quad\quad \textbf{C.} 508\quad\quad \textbf{D.} 340\par\vspace{2pt}
\textbf{\color{primary} Câu 31.} Một bài kiểm tra định kỳ nội dung Chương III nêu câu hỏi: Cho tam giác ABC có cạnh BC = 16 và góc A = 30°. Bán kính R của đường tròn ngoại tiếp tam giác ABC bằng. Hãy chọn đáp số đúng:\par
\hspace*{1.5em}\textbf{A.} R = 16$\sqrt{3}$\quad\quad \textbf{B.} R = 16\quad\quad \textbf{C.} R = 8.0\quad\quad \textbf{D.} R = 32\par\vspace{2pt}
\textbf{\color{primary} Câu 32.} Tìm lời giải đúng cho bài toán thông hiểu: Cho tam giác ABC có ba cạnh a = 5, b = 6, c = 7 (nửa chu vi p = 9). Theo công thức Heron, diện tích S của tam giác ABC bằng:\par
\hspace*{1.5em}\textbf{A.} 12$\sqrt{6}$\quad\quad \textbf{B.} 6$\sqrt{3}$\quad\quad \textbf{C.} 6$\sqrt{6}$\quad\quad \textbf{D.} 18\par\vspace{2pt}
\textbf{\color{primary} Câu 33.} Sử dụng mối liên hệ giữa các đại lượng trong Chương III: Cho tam giác ABC có AB = 15, AC = 17 và góc A = 120° (biết cos 120° = -1/2). Theo định lí côsin, bình phương cạnh BC$^2$ bằng:\par
\hspace*{1.5em}\textbf{A.} 259\quad\quad \textbf{B.} 1024\quad\quad \textbf{C.} 769\quad\quad \textbf{D.} 514\par\vspace{2pt}
\textbf{\color{primary} Câu 34.} Hãy tính toán và chọn phương án đúng đối với bài tập: Cho tam giác ABC có cạnh BC = 19 và góc A = 30°. Bán kính R của đường tròn ngoại tiếp tam giác ABC bằng:\par
\hspace*{1.5em}\textbf{A.} R = 9.5\quad\quad \textbf{B.} R = 19\quad\quad \textbf{C.} R = 19$\sqrt{3}$\quad\quad \textbf{D.} R = 38\par\vspace{2pt}
\textbf{\color{primary} Câu 35.} Dựa trên các tính chất trọng tâm của Chương III, xét bài toán: Cho tam giác ABC có ba cạnh a = 5, b = 6, c = 7 (nửa chu vi p = 9). Theo công thức Heron, diện tích S của tam giác ABC bằng:\par
\hspace*{1.5em}\textbf{A.} 12$\sqrt{6}$\quad\quad \textbf{B.} 18\quad\quad \textbf{C.} 6$\sqrt{3}$\quad\quad \textbf{D.} 6$\sqrt{6}$\par\vspace{2pt}
\textbf{\color{primary} Câu 36.} Thực hành giải toán (Hệ thức lượng trong tam giác): Cho tam giác ABC có AB = 18, AC = 20 và góc A = 120° (biết cos 120° = -1/2). Theo định lí côsin, bình phương cạnh BC$^2$ bằng. Đáp án nào dưới đây thỏa mãn?\par
\hspace*{1.5em}\textbf{A.} 1444\quad\quad \textbf{B.} 364\quad\quad \textbf{C.} 1084\quad\quad \textbf{D.} 724\par\vspace{2pt}
\textbf{\color{primary} Câu 37.} Kiểm tra kỹ năng biến đổi và suy luận qua bài toán: Cho tam giác ABC có cạnh BC = 22 và góc A = 30°. Bán kính R của đường tròn ngoại tiếp tam giác ABC bằng:\par
\hspace*{1.5em}\textbf{A.} R = 11.0\quad\quad \textbf{B.} R = 44\quad\quad \textbf{C.} R = 22\quad\quad \textbf{D.} R = 22$\sqrt{3}$\par\vspace{2pt}
\textbf{\color{primary} Câu 38.} Cho giả thiết thuộc phạm vi Chương III: Cho tam giác ABC có ba cạnh a = 5, b = 6, c = 7 (nửa chu vi p = 9). Theo công thức Heron, diện tích S của tam giác ABC bằng. Kết luận đúng là:\par
\hspace*{1.5em}\textbf{A.} 6$\sqrt{6}$\quad\quad \textbf{B.} 12$\sqrt{6}$\quad\quad \textbf{C.} 18\quad\quad \textbf{D.} 6$\sqrt{3}$\par\vspace{2pt}
\textbf{\color{primary} Câu 39.} Rút gọn hoặc tính giá trị theo yêu cầu sau của Chương III: Cho tam giác ABC có AB = 21, AC = 23 và góc A = 120° (biết cos 120° = -1/2). Theo định lí côsin, bình phương cạnh BC$^2$ bằng:\par
\hspace*{1.5em}\textbf{A.} 487\quad\quad \textbf{B.} 1453\quad\quad \textbf{C.} 1936\quad\quad \textbf{D.} 970\par\vspace{2pt}
\textbf{\color{primary} Câu 40.} Xác định phương án đúng sau khi thực hiện phép tính / suy luận: Cho tam giác ABC có cạnh BC = 25 và góc A = 30°. Bán kính R của đường tròn ngoại tiếp tam giác ABC bằng:\par
\hspace*{1.5em}\textbf{A.} R = 12.5\quad\quad \textbf{B.} R = 50\quad\quad \textbf{C.} R = 25$\sqrt{3}$\quad\quad \textbf{D.} R = 25\par\vspace{2pt}
\subsection*{\color{accent} Phần III. Mức độ Khó (Vận dụng \& Vận dụng cao -- Câu 41 đến Câu 50)}
\textbf{\color{primary} Câu 41.} Cho tam giác ABC có độ dài ba cạnh a = 13, b = 14, c = 15 (diện tích S = 84). Bán kính r của đường tròn nội tiếp tam giác ABC (r = S/p) bằng:\par
\hspace*{1.5em}\textbf{A.} r = 8,125\quad\quad \textbf{B.} r = 4\quad\quad \textbf{C.} r = 6\quad\quad \textbf{D.} r = 2\par\vspace{2pt}
\textbf{\color{primary} Câu 42.} Bài toán vận dụng (Chương III – Hệ thức lượng trong tam giác): Cho góc $\alpha$ tù (90° < $\alpha$ < 180°) thỏa mãn sin $\alpha$ = 3/5. Giá trị của biểu thức M = 5$\cdot$cos $\alpha$ + 10$\cdot$tan $\alpha$ bằng. Kết quả đúng là:\par
\hspace*{1.5em}\textbf{A.} 11,5\quad\quad \textbf{B.} -3,5\quad\quad \textbf{C.} -11,5\quad\quad \textbf{D.} 3,5\par\vspace{2pt}
\textbf{\color{primary} Câu 43.} Vận dụng kiến thức nâng cao trong phạm vi Chương III để giải bài toán: Cho tam giác ABC có ba cạnh tỉ lệ với 5 : 12 : 13, độ dài ba cạnh lần lượt là 15, 36, 39. Bán kính r của đường tròn nội tiếp tam giác ABC (r = S/p) bằng:\par
\hspace*{1.5em}\textbf{A.} r = 7.5\quad\quad \textbf{B.} r = 3\quad\quad \textbf{C.} r = 9\quad\quad \textbf{D.} r = 6\par\vspace{2pt}
\textbf{\color{primary} Câu 44.} Phân tích và giải bài toán tổng hợp sau của Chương III: Cho tam giác ABC có độ dài ba cạnh a = 13, b = 14, c = 15 (diện tích S = 84). Bán kính r của đường tròn nội tiếp tam giác ABC (r = S/p) bằng:\par
\hspace*{1.5em}\textbf{A.} r = 8,125\quad\quad \textbf{B.} r = 4\quad\quad \textbf{C.} r = 2\quad\quad \textbf{D.} r = 6\par\vspace{2pt}
\textbf{\color{primary} Câu 45.} Thử thách tư duy trong phạm vi kiến thức Chương III: Cho góc $\alpha$ tù (90° < $\alpha$ < 180°) thỏa mãn sin $\alpha$ = 3/5. Giá trị của biểu thức M = 5$\cdot$cos $\alpha$ + 10$\cdot$tan $\alpha$ bằng. Đáp số chính xác là:\par
\hspace*{1.5em}\textbf{A.} 3,5\quad\quad \textbf{B.} 11,5\quad\quad \textbf{C.} -11,5\quad\quad \textbf{D.} -3,5\par\vspace{2pt}
\textbf{\color{primary} Câu 46.} Sử dụng phương pháp giải toán của Chương III để tìm đáp án cho bài toán: Cho tam giác ABC có ba cạnh tỉ lệ với 5 : 12 : 13, độ dài ba cạnh lần lượt là 30, 72, 78. Bán kính r của đường tròn nội tiếp tam giác ABC (r = S/p) bằng:\par
\hspace*{1.5em}\textbf{A.} r = 18\quad\quad \textbf{B.} r = 6\quad\quad \textbf{C.} r = 15.0\quad\quad \textbf{D.} r = 12\par\vspace{2pt}
\textbf{\color{primary} Câu 47.} Xét bài toán vận dụng thực tiễn / nâng cao thuộc Chương III: Cho tam giác ABC có độ dài ba cạnh a = 13, b = 14, c = 15 (diện tích S = 84). Bán kính r của đường tròn nội tiếp tam giác ABC (r = S/p) bằng:\par
\hspace*{1.5em}\textbf{A.} r = 6\quad\quad \textbf{B.} r = 8,125\quad\quad \textbf{C.} r = 2\quad\quad \textbf{D.} r = 4\par\vspace{2pt}
\textbf{\color{primary} Câu 48.} Tìm kết quả cuối cùng của bài toán phân loại học sinh khá giỏi (Chương III): Cho góc $\alpha$ tù (90° < $\alpha$ < 180°) thỏa mãn sin $\alpha$ = 3/5. Giá trị của biểu thức M = 5$\cdot$cos $\alpha$ + 10$\cdot$tan $\alpha$ bằng:\par
\hspace*{1.5em}\textbf{A.} -11,5\quad\quad \textbf{B.} 3,5\quad\quad \textbf{C.} 11,5\quad\quad \textbf{D.} -3,5\par\vspace{2pt}
\textbf{\color{primary} Câu 49.} Kết hợp các tính chất đã học trong Chương III, hãy giải quyết bài toán: Cho tam giác ABC có ba cạnh tỉ lệ với 5 : 12 : 13, độ dài ba cạnh lần lượt là 45, 108, 117. Bán kính r của đường tròn nội tiếp tam giác ABC (r = S/p) bằng:\par
\hspace*{1.5em}\textbf{A.} r = 18\quad\quad \textbf{B.} r = 27\quad\quad \textbf{C.} r = 9\quad\quad \textbf{D.} r = 22.5\par\vspace{2pt}
\textbf{\color{primary} Câu 50.} Câu hỏi vận dụng cao cuối Chương III: Cho tam giác ABC có độ dài ba cạnh a = 13, b = 14, c = 15 (diện tích S = 84). Bán kính r của đường tròn nội tiếp tam giác ABC (r = S/p) bằng. Phương án đúng là:\par
\hspace*{1.5em}\textbf{A.} r = 4\quad\quad \textbf{B.} r = 6\quad\quad \textbf{C.} r = 2\quad\quad \textbf{D.} r = 8,125\par\vspace{2pt}
\subsection*{\color{primary} Bảng đáp án Chương III: Hệ thức lượng trong tam giác}
\begin{center}
\renewcommand{\arraystretch}{1.25}
\begin{tabularx}{\textwidth}{*{10}{>{\centering\arraybackslash}X}}
\toprule
\textbf{\color{primary} Cột 1} & \textbf{\color{primary} Cột 2} & \textbf{\color{primary} Cột 3} & \textbf{\color{primary} Cột 4} & \textbf{\color{primary} Cột 5} & \textbf{\color{primary} Cột 6} & \textbf{\color{primary} Cột 7} & \textbf{\color{primary} Cột 8} & \textbf{\color{primary} Cột 9} & \textbf{\color{primary} Cột 10} \\
\midrule
\textbf{1.C} & \textbf{2.A} & \textbf{3.D} & \textbf{4.D} & \textbf{5.A} & \textbf{6.B} & \textbf{7.C} & \textbf{8.D} & \textbf{9.D} & \textbf{10.B} \\
\rowcolor{rowtint} \textbf{11.C} & \textbf{12.D} & \textbf{13.D} & \textbf{14.B} & \textbf{15.B} & \textbf{16.D} & \textbf{17.A} & \textbf{18.D} & \textbf{19.A} & \textbf{20.A} \\
\textbf{21.A} & \textbf{22.C} & \textbf{23.D} & \textbf{24.B} & \textbf{25.B} & \textbf{26.D} & \textbf{27.B} & \textbf{28.C} & \textbf{29.C} & \textbf{30.C} \\
\rowcolor{rowtint} \textbf{31.B} & \textbf{32.C} & \textbf{33.C} & \textbf{34.B} & \textbf{35.D} & \textbf{36.C} & \textbf{37.C} & \textbf{38.A} & \textbf{39.B} & \textbf{40.D} \\
\textbf{41.B} & \textbf{42.C} & \textbf{43.D} & \textbf{44.B} & \textbf{45.C} & \textbf{46.D} & \textbf{47.D} & \textbf{48.A} & \textbf{49.A} & \textbf{50.A} \\
\bottomrule
\end{tabularx}
\end{center}\vspace{10pt}
\section*{\color{primary} CHƯƠNG IV: VECTƠ}
\textit{Trọng tâm kiến thức SGK: Các khái niệm mở đầu về vectơ, tổng và hiệu của hai vectơ, tích của một vectơ với một số, vectơ trong mặt phẳng tọa độ và tích vô hướng của hai vectơ.}\par\vspace{4pt}
\subsection*{\color{accent} Phần I. Mức độ Dễ (Nhận biết -- Câu 1 đến Câu 20)}
\textbf{\color{primary} Câu 1.} Với ba điểm A, B, C bất kỳ, đẳng thức nào sau đây là quy tắc ba điểm đối với phép cộng vectơ?\par
\hspace*{1.5em}\textbf{A.} $\overrightarrow{AB}$ + $\overrightarrow{AC}$ = $\overrightarrow{BC}$\quad\quad \textbf{B.} $\overrightarrow{AB}$ - $\overrightarrow{BC}$ = $\overrightarrow{AC}$\quad\quad \textbf{C.} $\overrightarrow{AB}$ + $\overrightarrow{BC}$ = $\overrightarrow{AC}$\quad\quad \textbf{D.} $\overrightarrow{BA}$ + $\overrightarrow{BC}$ = $\overrightarrow{AC}$\par\vspace{2pt}
\textbf{\color{primary} Câu 2.} Thực hiện yêu cầu nhận biết sau: Trong mặt phẳng tọa độ Oxy, cho hai điểm A(1; 3) và B(4; 7). Tọa độ của vectơ $\overrightarrow{AB}$ là.\par
\hspace*{1.5em}\textbf{A.} (5; 10)\quad\quad \textbf{B.} (-3; -4)\quad\quad \textbf{C.} (3; 4)\quad\quad \textbf{D.} (4; 3)\par\vspace{2pt}
\textbf{\color{primary} Câu 3.} Trong bài học thuộc Chương IV, xét nội dung: Cho hai vectơ $\vec{u}$ = (2; 3) và $\vec{v}$ = (3; -2). Tích vô hướng $\vec{u}$ $\cdot$ $\vec{v}$ bằng. Đáp án đúng là:\par
\hspace*{1.5em}\textbf{A.} (6; -6)\quad\quad \textbf{B.} 12\quad\quad \textbf{C.} 0\quad\quad \textbf{D.} 0 (1)\par\vspace{2pt}
\textbf{\color{primary} Câu 4.} Dựa vào kiến thức chuẩn SGK của Chương IV (Vectơ), với ba điểm A, B, C bất kỳ, đẳng thức nào sau đây là quy tắc ba điểm đối với phép cộng vectơ:\par
\hspace*{1.5em}\textbf{A.} $\overrightarrow{AB}$ + $\overrightarrow{AC}$ = $\overrightarrow{BC}$\quad\quad \textbf{B.} $\overrightarrow{AB}$ + $\overrightarrow{BC}$ = $\overrightarrow{AC}$\quad\quad \textbf{C.} $\overrightarrow{BA}$ + $\overrightarrow{BC}$ = $\overrightarrow{AC}$\quad\quad \textbf{D.} $\overrightarrow{AB}$ - $\overrightarrow{BC}$ = $\overrightarrow{AC}$\par\vspace{2pt}
\textbf{\color{primary} Câu 5.} Chọn kết quả đúng nhất khi xét: Trong mặt phẳng tọa độ Oxy, cho hai điểm A(1; 6) và B(4; 10). Tọa độ của vectơ $\overrightarrow{AB}$ là:\par
\hspace*{1.5em}\textbf{A.} (-3; -4)\quad\quad \textbf{B.} (3; 4)\quad\quad \textbf{C.} (5; 16)\quad\quad \textbf{D.} (4; 3)\par\vspace{2pt}
\textbf{\color{primary} Câu 6.} Khi làm bài tập cơ bản về vectơ, một học sinh gặp câu hỏi: Cho hai vectơ $\vec{u}$ = (2; 6) và $\vec{v}$ = (3; -2). Tích vô hướng $\vec{u}$ $\cdot$ $\vec{v}$ bằng. Kết quả nào dưới đây đúng?\par
\hspace*{1.5em}\textbf{A.} -6\quad\quad \textbf{B.} -6 (1)\quad\quad \textbf{C.} 18\quad\quad \textbf{D.} (6; -12)\par\vspace{2pt}
\textbf{\color{primary} Câu 7.} Xác định giá trị hoặc khẳng định đúng cho bài toán sau: Với ba điểm A, B, C bất kỳ, đẳng thức nào sau đây là quy tắc ba điểm đối với phép cộng vectơ:\par
\hspace*{1.5em}\textbf{A.} $\overrightarrow{AB}$ - $\overrightarrow{BC}$ = $\overrightarrow{AC}$\quad\quad \textbf{B.} $\overrightarrow{BA}$ + $\overrightarrow{BC}$ = $\overrightarrow{AC}$\quad\quad \textbf{C.} $\overrightarrow{AB}$ + $\overrightarrow{AC}$ = $\overrightarrow{BC}$\quad\quad \textbf{D.} $\overrightarrow{AB}$ + $\overrightarrow{BC}$ = $\overrightarrow{AC}$\par\vspace{2pt}
\textbf{\color{primary} Câu 8.} Áp dụng trực tiếp định nghĩa và quy tắc trong Chương IV: Trong mặt phẳng tọa độ Oxy, cho hai điểm A(1; 9) và B(4; 13). Tọa độ của vectơ $\overrightarrow{AB}$ là:\par
\hspace*{1.5em}\textbf{A.} (5; 22)\quad\quad \textbf{B.} (3; 4)\quad\quad \textbf{C.} (4; 3)\quad\quad \textbf{D.} (-3; -4)\par\vspace{2pt}
\textbf{\color{primary} Câu 9.} Kiểm tra nhanh kiến thức nền tảng: Cho hai vectơ $\vec{u}$ = (2; 9) và $\vec{v}$ = (3; -2). Tích vô hướng $\vec{u}$ $\cdot$ $\vec{v}$ bằng:\par
\hspace*{1.5em}\textbf{A.} (6; -18)\quad\quad \textbf{B.} -12\quad\quad \textbf{C.} -12 (1)\quad\quad \textbf{D.} 24\par\vspace{2pt}
\textbf{\color{primary} Câu 10.} Cho dữ kiện nhận biết sau đây: Với ba điểm A, B, C bất kỳ, đẳng thức nào sau đây là quy tắc ba điểm đối với phép cộng vectơ. Phương án chính xác là:\par
\hspace*{1.5em}\textbf{A.} $\overrightarrow{AB}$ + $\overrightarrow{BC}$ = $\overrightarrow{AC}$\quad\quad \textbf{B.} $\overrightarrow{AB}$ + $\overrightarrow{AC}$ = $\overrightarrow{BC}$\quad\quad \textbf{C.} $\overrightarrow{BA}$ + $\overrightarrow{BC}$ = $\overrightarrow{AC}$\quad\quad \textbf{D.} $\overrightarrow{AB}$ - $\overrightarrow{BC}$ = $\overrightarrow{AC}$\par\vspace{2pt}
\textbf{\color{primary} Câu 11.} Hãy chỉ ra kết quả đúng trong bốn lựa chọn A, B, C, D đối với yêu cầu: Trong mặt phẳng tọa độ Oxy, cho hai điểm A(1; 12) và B(4; 16). Tọa độ của vectơ $\overrightarrow{AB}$ là:\par
\hspace*{1.5em}\textbf{A.} (5; 28)\quad\quad \textbf{B.} (3; 4)\quad\quad \textbf{C.} (-3; -4)\quad\quad \textbf{D.} (4; 3)\par\vspace{2pt}
\textbf{\color{primary} Câu 12.} Câu hỏi ôn tập lý thuyết và kỹ năng cơ bản (Chương IV): Cho hai vectơ $\vec{u}$ = (2; 12) và $\vec{v}$ = (3; -2). Tích vô hướng $\vec{u}$ $\cdot$ $\vec{v}$ bằng:\par
\hspace*{1.5em}\textbf{A.} -18 (1)\quad\quad \textbf{B.} 30\quad\quad \textbf{C.} (6; -24)\quad\quad \textbf{D.} -18\par\vspace{2pt}
\textbf{\color{primary} Câu 13.} Theo nội dung sách giáo khoa Kết nối tri thức với cuộc sống (Chương IV), với ba điểm A, B, C bất kỳ, đẳng thức nào sau đây là quy tắc ba điểm đối với phép cộng vectơ:\par
\hspace*{1.5em}\textbf{A.} $\overrightarrow{AB}$ + $\overrightarrow{AC}$ = $\overrightarrow{BC}$\quad\quad \textbf{B.} $\overrightarrow{AB}$ - $\overrightarrow{BC}$ = $\overrightarrow{AC}$\quad\quad \textbf{C.} $\overrightarrow{AB}$ + $\overrightarrow{BC}$ = $\overrightarrow{AC}$\quad\quad \textbf{D.} $\overrightarrow{BA}$ + $\overrightarrow{BC}$ = $\overrightarrow{AC}$\par\vspace{2pt}
\textbf{\color{primary} Câu 14.} Tìm đáp số đúng của bài toán cơ bản sau: Trong mặt phẳng tọa độ Oxy, cho hai điểm A(1; 15) và B(4; 19). Tọa độ của vectơ $\overrightarrow{AB}$ là:\par
\hspace*{1.5em}\textbf{A.} (3; 4)\quad\quad \textbf{B.} (-3; -4)\quad\quad \textbf{C.} (4; 3)\quad\quad \textbf{D.} (5; 34)\par\vspace{2pt}
\textbf{\color{primary} Câu 15.} Xét đối tượng toán học trong Chương IV: Cho hai vectơ $\vec{u}$ = (2; 15) và $\vec{v}$ = (3; -2). Tích vô hướng $\vec{u}$ $\cdot$ $\vec{v}$ bằng. Giá trị hoặc tính chất tương ứng là:\par
\hspace*{1.5em}\textbf{A.} 36\quad\quad \textbf{B.} -24 (1)\quad\quad \textbf{C.} -24\quad\quad \textbf{D.} (6; -30)\par\vspace{2pt}
\textbf{\color{primary} Câu 16.} Hoàn thành phép xác định cơ bản sau đây: Với ba điểm A, B, C bất kỳ, đẳng thức nào sau đây là quy tắc ba điểm đối với phép cộng vectơ:\par
\hspace*{1.5em}\textbf{A.} $\overrightarrow{AB}$ + $\overrightarrow{BC}$ = $\overrightarrow{AC}$\quad\quad \textbf{B.} $\overrightarrow{BA}$ + $\overrightarrow{BC}$ = $\overrightarrow{AC}$\quad\quad \textbf{C.} $\overrightarrow{AB}$ - $\overrightarrow{BC}$ = $\overrightarrow{AC}$\quad\quad \textbf{D.} $\overrightarrow{AB}$ + $\overrightarrow{AC}$ = $\overrightarrow{BC}$\par\vspace{2pt}
\textbf{\color{primary} Câu 17.} Trong các khẳng định hoặc giá trị dưới đây, đâu là đáp án đúng cho câu hỏi: Trong mặt phẳng tọa độ Oxy, cho hai điểm A(1; 18) và B(4; 22). Tọa độ của vectơ $\overrightarrow{AB}$ là?\par
\hspace*{1.5em}\textbf{A.} (-3; -4)\quad\quad \textbf{B.} (4; 3)\quad\quad \textbf{C.} (3; 4)\quad\quad \textbf{D.} (5; 40)\par\vspace{2pt}
\textbf{\color{primary} Câu 18.} Bài tập khởi động Chương IV: Cho hai vectơ $\vec{u}$ = (2; 18) và $\vec{v}$ = (3; -2). Tích vô hướng $\vec{u}$ $\cdot$ $\vec{v}$ bằng:\par
\hspace*{1.5em}\textbf{A.} (6; -36)\quad\quad \textbf{B.} -30\quad\quad \textbf{C.} 42\quad\quad \textbf{D.} -30 (1)\par\vspace{2pt}
\textbf{\color{primary} Câu 19.} Quan sát dữ kiện đề bài: Với ba điểm A, B, C bất kỳ, đẳng thức nào sau đây là quy tắc ba điểm đối với phép cộng vectơ. Hãy chọn phương án đúng:\par
\hspace*{1.5em}\textbf{A.} $\overrightarrow{BA}$ + $\overrightarrow{BC}$ = $\overrightarrow{AC}$\quad\quad \textbf{B.} $\overrightarrow{AB}$ + $\overrightarrow{AC}$ = $\overrightarrow{BC}$\quad\quad \textbf{C.} $\overrightarrow{AB}$ - $\overrightarrow{BC}$ = $\overrightarrow{AC}$\quad\quad \textbf{D.} $\overrightarrow{AB}$ + $\overrightarrow{BC}$ = $\overrightarrow{AC}$\par\vspace{2pt}
\textbf{\color{primary} Câu 20.} Sử dụng công thức hoặc tính chất đầu bài của Chương IV để giải quyết yêu cầu: Trong mặt phẳng tọa độ Oxy, cho hai điểm A(1; 21) và B(4; 25). Tọa độ của vectơ $\overrightarrow{AB}$ là:\par
\hspace*{1.5em}\textbf{A.} (5; 46)\quad\quad \textbf{B.} (4; 3)\quad\quad \textbf{C.} (3; 4)\quad\quad \textbf{D.} (-3; -4)\par\vspace{2pt}
\subsection*{\color{accent} Phần II. Mức độ Trung bình (Thông hiểu -- Câu 21 đến Câu 40)}
\textbf{\color{primary} Câu 21.} Cho hình vuông ABCD có độ dài cạnh bằng 3. Độ dài của vectơ tổng |$\overrightarrow{AB}$ + $\overrightarrow{AD}$| bằng:\par
\hspace*{1.5em}\textbf{A.} 0\quad\quad \textbf{B.} 3$\sqrt{2}$\quad\quad \textbf{C.} 6\quad\quad \textbf{D.} 3\par\vspace{2pt}
\textbf{\color{primary} Câu 22.} Vận dụng quy tắc trong Chương IV (Vectơ) để giải bài toán: Trong mặt phẳng Oxy, tìm m để hai vectơ $\vec{a}$ = (3; -2) và $\vec{b}$ = (12; m) cùng phương:\par
\hspace*{1.5em}\textbf{A.} m = 18\quad\quad \textbf{B.} m = -18\quad\quad \textbf{C.} m = -8\quad\quad \textbf{D.} m = 8\par\vspace{2pt}
\textbf{\color{primary} Câu 23.} Biến đổi và tính toán theo kiến thức Chương IV: Cho hai vectơ $\vec{u}$ = (4; 3) và $\vec{v}$ = (-6; m). Giá trị của m để $\vec{u}$ vuông góc với $\vec{v}$ ($\vec{u}$ $\perp$ $\vec{v}$) là. Kết quả thu được là:\par
\hspace*{1.5em}\textbf{A.} m = -8\quad\quad \textbf{B.} m = 8\quad\quad \textbf{C.} m = 0\quad\quad \textbf{D.} m = 12\par\vspace{2pt}
\textbf{\color{primary} Câu 24.} Xét bài toán mức độ thông hiểu sau: Cho hình vuông ABCD có độ dài cạnh bằng 6. Độ dài của vectơ tổng |$\overrightarrow{AB}$ + $\overrightarrow{AD}$| bằng. Đáp án đúng là:\par
\hspace*{1.5em}\textbf{A.} 6$\sqrt{2}$\quad\quad \textbf{B.} 12\quad\quad \textbf{C.} 0\quad\quad \textbf{D.} 6\par\vspace{2pt}
\textbf{\color{primary} Câu 25.} Để hoàn thành phiếu học tập Chương IV, hãy xác định: Trong mặt phẳng Oxy, tìm m để hai vectơ $\vec{a}$ = (3; -2) và $\vec{b}$ = (12; m) cùng phương:\par
\hspace*{1.5em}\textbf{A.} m = -8\quad\quad \textbf{B.} m = 18\quad\quad \textbf{C.} m = -18\quad\quad \textbf{D.} m = 8\par\vspace{2pt}
\textbf{\color{primary} Câu 26.} Phân tích dữ kiện bài toán: Cho hai vectơ $\vec{u}$ = (7; 3) và $\vec{v}$ = (-6; m). Giá trị của m để $\vec{u}$ vuông góc với $\vec{v}$ ($\vec{u}$ $\perp$ $\vec{v}$) là. Giá trị cần tìm bằng:\par
\hspace*{1.5em}\textbf{A.} m = -14\quad\quad \textbf{B.} m = 0\quad\quad \textbf{C.} m = 14\quad\quad \textbf{D.} m = 21\par\vspace{2pt}
\textbf{\color{primary} Câu 27.} Áp dụng tính chất đã học ở Chương IV vào bài tập sau: Cho hình vuông ABCD có độ dài cạnh bằng 9. Độ dài của vectơ tổng |$\overrightarrow{AB}$ + $\overrightarrow{AD}$| bằng:\par
\hspace*{1.5em}\textbf{A.} 9$\sqrt{2}$\quad\quad \textbf{B.} 0\quad\quad \textbf{C.} 18\quad\quad \textbf{D.} 9\par\vspace{2pt}
\textbf{\color{primary} Câu 28.} Thực hiện các bước suy luận hợp lý để giải quyết yêu cầu: Trong mặt phẳng Oxy, tìm m để hai vectơ $\vec{a}$ = (3; -2) và $\vec{b}$ = (12; m) cùng phương:\par
\hspace*{1.5em}\textbf{A.} m = -18\quad\quad \textbf{B.} m = 18\quad\quad \textbf{C.} m = 8\quad\quad \textbf{D.} m = -8\par\vspace{2pt}
\textbf{\color{primary} Câu 29.} Cho bài toán tự luyện (Vectơ): Cho hai vectơ $\vec{u}$ = (10; 3) và $\vec{v}$ = (-6; m). Giá trị của m để $\vec{u}$ vuông góc với $\vec{v}$ ($\vec{u}$ $\perp$ $\vec{v}$) là. Phương án nào sau đây chính xác?\par
\hspace*{1.5em}\textbf{A.} m = 30\quad\quad \textbf{B.} m = 20\quad\quad \textbf{C.} m = 0\quad\quad \textbf{D.} m = -20\par\vspace{2pt}
\textbf{\color{primary} Câu 30.} Kết hợp các phép biến đổi trong Chương IV, hãy tìm kết quả của bài toán: Cho hình vuông ABCD có độ dài cạnh bằng 12. Độ dài của vectơ tổng |$\overrightarrow{AB}$ + $\overrightarrow{AD}$| bằng:\par
\hspace*{1.5em}\textbf{A.} 0\quad\quad \textbf{B.} 12$\sqrt{2}$\quad\quad \textbf{C.} 24\quad\quad \textbf{D.} 12\par\vspace{2pt}
\textbf{\color{primary} Câu 31.} Một bài kiểm tra định kỳ nội dung Chương IV nêu câu hỏi: Trong mặt phẳng Oxy, tìm m để hai vectơ $\vec{a}$ = (3; -2) và $\vec{b}$ = (12; m) cùng phương. Hãy chọn đáp số đúng:\par
\hspace*{1.5em}\textbf{A.} m = -8\quad\quad \textbf{B.} m = 8\quad\quad \textbf{C.} m = -18\quad\quad \textbf{D.} m = 18\par\vspace{2pt}
\textbf{\color{primary} Câu 32.} Tìm lời giải đúng cho bài toán thông hiểu: Cho hai vectơ $\vec{u}$ = (13; 3) và $\vec{v}$ = (-6; m). Giá trị của m để $\vec{u}$ vuông góc với $\vec{v}$ ($\vec{u}$ $\perp$ $\vec{v}$) là:\par
\hspace*{1.5em}\textbf{A.} m = -26\quad\quad \textbf{B.} m = 0\quad\quad \textbf{C.} m = 39\quad\quad \textbf{D.} m = 26\par\vspace{2pt}
\textbf{\color{primary} Câu 33.} Sử dụng mối liên hệ giữa các đại lượng trong Chương IV: Cho hình vuông ABCD có độ dài cạnh bằng 15. Độ dài của vectơ tổng |$\overrightarrow{AB}$ + $\overrightarrow{AD}$| bằng:\par
\hspace*{1.5em}\textbf{A.} 0\quad\quad \textbf{B.} 30\quad\quad \textbf{C.} 15$\sqrt{2}$\quad\quad \textbf{D.} 15\par\vspace{2pt}
\textbf{\color{primary} Câu 34.} Hãy tính toán và chọn phương án đúng đối với bài tập: Trong mặt phẳng Oxy, tìm m để hai vectơ $\vec{a}$ = (3; -2) và $\vec{b}$ = (12; m) cùng phương:\par
\hspace*{1.5em}\textbf{A.} m = -18\quad\quad \textbf{B.} m = 18\quad\quad \textbf{C.} m = 8\quad\quad \textbf{D.} m = -8\par\vspace{2pt}
\textbf{\color{primary} Câu 35.} Dựa trên các tính chất trọng tâm của Chương IV, xét bài toán: Cho hai vectơ $\vec{u}$ = (16; 3) và $\vec{v}$ = (-6; m). Giá trị của m để $\vec{u}$ vuông góc với $\vec{v}$ ($\vec{u}$ $\perp$ $\vec{v}$) là:\par
\hspace*{1.5em}\textbf{A.} m = 32\quad\quad \textbf{B.} m = 0\quad\quad \textbf{C.} m = 48\quad\quad \textbf{D.} m = -32\par\vspace{2pt}
\textbf{\color{primary} Câu 36.} Thực hành giải toán (Vectơ): Cho hình vuông ABCD có độ dài cạnh bằng 18. Độ dài của vectơ tổng |$\overrightarrow{AB}$ + $\overrightarrow{AD}$| bằng. Đáp án nào dưới đây thỏa mãn?\par
\hspace*{1.5em}\textbf{A.} 0\quad\quad \textbf{B.} 36\quad\quad \textbf{C.} 18\quad\quad \textbf{D.} 18$\sqrt{2}$\par\vspace{2pt}
\textbf{\color{primary} Câu 37.} Kiểm tra kỹ năng biến đổi và suy luận qua bài toán: Trong mặt phẳng Oxy, tìm m để hai vectơ $\vec{a}$ = (3; -2) và $\vec{b}$ = (12; m) cùng phương:\par
\hspace*{1.5em}\textbf{A.} m = -8\quad\quad \textbf{B.} m = -18\quad\quad \textbf{C.} m = 18\quad\quad \textbf{D.} m = 8\par\vspace{2pt}
\textbf{\color{primary} Câu 38.} Cho giả thiết thuộc phạm vi Chương IV: Cho hai vectơ $\vec{u}$ = (19; 3) và $\vec{v}$ = (-6; m). Giá trị của m để $\vec{u}$ vuông góc với $\vec{v}$ ($\vec{u}$ $\perp$ $\vec{v}$) là. Kết luận đúng là:\par
\hspace*{1.5em}\textbf{A.} m = 38\quad\quad \textbf{B.} m = 57\quad\quad \textbf{C.} m = -38\quad\quad \textbf{D.} m = 0\par\vspace{2pt}
\textbf{\color{primary} Câu 39.} Rút gọn hoặc tính giá trị theo yêu cầu sau của Chương IV: Cho hình vuông ABCD có độ dài cạnh bằng 21. Độ dài của vectơ tổng |$\overrightarrow{AB}$ + $\overrightarrow{AD}$| bằng:\par
\hspace*{1.5em}\textbf{A.} 42\quad\quad \textbf{B.} 0\quad\quad \textbf{C.} 21\quad\quad \textbf{D.} 21$\sqrt{2}$\par\vspace{2pt}
\textbf{\color{primary} Câu 40.} Xác định phương án đúng sau khi thực hiện phép tính / suy luận: Trong mặt phẳng Oxy, tìm m để hai vectơ $\vec{a}$ = (3; -2) và $\vec{b}$ = (12; m) cùng phương:\par
\hspace*{1.5em}\textbf{A.} m = 8\quad\quad \textbf{B.} m = -18\quad\quad \textbf{C.} m = 18\quad\quad \textbf{D.} m = -8\par\vspace{2pt}
\subsection*{\color{accent} Phần III. Mức độ Khó (Vận dụng \& Vận dụng cao -- Câu 41 đến Câu 50)}
\textbf{\color{primary} Câu 41.} Cho tam giác đều ABC có độ dài cạnh bằng 4. Tích vô hướng $\overrightarrow{AB}$ $\cdot$ $\overrightarrow{BC}$ (chú ý góc giữa hai vectơ $\overrightarrow{AB}$ và $\overrightarrow{BC}$ bằng 120°) bằng:\par
\hspace*{1.5em}\textbf{A.} 8\quad\quad \textbf{B.} -16\quad\quad \textbf{C.} -8\quad\quad \textbf{D.} 16\par\vspace{2pt}
\textbf{\color{primary} Câu 42.} Bài toán vận dụng (Chương IV – Vectơ): Trong tam giác ABC có G là trọng tâm và M là trung điểm của BC, biểu diễn vectơ $\overrightarrow{AG}$ theo hai vectơ $\overrightarrow{AB}$ và $\overrightarrow{AC}$ ta được. Kết quả đúng là:\par
\begin{itemize}\setlength{\itemsep}{1pt}\setlength{\parskip}{0pt}
  \item[] \textbf{A.} $\overrightarrow{AG}$ = 1/3 $\cdot$ $\overrightarrow{AB}$ + 2/3 $\cdot$ $\overrightarrow{AC}$
  \item[] \textbf{B.} $\overrightarrow{AG}$ = 1/2 $\cdot$ $\overrightarrow{AB}$ + 1/2 $\cdot$ $\overrightarrow{AC}$
  \item[] \textbf{C.} $\overrightarrow{AG}$ = 1/3 $\cdot$ $\overrightarrow{AB}$ + 1/3 $\cdot$ $\overrightarrow{AC}$
  \item[] \textbf{D.} $\overrightarrow{AG}$ = 2/3 $\cdot$ $\overrightarrow{AB}$ + 2/3 $\cdot$ $\overrightarrow{AC}$
\end{itemize}\vspace{2pt}
\textbf{\color{primary} Câu 43.} Vận dụng kiến thức nâng cao trong phạm vi Chương IV để giải bài toán: Cho tam giác đều ABC có độ dài cạnh bằng 8. Tích vô hướng $\overrightarrow{AB}$ $\cdot$ $\overrightarrow{BC}$ (chú ý góc giữa hai vectơ $\overrightarrow{AB}$ và $\overrightarrow{BC}$ bằng 120°) bằng:\par
\hspace*{1.5em}\textbf{A.} -32\quad\quad \textbf{B.} 32\quad\quad \textbf{C.} -64\quad\quad \textbf{D.} 64\par\vspace{2pt}
\textbf{\color{primary} Câu 44.} Phân tích và giải bài toán tổng hợp sau của Chương IV: Trong tam giác ABC có G là trọng tâm và M là trung điểm của BC, biểu diễn vectơ $\overrightarrow{AG}$ theo hai vectơ $\overrightarrow{AB}$ và $\overrightarrow{AC}$ ta được:\par
\begin{itemize}\setlength{\itemsep}{1pt}\setlength{\parskip}{0pt}
  \item[] \textbf{A.} $\overrightarrow{AG}$ = 1/3 $\cdot$ $\overrightarrow{AB}$ + 2/3 $\cdot$ $\overrightarrow{AC}$
  \item[] \textbf{B.} $\overrightarrow{AG}$ = 1/3 $\cdot$ $\overrightarrow{AB}$ + 1/3 $\cdot$ $\overrightarrow{AC}$
  \item[] \textbf{C.} $\overrightarrow{AG}$ = 1/2 $\cdot$ $\overrightarrow{AB}$ + 1/2 $\cdot$ $\overrightarrow{AC}$
  \item[] \textbf{D.} $\overrightarrow{AG}$ = 2/3 $\cdot$ $\overrightarrow{AB}$ + 2/3 $\cdot$ $\overrightarrow{AC}$
\end{itemize}\vspace{2pt}
\textbf{\color{primary} Câu 45.} Thử thách tư duy trong phạm vi kiến thức Chương IV: Cho tam giác đều ABC có độ dài cạnh bằng 12. Tích vô hướng $\overrightarrow{AB}$ $\cdot$ $\overrightarrow{BC}$ (chú ý góc giữa hai vectơ $\overrightarrow{AB}$ và $\overrightarrow{BC}$ bằng 120°) bằng. Đáp số chính xác là:\par
\hspace*{1.5em}\textbf{A.} 72\quad\quad \textbf{B.} -72\quad\quad \textbf{C.} -144\quad\quad \textbf{D.} 144\par\vspace{2pt}
\textbf{\color{primary} Câu 46.} Sử dụng phương pháp giải toán của Chương IV để tìm đáp án cho bài toán: Trong tam giác ABC có G là trọng tâm và M là trung điểm của BC, biểu diễn vectơ $\overrightarrow{AG}$ theo hai vectơ $\overrightarrow{AB}$ và $\overrightarrow{AC}$ ta được:\par
\begin{itemize}\setlength{\itemsep}{1pt}\setlength{\parskip}{0pt}
  \item[] \textbf{A.} $\overrightarrow{AG}$ = 1/3 $\cdot$ $\overrightarrow{AB}$ + 1/3 $\cdot$ $\overrightarrow{AC}$
  \item[] \textbf{B.} $\overrightarrow{AG}$ = 1/2 $\cdot$ $\overrightarrow{AB}$ + 1/2 $\cdot$ $\overrightarrow{AC}$
  \item[] \textbf{C.} $\overrightarrow{AG}$ = 2/3 $\cdot$ $\overrightarrow{AB}$ + 2/3 $\cdot$ $\overrightarrow{AC}$
  \item[] \textbf{D.} $\overrightarrow{AG}$ = 1/3 $\cdot$ $\overrightarrow{AB}$ + 2/3 $\cdot$ $\overrightarrow{AC}$
\end{itemize}\vspace{2pt}
\textbf{\color{primary} Câu 47.} Xét bài toán vận dụng thực tiễn / nâng cao thuộc Chương IV: Cho tam giác đều ABC có độ dài cạnh bằng 16. Tích vô hướng $\overrightarrow{AB}$ $\cdot$ $\overrightarrow{BC}$ (chú ý góc giữa hai vectơ $\overrightarrow{AB}$ và $\overrightarrow{BC}$ bằng 120°) bằng:\par
\hspace*{1.5em}\textbf{A.} 256\quad\quad \textbf{B.} -256\quad\quad \textbf{C.} 128\quad\quad \textbf{D.} -128\par\vspace{2pt}
\textbf{\color{primary} Câu 48.} Tìm kết quả cuối cùng của bài toán phân loại học sinh khá giỏi (Chương IV): Trong tam giác ABC có G là trọng tâm và M là trung điểm của BC, biểu diễn vectơ $\overrightarrow{AG}$ theo hai vectơ $\overrightarrow{AB}$ và $\overrightarrow{AC}$ ta được:\par
\begin{itemize}\setlength{\itemsep}{1pt}\setlength{\parskip}{0pt}
  \item[] \textbf{A.} $\overrightarrow{AG}$ = 1/3 $\cdot$ $\overrightarrow{AB}$ + 2/3 $\cdot$ $\overrightarrow{AC}$
  \item[] \textbf{B.} $\overrightarrow{AG}$ = 1/3 $\cdot$ $\overrightarrow{AB}$ + 1/3 $\cdot$ $\overrightarrow{AC}$
  \item[] \textbf{C.} $\overrightarrow{AG}$ = 2/3 $\cdot$ $\overrightarrow{AB}$ + 2/3 $\cdot$ $\overrightarrow{AC}$
  \item[] \textbf{D.} $\overrightarrow{AG}$ = 1/2 $\cdot$ $\overrightarrow{AB}$ + 1/2 $\cdot$ $\overrightarrow{AC}$
\end{itemize}\vspace{2pt}
\textbf{\color{primary} Câu 49.} Kết hợp các tính chất đã học trong Chương IV, hãy giải quyết bài toán: Cho tam giác đều ABC có độ dài cạnh bằng 20. Tích vô hướng $\overrightarrow{AB}$ $\cdot$ $\overrightarrow{BC}$ (chú ý góc giữa hai vectơ $\overrightarrow{AB}$ và $\overrightarrow{BC}$ bằng 120°) bằng:\par
\hspace*{1.5em}\textbf{A.} 200\quad\quad \textbf{B.} 400\quad\quad \textbf{C.} -200\quad\quad \textbf{D.} -400\par\vspace{2pt}
\textbf{\color{primary} Câu 50.} Câu hỏi vận dụng cao cuối Chương IV: Trong tam giác ABC có G là trọng tâm và M là trung điểm của BC, biểu diễn vectơ $\overrightarrow{AG}$ theo hai vectơ $\overrightarrow{AB}$ và $\overrightarrow{AC}$ ta được. Phương án đúng là:\par
\begin{itemize}\setlength{\itemsep}{1pt}\setlength{\parskip}{0pt}
  \item[] \textbf{A.} $\overrightarrow{AG}$ = 2/3 $\cdot$ $\overrightarrow{AB}$ + 2/3 $\cdot$ $\overrightarrow{AC}$
  \item[] \textbf{B.} $\overrightarrow{AG}$ = 1/3 $\cdot$ $\overrightarrow{AB}$ + 1/3 $\cdot$ $\overrightarrow{AC}$
  \item[] \textbf{C.} $\overrightarrow{AG}$ = 1/2 $\cdot$ $\overrightarrow{AB}$ + 1/2 $\cdot$ $\overrightarrow{AC}$
  \item[] \textbf{D.} $\overrightarrow{AG}$ = 1/3 $\cdot$ $\overrightarrow{AB}$ + 2/3 $\cdot$ $\overrightarrow{AC}$
\end{itemize}\vspace{2pt}
\subsection*{\color{primary} Bảng đáp án Chương IV: Vectơ}
\begin{center}
\renewcommand{\arraystretch}{1.25}
\begin{tabularx}{\textwidth}{*{10}{>{\centering\arraybackslash}X}}
\toprule
\textbf{\color{primary} Cột 1} & \textbf{\color{primary} Cột 2} & \textbf{\color{primary} Cột 3} & \textbf{\color{primary} Cột 4} & \textbf{\color{primary} Cột 5} & \textbf{\color{primary} Cột 6} & \textbf{\color{primary} Cột 7} & \textbf{\color{primary} Cột 8} & \textbf{\color{primary} Cột 9} & \textbf{\color{primary} Cột 10} \\
\midrule
\textbf{1.C} & \textbf{2.C} & \textbf{3.C} & \textbf{4.B} & \textbf{5.B} & \textbf{6.A} & \textbf{7.D} & \textbf{8.B} & \textbf{9.B} & \textbf{10.A} \\
\rowcolor{rowtint} \textbf{11.B} & \textbf{12.D} & \textbf{13.C} & \textbf{14.A} & \textbf{15.C} & \textbf{16.A} & \textbf{17.C} & \textbf{18.B} & \textbf{19.D} & \textbf{20.C} \\
\textbf{21.B} & \textbf{22.C} & \textbf{23.B} & \textbf{24.A} & \textbf{25.A} & \textbf{26.C} & \textbf{27.A} & \textbf{28.D} & \textbf{29.B} & \textbf{30.B} \\
\rowcolor{rowtint} \textbf{31.A} & \textbf{32.D} & \textbf{33.C} & \textbf{34.D} & \textbf{35.A} & \textbf{36.D} & \textbf{37.A} & \textbf{38.A} & \textbf{39.D} & \textbf{40.D} \\
\textbf{41.C} & \textbf{42.C} & \textbf{43.A} & \textbf{44.B} & \textbf{45.B} & \textbf{46.A} & \textbf{47.D} & \textbf{48.B} & \textbf{49.C} & \textbf{50.B} \\
\bottomrule
\end{tabularx}
\end{center}\vspace{10pt}
\section*{\color{primary} CHƯƠNG V: CÁC SỐ ĐẶC TRƯNG CỦA MẪU SỐ LIỆU KHÔNG GHÉP NHÓM}
\textit{Trọng tâm kiến thức SGK: Số gần đúng và sai số, các số đặc trưng đo xu thế trung tâm (số trung bình, trung vị, tứ phân vị, mốt) và các số đặc trưng đo độ phân tán (khoảng biến thiên, khoảng tứ phân vị, phương sai, độ lệch chuẩn).}\par\vspace{4pt}
\subsection*{\color{accent} Phần I. Mức độ Dễ (Nhận biết -- Câu 1 đến Câu 20)}
\textbf{\color{primary} Câu 1.} Cho mẫu số liệu gồm 5 giá trị: 2, 4, 6, 8, 10. Số trung bình cộng $\bar{x}$ và trung vị $M_e$ của mẫu số liệu đó bằng:\par
\hspace*{1.5em}\textbf{A.} $\bar{x}$ = $M_e$ = 6\quad\quad \textbf{B.} $\bar{x}$ = 6; $M_e$ = 8\quad\quad \textbf{C.} $\bar{x}$ = 30; $M_e$ = 6\quad\quad \textbf{D.} $\bar{x}$ = 5; $M_e$ = 6\par\vspace{2pt}
\textbf{\color{primary} Câu 2.} Thực hiện yêu cầu nhận biết sau: Khoảng biến thiên R của mẫu số liệu: 6, 13, 4, 20, 9 bằng.\par
\hspace*{1.5em}\textbf{A.} 16\quad\quad \textbf{B.} 20\quad\quad \textbf{C.} 14\quad\quad \textbf{D.} 17\par\vspace{2pt}
\textbf{\color{primary} Câu 3.} Trong bài học thuộc Chương V, xét nội dung: Mốt ($M_o$) của mẫu số liệu: 4; 5; 8; 5; 8; 8; 8 là. Đáp án đúng là:\par
\hspace*{1.5em}\textbf{A.} 8 (1)\quad\quad \textbf{B.} 8\quad\quad \textbf{C.} 5\quad\quad \textbf{D.} 4\par\vspace{2pt}
\textbf{\color{primary} Câu 4.} Dựa vào kiến thức chuẩn SGK của Chương V (Các số đặc trưng của mẫu số liệu không ghép nhóm), cho mẫu số liệu gồm 5 giá trị: 5, 7, 9, 11, 13. Số trung bình cộng $\bar{x}$ và trung vị $M_e$ của mẫu số liệu đó bằng:\par
\hspace*{1.5em}\textbf{A.} $\bar{x}$ = 9; $M_e$ = 11\quad\quad \textbf{B.} $\bar{x}$ = $M_e$ = 9\quad\quad \textbf{C.} $\bar{x}$ = 8; $M_e$ = 9\quad\quad \textbf{D.} $\bar{x}$ = 45; $M_e$ = 9\par\vspace{2pt}
\textbf{\color{primary} Câu 5.} Chọn kết quả đúng nhất khi xét: Khoảng biến thiên R của mẫu số liệu: 9, 16, 7, 23, 12 bằng:\par
\hspace*{1.5em}\textbf{A.} 14\quad\quad \textbf{B.} 17\quad\quad \textbf{C.} 16\quad\quad \textbf{D.} 23\par\vspace{2pt}
\textbf{\color{primary} Câu 6.} Khi làm bài tập cơ bản về các số đặc trưng của mẫu số liệu không ghép nhóm, một học sinh gặp câu hỏi: Mốt ($M_o$) của mẫu số liệu: 4; 5; 11; 5; 11; 11; 8 là. Kết quả nào dưới đây đúng?\par
\hspace*{1.5em}\textbf{A.} 5\quad\quad \textbf{B.} 4\quad\quad \textbf{C.} 8\quad\quad \textbf{D.} 11\par\vspace{2pt}
\textbf{\color{primary} Câu 7.} Xác định giá trị hoặc khẳng định đúng cho bài toán sau: Cho mẫu số liệu gồm 5 giá trị: 8, 10, 12, 14, 16. Số trung bình cộng $\bar{x}$ và trung vị $M_e$ của mẫu số liệu đó bằng:\par
\hspace*{1.5em}\textbf{A.} $\bar{x}$ = 12; $M_e$ = 14\quad\quad \textbf{B.} $\bar{x}$ = 11; $M_e$ = 12\quad\quad \textbf{C.} $\bar{x}$ = $M_e$ = 12\quad\quad \textbf{D.} $\bar{x}$ = 60; $M_e$ = 12\par\vspace{2pt}
\textbf{\color{primary} Câu 8.} Áp dụng trực tiếp định nghĩa và quy tắc trong Chương V: Khoảng biến thiên R của mẫu số liệu: 12, 19, 10, 26, 15 bằng:\par
\hspace*{1.5em}\textbf{A.} 14\quad\quad \textbf{B.} 16\quad\quad \textbf{C.} 26\quad\quad \textbf{D.} 17\par\vspace{2pt}
\textbf{\color{primary} Câu 9.} Kiểm tra nhanh kiến thức nền tảng: Mốt ($M_o$) của mẫu số liệu: 4; 5; 14; 5; 14; 14; 8 là:\par
\hspace*{1.5em}\textbf{A.} 5\quad\quad \textbf{B.} 8\quad\quad \textbf{C.} 4\quad\quad \textbf{D.} 14\par\vspace{2pt}
\textbf{\color{primary} Câu 10.} Cho dữ kiện nhận biết sau đây: Cho mẫu số liệu gồm 5 giá trị: 11, 13, 15, 17, 19. Số trung bình cộng $\bar{x}$ và trung vị $M_e$ của mẫu số liệu đó bằng. Phương án chính xác là:\par
\hspace*{1.5em}\textbf{A.} $\bar{x}$ = 15; $M_e$ = 17\quad\quad \textbf{B.} $\bar{x}$ = $M_e$ = 15\quad\quad \textbf{C.} $\bar{x}$ = 75; $M_e$ = 15\quad\quad \textbf{D.} $\bar{x}$ = 14; $M_e$ = 15\par\vspace{2pt}
\textbf{\color{primary} Câu 11.} Hãy chỉ ra kết quả đúng trong bốn lựa chọn A, B, C, D đối với yêu cầu: Khoảng biến thiên R của mẫu số liệu: 15, 22, 13, 29, 18 bằng:\par
\hspace*{1.5em}\textbf{A.} 16\quad\quad \textbf{B.} 29\quad\quad \textbf{C.} 17\quad\quad \textbf{D.} 14\par\vspace{2pt}
\textbf{\color{primary} Câu 12.} Câu hỏi ôn tập lý thuyết và kỹ năng cơ bản (Chương V): Mốt ($M_o$) của mẫu số liệu: 4; 5; 17; 5; 17; 17; 8 là:\par
\hspace*{1.5em}\textbf{A.} 5\quad\quad \textbf{B.} 4\quad\quad \textbf{C.} 8\quad\quad \textbf{D.} 17\par\vspace{2pt}
\textbf{\color{primary} Câu 13.} Theo nội dung sách giáo khoa Kết nối tri thức với cuộc sống (Chương V), cho mẫu số liệu gồm 5 giá trị: 14, 16, 18, 20, 22. Số trung bình cộng $\bar{x}$ và trung vị $M_e$ của mẫu số liệu đó bằng:\par
\hspace*{1.5em}\textbf{A.} $\bar{x}$ = $M_e$ = 18\quad\quad \textbf{B.} $\bar{x}$ = 90; $M_e$ = 18\quad\quad \textbf{C.} $\bar{x}$ = 17; $M_e$ = 18\quad\quad \textbf{D.} $\bar{x}$ = 18; $M_e$ = 20\par\vspace{2pt}
\textbf{\color{primary} Câu 14.} Tìm đáp số đúng của bài toán cơ bản sau: Khoảng biến thiên R của mẫu số liệu: 18, 25, 16, 32, 21 bằng:\par
\hspace*{1.5em}\textbf{A.} 17\quad\quad \textbf{B.} 16\quad\quad \textbf{C.} 32\quad\quad \textbf{D.} 14\par\vspace{2pt}
\textbf{\color{primary} Câu 15.} Xét đối tượng toán học trong Chương V: Mốt ($M_o$) của mẫu số liệu: 4; 5; 20; 5; 20; 20; 8 là. Giá trị hoặc tính chất tương ứng là:\par
\hspace*{1.5em}\textbf{A.} 20\quad\quad \textbf{B.} 8\quad\quad \textbf{C.} 4\quad\quad \textbf{D.} 5\par\vspace{2pt}
\textbf{\color{primary} Câu 16.} Hoàn thành phép xác định cơ bản sau đây: Cho mẫu số liệu gồm 5 giá trị: 17, 19, 21, 23, 25. Số trung bình cộng $\bar{x}$ và trung vị $M_e$ của mẫu số liệu đó bằng:\par
\hspace*{1.5em}\textbf{A.} $\bar{x}$ = 20; $M_e$ = 21\quad\quad \textbf{B.} $\bar{x}$ = $M_e$ = 21\quad\quad \textbf{C.} $\bar{x}$ = 105; $M_e$ = 21\quad\quad \textbf{D.} $\bar{x}$ = 21; $M_e$ = 23\par\vspace{2pt}
\textbf{\color{primary} Câu 17.} Trong các khẳng định hoặc giá trị dưới đây, đâu là đáp án đúng cho câu hỏi: Khoảng biến thiên R của mẫu số liệu: 21, 28, 19, 35, 24 bằng?\par
\hspace*{1.5em}\textbf{A.} 35\quad\quad \textbf{B.} 16\quad\quad \textbf{C.} 17\quad\quad \textbf{D.} 14\par\vspace{2pt}
\textbf{\color{primary} Câu 18.} Bài tập khởi động Chương V: Mốt ($M_o$) của mẫu số liệu: 4; 5; 23; 5; 23; 23; 8 là:\par
\hspace*{1.5em}\textbf{A.} 4\quad\quad \textbf{B.} 5\quad\quad \textbf{C.} 8\quad\quad \textbf{D.} 23\par\vspace{2pt}
\textbf{\color{primary} Câu 19.} Quan sát dữ kiện đề bài: Cho mẫu số liệu gồm 5 giá trị: 20, 22, 24, 26, 28. Số trung bình cộng $\bar{x}$ và trung vị $M_e$ của mẫu số liệu đó bằng. Hãy chọn phương án đúng:\par
\hspace*{1.5em}\textbf{A.} $\bar{x}$ = 24; $M_e$ = 26\quad\quad \textbf{B.} $\bar{x}$ = 23; $M_e$ = 24\quad\quad \textbf{C.} $\bar{x}$ = $M_e$ = 24\quad\quad \textbf{D.} $\bar{x}$ = 120; $M_e$ = 24\par\vspace{2pt}
\textbf{\color{primary} Câu 20.} Sử dụng công thức hoặc tính chất đầu bài của Chương V để giải quyết yêu cầu: Khoảng biến thiên R của mẫu số liệu: 24, 31, 22, 38, 27 bằng:\par
\hspace*{1.5em}\textbf{A.} 38\quad\quad \textbf{B.} 16\quad\quad \textbf{C.} 14\quad\quad \textbf{D.} 17\par\vspace{2pt}
\subsection*{\color{accent} Phần II. Mức độ Trung bình (Thông hiểu -- Câu 21 đến Câu 40)}
\textbf{\color{primary} Câu 21.} Quy tròn số gần đúng a = 12,4682 với độ chính xác d = 0,03 (quy tròn đến hàng phần mười) ta được:\par
\hspace*{1.5em}\textbf{A.} 12,4\quad\quad \textbf{B.} 12,47\quad\quad \textbf{C.} 12\quad\quad \textbf{D.} 12,5\par\vspace{2pt}
\textbf{\color{primary} Câu 22.} Vận dụng quy tắc trong Chương V (Các số đặc trưng của mẫu số liệu không ghép nhóm) để giải bài toán: Cho mẫu số liệu đã sắp xếp tăng dần gồm 7 số: 2; 4; 5; 7; 9; 10; 13. Khoảng tứ phân vị $\Delta_Q$ = Q$_3$ - Q$_1$ bằng:\par
\hspace*{1.5em}\textbf{A.} 6\quad\quad \textbf{B.} 4\quad\quad \textbf{C.} 5\quad\quad \textbf{D.} 11\par\vspace{2pt}
\textbf{\color{primary} Câu 23.} Biến đổi và tính toán theo kiến thức Chương V: Nếu phương sai của một mẫu số liệu bằng s$^2$ = 16 thì độ lệch chuẩn s của mẫu số liệu đó bằng. Kết quả thu được là:\par
\hspace*{1.5em}\textbf{A.} s = 256\quad\quad \textbf{B.} s = $\pm$4\quad\quad \textbf{C.} s = 4\quad\quad \textbf{D.} s = 2.0\par\vspace{2pt}
\textbf{\color{primary} Câu 24.} Xét bài toán mức độ thông hiểu sau: Quy tròn số gần đúng a = 15,4682 với độ chính xác d = 0,03 (quy tròn đến hàng phần mười) ta được. Đáp án đúng là:\par
\hspace*{1.5em}\textbf{A.} 15,4\quad\quad \textbf{B.} 15,47\quad\quad \textbf{C.} 15,5\quad\quad \textbf{D.} 15\par\vspace{2pt}
\textbf{\color{primary} Câu 25.} Để hoàn thành phiếu học tập Chương V, hãy xác định: Cho mẫu số liệu đã sắp xếp tăng dần gồm 7 số: 5; 7; 8; 10; 12; 13; 16. Khoảng tứ phân vị $\Delta_Q$ = Q$_3$ - Q$_1$ bằng:\par
\hspace*{1.5em}\textbf{A.} 6\quad\quad \textbf{B.} 11\quad\quad \textbf{C.} 4\quad\quad \textbf{D.} 5\par\vspace{2pt}
\textbf{\color{primary} Câu 26.} Phân tích dữ kiện bài toán: Nếu phương sai của một mẫu số liệu bằng s$^2$ = 49 thì độ lệch chuẩn s của mẫu số liệu đó bằng. Giá trị cần tìm bằng:\par
\hspace*{1.5em}\textbf{A.} s = 7\quad\quad \textbf{B.} s = $\pm$7\quad\quad \textbf{C.} s = 2401\quad\quad \textbf{D.} s = 3.5\par\vspace{2pt}
\textbf{\color{primary} Câu 27.} Áp dụng tính chất đã học ở Chương V vào bài tập sau: Quy tròn số gần đúng a = 18,4682 với độ chính xác d = 0,03 (quy tròn đến hàng phần mười) ta được:\par
\hspace*{1.5em}\textbf{A.} 18\quad\quad \textbf{B.} 18,4\quad\quad \textbf{C.} 18,5\quad\quad \textbf{D.} 18,47\par\vspace{2pt}
\textbf{\color{primary} Câu 28.} Thực hiện các bước suy luận hợp lý để giải quyết yêu cầu: Cho mẫu số liệu đã sắp xếp tăng dần gồm 7 số: 8; 10; 11; 13; 15; 16; 19. Khoảng tứ phân vị $\Delta_Q$ = Q$_3$ - Q$_1$ bằng:\par
\hspace*{1.5em}\textbf{A.} 6\quad\quad \textbf{B.} 5\quad\quad \textbf{C.} 4\quad\quad \textbf{D.} 11\par\vspace{2pt}
\textbf{\color{primary} Câu 29.} Cho bài toán tự luyện (Các số đặc trưng của mẫu số liệu không ghép nhóm): Nếu phương sai của một mẫu số liệu bằng s$^2$ = 100 thì độ lệch chuẩn s của mẫu số liệu đó bằng. Phương án nào sau đây chính xác?\par
\hspace*{1.5em}\textbf{A.} s = 5.0\quad\quad \textbf{B.} s = $\pm$10\quad\quad \textbf{C.} s = 10\quad\quad \textbf{D.} s = 10000\par\vspace{2pt}
\textbf{\color{primary} Câu 30.} Kết hợp các phép biến đổi trong Chương V, hãy tìm kết quả của bài toán: Quy tròn số gần đúng a = 21,4682 với độ chính xác d = 0,03 (quy tròn đến hàng phần mười) ta được:\par
\hspace*{1.5em}\textbf{A.} 21,47\quad\quad \textbf{B.} 21,5\quad\quad \textbf{C.} 21\quad\quad \textbf{D.} 21,4\par\vspace{2pt}
\textbf{\color{primary} Câu 31.} Một bài kiểm tra định kỳ nội dung Chương V nêu câu hỏi: Cho mẫu số liệu đã sắp xếp tăng dần gồm 7 số: 11; 13; 14; 16; 18; 19; 22. Khoảng tứ phân vị $\Delta_Q$ = Q$_3$ - Q$_1$ bằng. Hãy chọn đáp số đúng:\par
\hspace*{1.5em}\textbf{A.} 5\quad\quad \textbf{B.} 4\quad\quad \textbf{C.} 11\quad\quad \textbf{D.} 6\par\vspace{2pt}
\textbf{\color{primary} Câu 32.} Tìm lời giải đúng cho bài toán thông hiểu: Nếu phương sai của một mẫu số liệu bằng s$^2$ = 169 thì độ lệch chuẩn s của mẫu số liệu đó bằng:\par
\hspace*{1.5em}\textbf{A.} s = 28561\quad\quad \textbf{B.} s = $\pm$13\quad\quad \textbf{C.} s = 13\quad\quad \textbf{D.} s = 6.5\par\vspace{2pt}
\textbf{\color{primary} Câu 33.} Sử dụng mối liên hệ giữa các đại lượng trong Chương V: Quy tròn số gần đúng a = 24,4682 với độ chính xác d = 0,03 (quy tròn đến hàng phần mười) ta được:\par
\hspace*{1.5em}\textbf{A.} 24,47\quad\quad \textbf{B.} 24,4\quad\quad \textbf{C.} 24\quad\quad \textbf{D.} 24,5\par\vspace{2pt}
\textbf{\color{primary} Câu 34.} Hãy tính toán và chọn phương án đúng đối với bài tập: Cho mẫu số liệu đã sắp xếp tăng dần gồm 7 số: 14; 16; 17; 19; 21; 22; 25. Khoảng tứ phân vị $\Delta_Q$ = Q$_3$ - Q$_1$ bằng:\par
\hspace*{1.5em}\textbf{A.} 5\quad\quad \textbf{B.} 4\quad\quad \textbf{C.} 11\quad\quad \textbf{D.} 6\par\vspace{2pt}
\textbf{\color{primary} Câu 35.} Dựa trên các tính chất trọng tâm của Chương V, xét bài toán: Nếu phương sai của một mẫu số liệu bằng s$^2$ = 256 thì độ lệch chuẩn s của mẫu số liệu đó bằng:\par
\hspace*{1.5em}\textbf{A.} s = 16\quad\quad \textbf{B.} s = 8.0\quad\quad \textbf{C.} s = 65536\quad\quad \textbf{D.} s = $\pm$16\par\vspace{2pt}
\textbf{\color{primary} Câu 36.} Thực hành giải toán (Các số đặc trưng của mẫu số liệu không ghép nhóm): Quy tròn số gần đúng a = 27,4682 với độ chính xác d = 0,03 (quy tròn đến hàng phần mười) ta được. Đáp án nào dưới đây thỏa mãn?\par
\hspace*{1.5em}\textbf{A.} 27,5\quad\quad \textbf{B.} 27\quad\quad \textbf{C.} 27,4\quad\quad \textbf{D.} 27,47\par\vspace{2pt}
\textbf{\color{primary} Câu 37.} Kiểm tra kỹ năng biến đổi và suy luận qua bài toán: Cho mẫu số liệu đã sắp xếp tăng dần gồm 7 số: 17; 19; 20; 22; 24; 25; 28. Khoảng tứ phân vị $\Delta_Q$ = Q$_3$ - Q$_1$ bằng:\par
\hspace*{1.5em}\textbf{A.} 6\quad\quad \textbf{B.} 11\quad\quad \textbf{C.} 5\quad\quad \textbf{D.} 4\par\vspace{2pt}
\textbf{\color{primary} Câu 38.} Cho giả thiết thuộc phạm vi Chương V: Nếu phương sai của một mẫu số liệu bằng s$^2$ = 361 thì độ lệch chuẩn s của mẫu số liệu đó bằng. Kết luận đúng là:\par
\hspace*{1.5em}\textbf{A.} s = 130321\quad\quad \textbf{B.} s = 19\quad\quad \textbf{C.} s = 9.5\quad\quad \textbf{D.} s = $\pm$19\par\vspace{2pt}
\textbf{\color{primary} Câu 39.} Rút gọn hoặc tính giá trị theo yêu cầu sau của Chương V: Quy tròn số gần đúng a = 30,4682 với độ chính xác d = 0,03 (quy tròn đến hàng phần mười) ta được:\par
\hspace*{1.5em}\textbf{A.} 30,4\quad\quad \textbf{B.} 30,47\quad\quad \textbf{C.} 30\quad\quad \textbf{D.} 30,5\par\vspace{2pt}
\textbf{\color{primary} Câu 40.} Xác định phương án đúng sau khi thực hiện phép tính / suy luận: Cho mẫu số liệu đã sắp xếp tăng dần gồm 7 số: 20; 22; 23; 25; 27; 28; 31. Khoảng tứ phân vị $\Delta_Q$ = Q$_3$ - Q$_1$ bằng:\par
\hspace*{1.5em}\textbf{A.} 4\quad\quad \textbf{B.} 11\quad\quad \textbf{C.} 5\quad\quad \textbf{D.} 6\par\vspace{2pt}
\subsection*{\color{accent} Phần III. Mức độ Khó (Vận dụng \& Vận dụng cao -- Câu 41 đến Câu 50)}
\textbf{\color{primary} Câu 41.} Cho mẫu số liệu gồm 4 số: 2, 4, 6, 8. Phương sai s$^2$ của mẫu số liệu này bằng:\par
\hspace*{1.5em}\textbf{A.} 5\quad\quad \textbf{B.} 2,5\quad\quad \textbf{C.} 4\quad\quad \textbf{D.} $\sqrt{5}$\par\vspace{2pt}
\textbf{\color{primary} Câu 42.} Bài toán vận dụng (Chương V – Các số đặc trưng của mẫu số liệu không ghép nhóm): Trong một mẫu số liệu có Q$_1$ = 12 và Q$_3$ = 20 (khoảng tứ phân vị $\Delta_Q$ = 8). Một giá trị x được coi là giá trị ngoại lệ (bất thường) nếu x > Q$_3$ + 1,5$\cdot$$\Delta_Q$, tức là khi. Kết quả đúng là:\par
\hspace*{1.5em}\textbf{A.} x > 24\quad\quad \textbf{B.} x > 36\quad\quad \textbf{C.} x > 32\quad\quad \textbf{D.} x > 28\par\vspace{2pt}
\textbf{\color{primary} Câu 43.} Vận dụng kiến thức nâng cao trong phạm vi Chương V để giải bài toán: Cho mẫu số liệu gồm 4 số: 4, 6, 8, 10. Phương sai s$^2$ của mẫu số liệu này bằng:\par
\hspace*{1.5em}\textbf{A.} $\sqrt{5}$\quad\quad \textbf{B.} 2,5\quad\quad \textbf{C.} 5\quad\quad \textbf{D.} 4\par\vspace{2pt}
\textbf{\color{primary} Câu 44.} Phân tích và giải bài toán tổng hợp sau của Chương V: Trong một mẫu số liệu có Q$_1$ = 16 và Q$_3$ = 24 (khoảng tứ phân vị $\Delta_Q$ = 8). Một giá trị x được coi là giá trị ngoại lệ (bất thường) nếu x > Q$_3$ + 1,5$\cdot$$\Delta_Q$, tức là khi:\par
\hspace*{1.5em}\textbf{A.} x > 40\quad\quad \textbf{B.} x > 32\quad\quad \textbf{C.} x > 28\quad\quad \textbf{D.} x > 36\par\vspace{2pt}
\textbf{\color{primary} Câu 45.} Thử thách tư duy trong phạm vi kiến thức Chương V: Cho mẫu số liệu gồm 4 số: 6, 8, 10, 12. Phương sai s$^2$ của mẫu số liệu này bằng. Đáp số chính xác là:\par
\hspace*{1.5em}\textbf{A.} 5\quad\quad \textbf{B.} 4\quad\quad \textbf{C.} 2,5\quad\quad \textbf{D.} $\sqrt{5}$\par\vspace{2pt}
\textbf{\color{primary} Câu 46.} Sử dụng phương pháp giải toán của Chương V để tìm đáp án cho bài toán: Trong một mẫu số liệu có Q$_1$ = 20 và Q$_3$ = 28 (khoảng tứ phân vị $\Delta_Q$ = 8). Một giá trị x được coi là giá trị ngoại lệ (bất thường) nếu x > Q$_3$ + 1,5$\cdot$$\Delta_Q$, tức là khi:\par
\hspace*{1.5em}\textbf{A.} x > 44\quad\quad \textbf{B.} x > 32\quad\quad \textbf{C.} x > 36\quad\quad \textbf{D.} x > 40\par\vspace{2pt}
\textbf{\color{primary} Câu 47.} Xét bài toán vận dụng thực tiễn / nâng cao thuộc Chương V: Cho mẫu số liệu gồm 4 số: 8, 10, 12, 14. Phương sai s$^2$ của mẫu số liệu này bằng:\par
\hspace*{1.5em}\textbf{A.} 4\quad\quad \textbf{B.} 5\quad\quad \textbf{C.} 2,5\quad\quad \textbf{D.} $\sqrt{5}$\par\vspace{2pt}
\textbf{\color{primary} Câu 48.} Tìm kết quả cuối cùng của bài toán phân loại học sinh khá giỏi (Chương V): Trong một mẫu số liệu có Q$_1$ = 24 và Q$_3$ = 32 (khoảng tứ phân vị $\Delta_Q$ = 8). Một giá trị x được coi là giá trị ngoại lệ (bất thường) nếu x > Q$_3$ + 1,5$\cdot$$\Delta_Q$, tức là khi:\par
\hspace*{1.5em}\textbf{A.} x > 48\quad\quad \textbf{B.} x > 36\quad\quad \textbf{C.} x > 44\quad\quad \textbf{D.} x > 40\par\vspace{2pt}
\textbf{\color{primary} Câu 49.} Kết hợp các tính chất đã học trong Chương V, hãy giải quyết bài toán: Cho mẫu số liệu gồm 4 số: 10, 12, 14, 16. Phương sai s$^2$ của mẫu số liệu này bằng:\par
\hspace*{1.5em}\textbf{A.} 5\quad\quad \textbf{B.} 2,5\quad\quad \textbf{C.} $\sqrt{5}$\quad\quad \textbf{D.} 4\par\vspace{2pt}
\textbf{\color{primary} Câu 50.} Câu hỏi vận dụng cao cuối Chương V: Trong một mẫu số liệu có Q$_1$ = 28 và Q$_3$ = 36 (khoảng tứ phân vị $\Delta_Q$ = 8). Một giá trị x được coi là giá trị ngoại lệ (bất thường) nếu x > Q$_3$ + 1,5$\cdot$$\Delta_Q$, tức là khi. Phương án đúng là:\par
\hspace*{1.5em}\textbf{A.} x > 44\quad\quad \textbf{B.} x > 48\quad\quad \textbf{C.} x > 52\quad\quad \textbf{D.} x > 40\par\vspace{2pt}
\subsection*{\color{primary} Bảng đáp án Chương V: Các số đặc trưng của mẫu số liệu không ghép nhóm}
\begin{center}
\renewcommand{\arraystretch}{1.25}
\begin{tabularx}{\textwidth}{*{10}{>{\centering\arraybackslash}X}}
\toprule
\textbf{\color{primary} Cột 1} & \textbf{\color{primary} Cột 2} & \textbf{\color{primary} Cột 3} & \textbf{\color{primary} Cột 4} & \textbf{\color{primary} Cột 5} & \textbf{\color{primary} Cột 6} & \textbf{\color{primary} Cột 7} & \textbf{\color{primary} Cột 8} & \textbf{\color{primary} Cột 9} & \textbf{\color{primary} Cột 10} \\
\midrule
\textbf{1.A} & \textbf{2.A} & \textbf{3.B} & \textbf{4.B} & \textbf{5.C} & \textbf{6.D} & \textbf{7.C} & \textbf{8.B} & \textbf{9.D} & \textbf{10.B} \\
\rowcolor{rowtint} \textbf{11.A} & \textbf{12.D} & \textbf{13.A} & \textbf{14.B} & \textbf{15.A} & \textbf{16.B} & \textbf{17.B} & \textbf{18.D} & \textbf{19.C} & \textbf{20.B} \\
\textbf{21.D} & \textbf{22.A} & \textbf{23.C} & \textbf{24.C} & \textbf{25.A} & \textbf{26.A} & \textbf{27.C} & \textbf{28.A} & \textbf{29.C} & \textbf{30.B} \\
\rowcolor{rowtint} \textbf{31.D} & \textbf{32.C} & \textbf{33.D} & \textbf{34.D} & \textbf{35.A} & \textbf{36.A} & \textbf{37.A} & \textbf{38.B} & \textbf{39.D} & \textbf{40.D} \\
\textbf{41.A} & \textbf{42.C} & \textbf{43.C} & \textbf{44.D} & \textbf{45.A} & \textbf{46.D} & \textbf{47.B} & \textbf{48.C} & \textbf{49.A} & \textbf{50.B} \\
\bottomrule
\end{tabularx}
\end{center}\vspace{10pt}
\section*{\color{primary} CHƯƠNG VI: HÀM SỐ, ĐỒ THỊ VÀ ỨNG DỤNG}
\textit{Trọng tâm kiến thức SGK: Hàm số, tập xác định và sự biến thiên, hàm số bậc hai y = ax$^2$ + bx + c (a $\ne$ 0), dấu của tam thức bậc hai và phương trình quy về phương trình bậc hai.}\par\vspace{4pt}
\subsection*{\color{accent} Phần I. Mức độ Dễ (Nhận biết -- Câu 1 đến Câu 20)}
\textbf{\color{primary} Câu 1.} Tập xác định D của hàm số y = (2x + 1)/(x - 3) là:\par
\hspace*{1.5em}\textbf{A.} D = $\mathbb{R}$ \textbackslash\{\} \{3\}\quad\quad \textbf{B.} D = $\mathbb{R}$ \textbackslash\{\} \{-3\}\quad\quad \textbf{C.} D = $\mathbb{R}$\quad\quad \textbf{D.} D = (3; +$\infty$)\par\vspace{2pt}
\textbf{\color{primary} Câu 2.} Thực hiện yêu cầu nhận biết sau: Hoành độ đỉnh và trục đối xứng của parabol (P): y = x$^2$ - 6x + 6 là.\par
\hspace*{1.5em}\textbf{A.} x = 3\quad\quad \textbf{B.} x = -3\quad\quad \textbf{C.} x = 6\quad\quad \textbf{D.} x = -(i + 2)\par\vspace{2pt}
\textbf{\color{primary} Câu 3.} Trong bài học thuộc Chương VI, xét nội dung: Hàm số bậc hai y = x$^2$ - 10x + 1 đồng biến trên khoảng nào sau đây. Đáp án đúng là:\par
\hspace*{1.5em}\textbf{A.} (-5; +$\infty$)\quad\quad \textbf{B.} (5; +$\infty$)\quad\quad \textbf{C.} (-$\infty$; +$\infty$)\quad\quad \textbf{D.} (-$\infty$; 5)\par\vspace{2pt}
\textbf{\color{primary} Câu 4.} Dựa vào kiến thức chuẩn SGK của Chương VI (Hàm số, đồ thị và ứng dụng), tập xác định D của hàm số y = (2x + 1)/(x - 6) là:\par
\hspace*{1.5em}\textbf{A.} D = $\mathbb{R}$ \textbackslash\{\} \{-6\}\quad\quad \textbf{B.} D = $\mathbb{R}$ \textbackslash\{\} \{6\}\quad\quad \textbf{C.} D = (6; +$\infty$)\quad\quad \textbf{D.} D = $\mathbb{R}$\par\vspace{2pt}
\textbf{\color{primary} Câu 5.} Chọn kết quả đúng nhất khi xét: Hoành độ đỉnh và trục đối xứng của parabol (P): y = x$^2$ - 12x + 9 là:\par
\hspace*{1.5em}\textbf{A.} x = 12\quad\quad \textbf{B.} x = 6\quad\quad \textbf{C.} x = -(i + 2)\quad\quad \textbf{D.} x = -6\par\vspace{2pt}
\textbf{\color{primary} Câu 6.} Khi làm bài tập cơ bản về hàm số, đồ thị và ứng dụng, một học sinh gặp câu hỏi: Hàm số bậc hai y = x$^2$ - 16x + 1 đồng biến trên khoảng nào sau đây. Kết quả nào dưới đây đúng?\par
\hspace*{1.5em}\textbf{A.} (-$\infty$; 8)\quad\quad \textbf{B.} (8; +$\infty$)\quad\quad \textbf{C.} (-$\infty$; +$\infty$)\quad\quad \textbf{D.} (-8; +$\infty$)\par\vspace{2pt}
\textbf{\color{primary} Câu 7.} Xác định giá trị hoặc khẳng định đúng cho bài toán sau: Tập xác định D của hàm số y = (2x + 1)/(x - 9) là:\par
\hspace*{1.5em}\textbf{A.} D = $\mathbb{R}$\quad\quad \textbf{B.} D = $\mathbb{R}$ \textbackslash\{\} \{-9\}\quad\quad \textbf{C.} D = (9; +$\infty$)\quad\quad \textbf{D.} D = $\mathbb{R}$ \textbackslash\{\} \{9\}\par\vspace{2pt}
\textbf{\color{primary} Câu 8.} Áp dụng trực tiếp định nghĩa và quy tắc trong Chương VI: Hoành độ đỉnh và trục đối xứng của parabol (P): y = x$^2$ - 18x + 12 là:\par
\hspace*{1.5em}\textbf{A.} x = 18\quad\quad \textbf{B.} x = -(i + 2)\quad\quad \textbf{C.} x = 9\quad\quad \textbf{D.} x = -9\par\vspace{2pt}
\textbf{\color{primary} Câu 9.} Kiểm tra nhanh kiến thức nền tảng: Hàm số bậc hai y = x$^2$ - 22x + 1 đồng biến trên khoảng nào sau đây?\par
\hspace*{1.5em}\textbf{A.} (-$\infty$; 11)\quad\quad \textbf{B.} (-$\infty$; +$\infty$)\quad\quad \textbf{C.} (-11; +$\infty$)\quad\quad \textbf{D.} (11; +$\infty$)\par\vspace{2pt}
\textbf{\color{primary} Câu 10.} Cho dữ kiện nhận biết sau đây: Tập xác định D của hàm số y = (2x + 1)/(x - 12) là. Phương án chính xác là:\par
\hspace*{1.5em}\textbf{A.} D = $\mathbb{R}$ \textbackslash\{\} \{12\}\quad\quad \textbf{B.} D = $\mathbb{R}$\quad\quad \textbf{C.} D = $\mathbb{R}$ \textbackslash\{\} \{-12\}\quad\quad \textbf{D.} D = (12; +$\infty$)\par\vspace{2pt}
\textbf{\color{primary} Câu 11.} Hãy chỉ ra kết quả đúng trong bốn lựa chọn A, B, C, D đối với yêu cầu: Hoành độ đỉnh và trục đối xứng của parabol (P): y = x$^2$ - 24x + 15 là:\par
\hspace*{1.5em}\textbf{A.} x = -(i + 2)\quad\quad \textbf{B.} x = -12\quad\quad \textbf{C.} x = 12\quad\quad \textbf{D.} x = 24\par\vspace{2pt}
\textbf{\color{primary} Câu 12.} Câu hỏi ôn tập lý thuyết và kỹ năng cơ bản (Chương VI): Hàm số bậc hai y = x$^2$ - 28x + 1 đồng biến trên khoảng nào sau đây:\par
\hspace*{1.5em}\textbf{A.} (-$\infty$; 14)\quad\quad \textbf{B.} (-$\infty$; +$\infty$)\quad\quad \textbf{C.} (-14; +$\infty$)\quad\quad \textbf{D.} (14; +$\infty$)\par\vspace{2pt}
\textbf{\color{primary} Câu 13.} Theo nội dung sách giáo khoa Kết nối tri thức với cuộc sống (Chương VI), tập xác định D của hàm số y = (2x + 1)/(x - 15) là:\par
\hspace*{1.5em}\textbf{A.} D = $\mathbb{R}$ \textbackslash\{\} \{-15\}\quad\quad \textbf{B.} D = $\mathbb{R}$\quad\quad \textbf{C.} D = (15; +$\infty$)\quad\quad \textbf{D.} D = $\mathbb{R}$ \textbackslash\{\} \{15\}\par\vspace{2pt}
\textbf{\color{primary} Câu 14.} Tìm đáp số đúng của bài toán cơ bản sau: Hoành độ đỉnh và trục đối xứng của parabol (P): y = x$^2$ - 30x + 18 là:\par
\hspace*{1.5em}\textbf{A.} x = 30\quad\quad \textbf{B.} x = 15\quad\quad \textbf{C.} x = -15\quad\quad \textbf{D.} x = -(i + 2)\par\vspace{2pt}
\textbf{\color{primary} Câu 15.} Xét đối tượng toán học trong Chương VI: Hàm số bậc hai y = x$^2$ - 34x + 1 đồng biến trên khoảng nào sau đây. Giá trị hoặc tính chất tương ứng là:\par
\hspace*{1.5em}\textbf{A.} (-$\infty$; +$\infty$)\quad\quad \textbf{B.} (-$\infty$; 17)\quad\quad \textbf{C.} (-17; +$\infty$)\quad\quad \textbf{D.} (17; +$\infty$)\par\vspace{2pt}
\textbf{\color{primary} Câu 16.} Hoàn thành phép xác định cơ bản sau đây: Tập xác định D của hàm số y = (2x + 1)/(x - 18) là:\par
\hspace*{1.5em}\textbf{A.} D = $\mathbb{R}$ \textbackslash\{\} \{-18\}\quad\quad \textbf{B.} D = (18; +$\infty$)\quad\quad \textbf{C.} D = $\mathbb{R}$ \textbackslash\{\} \{18\}\quad\quad \textbf{D.} D = $\mathbb{R}$\par\vspace{2pt}
\textbf{\color{primary} Câu 17.} Trong các khẳng định hoặc giá trị dưới đây, đâu là đáp án đúng cho câu hỏi: Hoành độ đỉnh và trục đối xứng của parabol (P): y = x$^2$ - 36x + 21 là?\par
\hspace*{1.5em}\textbf{A.} x = 36\quad\quad \textbf{B.} x = -(i + 2)\quad\quad \textbf{C.} x = -18\quad\quad \textbf{D.} x = 18\par\vspace{2pt}
\textbf{\color{primary} Câu 18.} Bài tập khởi động Chương VI: Hàm số bậc hai y = x$^2$ - 40x + 1 đồng biến trên khoảng nào sau đây:\par
\hspace*{1.5em}\textbf{A.} (20; +$\infty$)\quad\quad \textbf{B.} (-$\infty$; 20)\quad\quad \textbf{C.} (-20; +$\infty$)\quad\quad \textbf{D.} (-$\infty$; +$\infty$)\par\vspace{2pt}
\textbf{\color{primary} Câu 19.} Quan sát dữ kiện đề bài: Tập xác định D của hàm số y = (2x + 1)/(x - 21) là. Hãy chọn phương án đúng:\par
\hspace*{1.5em}\textbf{A.} D = $\mathbb{R}$ \textbackslash\{\} \{-21\}\quad\quad \textbf{B.} D = $\mathbb{R}$\quad\quad \textbf{C.} D = (21; +$\infty$)\quad\quad \textbf{D.} D = $\mathbb{R}$ \textbackslash\{\} \{21\}\par\vspace{2pt}
\textbf{\color{primary} Câu 20.} Sử dụng công thức hoặc tính chất đầu bài của Chương VI để giải quyết yêu cầu: Hoành độ đỉnh và trục đối xứng của parabol (P): y = x$^2$ - 42x + 24 là:\par
\hspace*{1.5em}\textbf{A.} x = 42\quad\quad \textbf{B.} x = -(i + 2)\quad\quad \textbf{C.} x = 21\quad\quad \textbf{D.} x = -21\par\vspace{2pt}
\subsection*{\color{accent} Phần II. Mức độ Trung bình (Thông hiểu -- Câu 21 đến Câu 40)}
\textbf{\color{primary} Câu 21.} Tập nghiệm của bất phương trình bậc hai x$^2$ - 5x + 4 $\le$ 0 là:\par
\hspace*{1.5em}\textbf{A.} S = [1; 4]\quad\quad \textbf{B.} S = (1; 4)\quad\quad \textbf{C.} S = (-$\infty$; 1] $\cup$ [4; +$\infty$)\quad\quad \textbf{D.} S = [-4; -1]\par\vspace{2pt}
\textbf{\color{primary} Câu 22.} Vận dụng quy tắc trong Chương VI (Hàm số, đồ thị và ứng dụng) để giải bài toán: Tập xác định của hàm số y = $\sqrt{x - 2}$ + 1/(x - 6) là:\par
\hspace*{1.5em}\textbf{A.} [2; +$\infty$) \textbackslash\{\} \{6\}\quad\quad \textbf{B.} [2; 6]\quad\quad \textbf{C.} (2; 6)\quad\quad \textbf{D.} $\mathbb{R}$ \textbackslash\{\} \{2; 6\}\par\vspace{2pt}
\textbf{\color{primary} Câu 23.} Biến đổi và tính toán theo kiến thức Chương VI: Tổng các nghiệm thực của phương trình chứa căn $\sqrt{x^2 + 9}$ = $\sqrt{2x + 12}$ bằng. Kết quả thu được là:\par
\hspace*{1.5em}\textbf{A.} -2\quad\quad \textbf{B.} 3\quad\quad \textbf{C.} 2\quad\quad \textbf{D.} 4\par\vspace{2pt}
\textbf{\color{primary} Câu 24.} Xét bài toán mức độ thông hiểu sau: Tập nghiệm của bất phương trình bậc hai x$^2$ - 11x + 28 $\le$ 0 là. Đáp án đúng là:\par
\hspace*{1.5em}\textbf{A.} S = [4; 7]\quad\quad \textbf{B.} S = (-$\infty$; 4] $\cup$ [7; +$\infty$)\quad\quad \textbf{C.} S = (4; 7)\quad\quad \textbf{D.} S = [-7; -4]\par\vspace{2pt}
\textbf{\color{primary} Câu 25.} Để hoàn thành phiếu học tập Chương VI, hãy xác định: Tập xác định của hàm số y = $\sqrt{x - 5}$ + 1/(x - 9) là:\par
\hspace*{1.5em}\textbf{A.} [5; 9]\quad\quad \textbf{B.} (5; 9)\quad\quad \textbf{C.} $\mathbb{R}$ \textbackslash\{\} \{5; 9\}\quad\quad \textbf{D.} [5; +$\infty$) \textbackslash\{\} \{9\}\par\vspace{2pt}
\textbf{\color{primary} Câu 26.} Phân tích dữ kiện bài toán: Tổng các nghiệm thực của phương trình chứa căn $\sqrt{x^2 + 15}$ = $\sqrt{2x + 18}$ bằng. Giá trị cần tìm bằng:\par
\hspace*{1.5em}\textbf{A.} 3\quad\quad \textbf{B.} 4\quad\quad \textbf{C.} 2\quad\quad \textbf{D.} -2\par\vspace{2pt}
\textbf{\color{primary} Câu 27.} Áp dụng tính chất đã học ở Chương VI vào bài tập sau: Tập nghiệm của bất phương trình bậc hai x$^2$ - 17x + 70 $\le$ 0 là:\par
\hspace*{1.5em}\textbf{A.} S = (-$\infty$; 7] $\cup$ [10; +$\infty$)\quad\quad \textbf{B.} S = (7; 10)\quad\quad \textbf{C.} S = [7; 10]\quad\quad \textbf{D.} S = [-10; -7]\par\vspace{2pt}
\textbf{\color{primary} Câu 28.} Thực hiện các bước suy luận hợp lý để giải quyết yêu cầu: Tập xác định của hàm số y = $\sqrt{x - 8}$ + 1/(x - 12) là:\par
\hspace*{1.5em}\textbf{A.} [8; 12]\quad\quad \textbf{B.} $\mathbb{R}$ \textbackslash\{\} \{8; 12\}\quad\quad \textbf{C.} [8; +$\infty$) \textbackslash\{\} \{12\}\quad\quad \textbf{D.} (8; 12)\par\vspace{2pt}
\textbf{\color{primary} Câu 29.} Cho bài toán tự luyện (Hàm số, đồ thị và ứng dụng): Tổng các nghiệm thực của phương trình chứa căn $\sqrt{x^2 + 21}$ = $\sqrt{2x + 24}$ bằng. Phương án nào sau đây chính xác?\par
\hspace*{1.5em}\textbf{A.} 4\quad\quad \textbf{B.} -2\quad\quad \textbf{C.} 2\quad\quad \textbf{D.} 3\par\vspace{2pt}
\textbf{\color{primary} Câu 30.} Kết hợp các phép biến đổi trong Chương VI, hãy tìm kết quả của bài toán: Tập nghiệm của bất phương trình bậc hai x$^2$ - 23x + 130 $\le$ 0 là:\par
\hspace*{1.5em}\textbf{A.} S = [-13; -10]\quad\quad \textbf{B.} S = (10; 13)\quad\quad \textbf{C.} S = (-$\infty$; 10] $\cup$ [13; +$\infty$)\quad\quad \textbf{D.} S = [10; 13]\par\vspace{2pt}
\textbf{\color{primary} Câu 31.} Một bài kiểm tra định kỳ nội dung Chương VI nêu câu hỏi: Tập xác định của hàm số y = $\sqrt{x - 11}$ + 1/(x - 15) là. Hãy chọn đáp số đúng:\par
\hspace*{1.5em}\textbf{A.} (11; 15)\quad\quad \textbf{B.} [11; 15]\quad\quad \textbf{C.} $\mathbb{R}$ \textbackslash\{\} \{11; 15\}\quad\quad \textbf{D.} [11; +$\infty$) \textbackslash\{\} \{15\}\par\vspace{2pt}
\textbf{\color{primary} Câu 32.} Tìm lời giải đúng cho bài toán thông hiểu: Tổng các nghiệm thực của phương trình chứa căn $\sqrt{x^2 + 27}$ = $\sqrt{2x + 30}$ bằng:\par
\hspace*{1.5em}\textbf{A.} -2\quad\quad \textbf{B.} 2\quad\quad \textbf{C.} 4\quad\quad \textbf{D.} 3\par\vspace{2pt}
\textbf{\color{primary} Câu 33.} Sử dụng mối liên hệ giữa các đại lượng trong Chương VI: Tập nghiệm của bất phương trình bậc hai x$^2$ - 29x + 208 $\le$ 0 là:\par
\hspace*{1.5em}\textbf{A.} S = [13; 16]\quad\quad \textbf{B.} S = [-16; -13]\quad\quad \textbf{C.} S = (13; 16)\quad\quad \textbf{D.} S = (-$\infty$; 13] $\cup$ [16; +$\infty$)\par\vspace{2pt}
\textbf{\color{primary} Câu 34.} Hãy tính toán và chọn phương án đúng đối với bài tập: Tập xác định của hàm số y = $\sqrt{x - 14}$ + 1/(x - 18) là:\par
\hspace*{1.5em}\textbf{A.} [14; 18]\quad\quad \textbf{B.} $\mathbb{R}$ \textbackslash\{\} \{14; 18\}\quad\quad \textbf{C.} (14; 18)\quad\quad \textbf{D.} [14; +$\infty$) \textbackslash\{\} \{18\}\par\vspace{2pt}
\textbf{\color{primary} Câu 35.} Dựa trên các tính chất trọng tâm của Chương VI, xét bài toán: Tổng các nghiệm thực của phương trình chứa căn $\sqrt{x^2 + 33}$ = $\sqrt{2x + 36}$ bằng:\par
\hspace*{1.5em}\textbf{A.} 4\quad\quad \textbf{B.} 3\quad\quad \textbf{C.} -2\quad\quad \textbf{D.} 2\par\vspace{2pt}
\textbf{\color{primary} Câu 36.} Thực hành giải toán (Hàm số, đồ thị và ứng dụng): Tập nghiệm của bất phương trình bậc hai x$^2$ - 35x + 304 $\le$ 0 là. Đáp án nào dưới đây thỏa mãn?\par
\hspace*{1.5em}\textbf{A.} S = (-$\infty$; 16] $\cup$ [19; +$\infty$)\quad\quad \textbf{B.} S = [-19; -16]\quad\quad \textbf{C.} S = [16; 19]\quad\quad \textbf{D.} S = (16; 19)\par\vspace{2pt}
\textbf{\color{primary} Câu 37.} Kiểm tra kỹ năng biến đổi và suy luận qua bài toán: Tập xác định của hàm số y = $\sqrt{x - 17}$ + 1/(x - 21) là:\par
\hspace*{1.5em}\textbf{A.} [17; +$\infty$) \textbackslash\{\} \{21\}\quad\quad \textbf{B.} (17; 21)\quad\quad \textbf{C.} [17; 21]\quad\quad \textbf{D.} $\mathbb{R}$ \textbackslash\{\} \{17; 21\}\par\vspace{2pt}
\textbf{\color{primary} Câu 38.} Cho giả thiết thuộc phạm vi Chương VI: Tổng các nghiệm thực của phương trình chứa căn $\sqrt{x^2 + 39}$ = $\sqrt{2x + 42}$ bằng. Kết luận đúng là:\par
\hspace*{1.5em}\textbf{A.} 4\quad\quad \textbf{B.} 2\quad\quad \textbf{C.} -2\quad\quad \textbf{D.} 3\par\vspace{2pt}
\textbf{\color{primary} Câu 39.} Rút gọn hoặc tính giá trị theo yêu cầu sau của Chương VI: Tập nghiệm của bất phương trình bậc hai x$^2$ - 41x + 418 $\le$ 0 là:\par
\hspace*{1.5em}\textbf{A.} S = [19; 22]\quad\quad \textbf{B.} S = [-22; -19]\quad\quad \textbf{C.} S = (19; 22)\quad\quad \textbf{D.} S = (-$\infty$; 19] $\cup$ [22; +$\infty$)\par\vspace{2pt}
\textbf{\color{primary} Câu 40.} Xác định phương án đúng sau khi thực hiện phép tính / suy luận: Tập xác định của hàm số y = $\sqrt{x - 20}$ + 1/(x - 24) là:\par
\hspace*{1.5em}\textbf{A.} (20; 24)\quad\quad \textbf{B.} [20; +$\infty$) \textbackslash\{\} \{24\}\quad\quad \textbf{C.} [20; 24]\quad\quad \textbf{D.} $\mathbb{R}$ \textbackslash\{\} \{20; 24\}\par\vspace{2pt}
\subsection*{\color{accent} Phần III. Mức độ Khó (Vận dụng \& Vận dụng cao -- Câu 41 đến Câu 50)}
\textbf{\color{primary} Câu 41.} Tìm tất cả các giá trị của tham số m để tam thức bậc hai f(x) = x$^2$ - 2mx + m + 2 luôn dương (f(x) > 0, $\forall$x $\in$ $\mathbb{R}$):\par
\hspace*{1.5em}\textbf{A.} -1 < m < 2\quad\quad \textbf{B.} m > 2\quad\quad \textbf{C.} -2 < m < 1\quad\quad \textbf{D.} m $\le$ -1 hoặc m $\ge$ 2\par\vspace{2pt}
\textbf{\color{primary} Câu 42.} Bài toán vận dụng (Chương VI – Hàm số, đồ thị và ứng dụng): Parabol (P): y = ax$^2$ + bx + 2 đi qua điểm A(1; 6) và có trục đối xứng là đường thẳng x = 1. Giá trị của tổng a + b bằng. Kết quả đúng là:\par
\hspace*{1.5em}\textbf{A.} 4\quad\quad \textbf{B.} 6\quad\quad \textbf{C.} 8\quad\quad \textbf{D.} -4\par\vspace{2pt}
\textbf{\color{primary} Câu 43.} Vận dụng kiến thức nâng cao trong phạm vi Chương VI để giải bài toán: Tìm tất cả các giá trị của tham số m để tam thức bậc hai f(x) = x$^2$ - 2mx + m + 2 luôn dương (f(x) > 0, $\forall$x $\in$ $\mathbb{R}$):\par
\hspace*{1.5em}\textbf{A.} m > 2\quad\quad \textbf{B.} -2 < m < 1\quad\quad \textbf{C.} m $\le$ -1 hoặc m $\ge$ 2\quad\quad \textbf{D.} -1 < m < 2\par\vspace{2pt}
\textbf{\color{primary} Câu 44.} Phân tích và giải bài toán tổng hợp sau của Chương VI: Parabol (P): y = ax$^2$ + bx + 2 đi qua điểm A(1; 8) và có trục đối xứng là đường thẳng x = 1. Giá trị của tổng a + b bằng:\par
\hspace*{1.5em}\textbf{A.} -6\quad\quad \textbf{B.} 6\quad\quad \textbf{C.} 12\quad\quad \textbf{D.} 8\par\vspace{2pt}
\textbf{\color{primary} Câu 45.} Thử thách tư duy trong phạm vi kiến thức Chương VI: Tìm tất cả các giá trị của tham số m để tam thức bậc hai f(x) = x$^2$ - 2mx + m + 2 luôn dương (f(x) > 0, $\forall$x $\in$ $\mathbb{R}$). Đáp số chính xác là:\par
\hspace*{1.5em}\textbf{A.} m > 2\quad\quad \textbf{B.} -1 < m < 2\quad\quad \textbf{C.} -2 < m < 1\quad\quad \textbf{D.} m $\le$ -1 hoặc m $\ge$ 2\par\vspace{2pt}
\textbf{\color{primary} Câu 46.} Sử dụng phương pháp giải toán của Chương VI để tìm đáp án cho bài toán: Parabol (P): y = ax$^2$ + bx + 2 đi qua điểm A(1; 10) và có trục đối xứng là đường thẳng x = 1. Giá trị của tổng a + b bằng:\par
\hspace*{1.5em}\textbf{A.} 16\quad\quad \textbf{B.} 8\quad\quad \textbf{C.} -8\quad\quad \textbf{D.} 10\par\vspace{2pt}
\textbf{\color{primary} Câu 47.} Xét bài toán vận dụng thực tiễn / nâng cao thuộc Chương VI: Tìm tất cả các giá trị của tham số m để tam thức bậc hai f(x) = x$^2$ - 2mx + m + 2 luôn dương (f(x) > 0, $\forall$x $\in$ $\mathbb{R}$):\par
\hspace*{1.5em}\textbf{A.} -2 < m < 1\quad\quad \textbf{B.} m $\le$ -1 hoặc m $\ge$ 2\quad\quad \textbf{C.} m > 2\quad\quad \textbf{D.} -1 < m < 2\par\vspace{2pt}
\textbf{\color{primary} Câu 48.} Tìm kết quả cuối cùng của bài toán phân loại học sinh khá giỏi (Chương VI): Parabol (P): y = ax$^2$ + bx + 2 đi qua điểm A(1; 12) và có trục đối xứng là đường thẳng x = 1. Giá trị của tổng a + b bằng:\par
\hspace*{1.5em}\textbf{A.} 20\quad\quad \textbf{B.} 12\quad\quad \textbf{C.} 10\quad\quad \textbf{D.} -10\par\vspace{2pt}
\textbf{\color{primary} Câu 49.} Kết hợp các tính chất đã học trong Chương VI, hãy giải quyết bài toán: Tìm tất cả các giá trị của tham số m để tam thức bậc hai f(x) = x$^2$ - 2mx + m + 2 luôn dương (f(x) > 0, $\forall$x $\in$ $\mathbb{R}$):\par
\hspace*{1.5em}\textbf{A.} -2 < m < 1\quad\quad \textbf{B.} m $\le$ -1 hoặc m $\ge$ 2\quad\quad \textbf{C.} m > 2\quad\quad \textbf{D.} -1 < m < 2\par\vspace{2pt}
\textbf{\color{primary} Câu 50.} Câu hỏi vận dụng cao cuối Chương VI: Parabol (P): y = ax$^2$ + bx + 2 đi qua điểm A(1; 14) và có trục đối xứng là đường thẳng x = 1. Giá trị của tổng a + b bằng. Phương án đúng là:\par
\hspace*{1.5em}\textbf{A.} 24\quad\quad \textbf{B.} -12\quad\quad \textbf{C.} 14\quad\quad \textbf{D.} 12\par\vspace{2pt}
\subsection*{\color{primary} Bảng đáp án Chương VI: Hàm số, đồ thị và ứng dụng}
\begin{center}
\renewcommand{\arraystretch}{1.25}
\begin{tabularx}{\textwidth}{*{10}{>{\centering\arraybackslash}X}}
\toprule
\textbf{\color{primary} Cột 1} & \textbf{\color{primary} Cột 2} & \textbf{\color{primary} Cột 3} & \textbf{\color{primary} Cột 4} & \textbf{\color{primary} Cột 5} & \textbf{\color{primary} Cột 6} & \textbf{\color{primary} Cột 7} & \textbf{\color{primary} Cột 8} & \textbf{\color{primary} Cột 9} & \textbf{\color{primary} Cột 10} \\
\midrule
\textbf{1.A} & \textbf{2.A} & \textbf{3.B} & \textbf{4.B} & \textbf{5.B} & \textbf{6.B} & \textbf{7.D} & \textbf{8.C} & \textbf{9.D} & \textbf{10.A} \\
\rowcolor{rowtint} \textbf{11.C} & \textbf{12.D} & \textbf{13.D} & \textbf{14.B} & \textbf{15.D} & \textbf{16.C} & \textbf{17.D} & \textbf{18.A} & \textbf{19.D} & \textbf{20.C} \\
\textbf{21.A} & \textbf{22.A} & \textbf{23.C} & \textbf{24.A} & \textbf{25.D} & \textbf{26.C} & \textbf{27.C} & \textbf{28.C} & \textbf{29.C} & \textbf{30.D} \\
\rowcolor{rowtint} \textbf{31.D} & \textbf{32.B} & \textbf{33.A} & \textbf{34.D} & \textbf{35.D} & \textbf{36.C} & \textbf{37.A} & \textbf{38.B} & \textbf{39.A} & \textbf{40.B} \\
\textbf{41.A} & \textbf{42.A} & \textbf{43.D} & \textbf{44.B} & \textbf{45.B} & \textbf{46.B} & \textbf{47.D} & \textbf{48.C} & \textbf{49.D} & \textbf{50.D} \\
\bottomrule
\end{tabularx}
\end{center}\vspace{10pt}
\section*{\color{primary} CHƯƠNG VII: PHƯƠNG PHÁP TỌA ĐỘ TRONG MẶT PHẲNG}
\textit{Trọng tâm kiến thức SGK: Phương trình đường thẳng (tham số, tổng quát), vị trí tương đối, góc và khoảng cách từ một điểm đến đường thẳng, phương trình đường tròn và ba đường conic (elip, hypebol, parabol).}\par\vspace{4pt}
\subsection*{\color{accent} Phần I. Mức độ Dễ (Nhận biết -- Câu 1 đến Câu 20)}
\textbf{\color{primary} Câu 1.} Cho đường thẳng $\Delta$ có phương trình tổng quát: 2x - 3y + 7 = 0. Một vectơ pháp tuyến của $\Delta$ là:\par
\hspace*{1.5em}\textbf{A.} $\vec{n}$ = (3; 2)\quad\quad \textbf{B.} $\vec{n}$ = (-3; 2)\quad\quad \textbf{C.} $\vec{n}$ = (2; -3)\quad\quad \textbf{D.} $\vec{n}$ = (2; 3)\par\vspace{2pt}
\textbf{\color{primary} Câu 2.} Thực hiện yêu cầu nhận biết sau: Tâm I và bán kính R của đường tròn (C): (x - 3)$^2$ + (y + 3)$^2$ = 25 lần lượt là.\par
\hspace*{1.5em}\textbf{A.} I(3; -3), R = 25\quad\quad \textbf{B.} I(-3; 3), R = 5\quad\quad \textbf{C.} I(3; 3), R = 5\quad\quad \textbf{D.} I(3; -3), R = 5\par\vspace{2pt}
\textbf{\color{primary} Câu 3.} Trong bài học thuộc Chương VII, xét nội dung: Một vectơ chỉ phương của đường thẳng có phương trình tham số x = 2 + 5t, y = -1 - 4t là. Đáp án đúng là:\par
\hspace*{1.5em}\textbf{A.} $\vec{u}$ = (5; -4)\quad\quad \textbf{B.} $\vec{u}$ = (5; 4)\quad\quad \textbf{C.} $\vec{u}$ = (4; 5)\quad\quad \textbf{D.} $\vec{u}$ = (2; -1)\par\vspace{2pt}
\textbf{\color{primary} Câu 4.} Dựa vào kiến thức chuẩn SGK của Chương VII (Phương pháp tọa độ trong mặt phẳng), cho đường thẳng $\Delta$ có phương trình tổng quát: 5x - 3y + 7 = 0. Một vectơ pháp tuyến của $\Delta$ là:\par
\hspace*{1.5em}\textbf{A.} $\vec{n}$ = (3; 5)\quad\quad \textbf{B.} $\vec{n}$ = (5; 3)\quad\quad \textbf{C.} $\vec{n}$ = (5; -3)\quad\quad \textbf{D.} $\vec{n}$ = (-3; 5)\par\vspace{2pt}
\textbf{\color{primary} Câu 5.} Chọn kết quả đúng nhất khi xét: Tâm I và bán kính R của đường tròn (C): (x - 6)$^2$ + (y + 3)$^2$ = 64 lần lượt là:\par
\hspace*{1.5em}\textbf{A.} I(-6; 3), R = 8\quad\quad \textbf{B.} I(6; -3), R = 64\quad\quad \textbf{C.} I(6; -3), R = 8\quad\quad \textbf{D.} I(6; 3), R = 8\par\vspace{2pt}
\textbf{\color{primary} Câu 6.} Khi làm bài tập cơ bản về phương pháp tọa độ trong mặt phẳng, một học sinh gặp câu hỏi: Một vectơ chỉ phương của đường thẳng có phương trình tham số x = 2 + 8t, y = -1 - 4t là. Kết quả nào dưới đây đúng?\par
\hspace*{1.5em}\textbf{A.} $\vec{u}$ = (8; 4)\quad\quad \textbf{B.} $\vec{u}$ = (8; -4)\quad\quad \textbf{C.} $\vec{u}$ = (4; 8)\quad\quad \textbf{D.} $\vec{u}$ = (2; -1)\par\vspace{2pt}
\textbf{\color{primary} Câu 7.} Xác định giá trị hoặc khẳng định đúng cho bài toán sau: Cho đường thẳng $\Delta$ có phương trình tổng quát: 8x - 3y + 7 = 0. Một vectơ pháp tuyến của $\Delta$ là:\par
\hspace*{1.5em}\textbf{A.} $\vec{n}$ = (8; -3)\quad\quad \textbf{B.} $\vec{n}$ = (8; 3)\quad\quad \textbf{C.} $\vec{n}$ = (-3; 8)\quad\quad \textbf{D.} $\vec{n}$ = (3; 8)\par\vspace{2pt}
\textbf{\color{primary} Câu 8.} Áp dụng trực tiếp định nghĩa và quy tắc trong Chương VII: Tâm I và bán kính R của đường tròn (C): (x - 9)$^2$ + (y + 3)$^2$ = 121 lần lượt là:\par
\hspace*{1.5em}\textbf{A.} I(9; -3), R = 121\quad\quad \textbf{B.} I(-9; 3), R = 11\quad\quad \textbf{C.} I(9; 3), R = 11\quad\quad \textbf{D.} I(9; -3), R = 11\par\vspace{2pt}
\textbf{\color{primary} Câu 9.} Kiểm tra nhanh kiến thức nền tảng: Một vectơ chỉ phương của đường thẳng có phương trình tham số x = 2 + 11t, y = -1 - 4t là:\par
\hspace*{1.5em}\textbf{A.} $\vec{u}$ = (2; -1)\quad\quad \textbf{B.} $\vec{u}$ = (4; 11)\quad\quad \textbf{C.} $\vec{u}$ = (11; 4)\quad\quad \textbf{D.} $\vec{u}$ = (11; -4)\par\vspace{2pt}
\textbf{\color{primary} Câu 10.} Cho dữ kiện nhận biết sau đây: Cho đường thẳng $\Delta$ có phương trình tổng quát: 11x - 3y + 7 = 0. Một vectơ pháp tuyến của $\Delta$ là. Phương án chính xác là:\par
\hspace*{1.5em}\textbf{A.} $\vec{n}$ = (-3; 11)\quad\quad \textbf{B.} $\vec{n}$ = (11; 3)\quad\quad \textbf{C.} $\vec{n}$ = (3; 11)\quad\quad \textbf{D.} $\vec{n}$ = (11; -3)\par\vspace{2pt}
\textbf{\color{primary} Câu 11.} Hãy chỉ ra kết quả đúng trong bốn lựa chọn A, B, C, D đối với yêu cầu: Tâm I và bán kính R của đường tròn (C): (x - 12)$^2$ + (y + 3)$^2$ = 196 lần lượt là:\par
\hspace*{1.5em}\textbf{A.} I(12; -3), R = 14\quad\quad \textbf{B.} I(12; -3), R = 196\quad\quad \textbf{C.} I(-12; 3), R = 14\quad\quad \textbf{D.} I(12; 3), R = 14\par\vspace{2pt}
\textbf{\color{primary} Câu 12.} Câu hỏi ôn tập lý thuyết và kỹ năng cơ bản (Chương VII): Một vectơ chỉ phương của đường thẳng có phương trình tham số x = 2 + 14t, y = -1 - 4t là:\par
\hspace*{1.5em}\textbf{A.} $\vec{u}$ = (14; 4)\quad\quad \textbf{B.} $\vec{u}$ = (14; -4)\quad\quad \textbf{C.} $\vec{u}$ = (4; 14)\quad\quad \textbf{D.} $\vec{u}$ = (2; -1)\par\vspace{2pt}
\textbf{\color{primary} Câu 13.} Theo nội dung sách giáo khoa Kết nối tri thức với cuộc sống (Chương VII), cho đường thẳng $\Delta$ có phương trình tổng quát: 14x - 3y + 7 = 0. Một vectơ pháp tuyến của $\Delta$ là:\par
\hspace*{1.5em}\textbf{A.} $\vec{n}$ = (14; 3)\quad\quad \textbf{B.} $\vec{n}$ = (-3; 14)\quad\quad \textbf{C.} $\vec{n}$ = (14; -3)\quad\quad \textbf{D.} $\vec{n}$ = (3; 14)\par\vspace{2pt}
\textbf{\color{primary} Câu 14.} Tìm đáp số đúng của bài toán cơ bản sau: Tâm I và bán kính R của đường tròn (C): (x - 15)$^2$ + (y + 3)$^2$ = 289 lần lượt là:\par
\hspace*{1.5em}\textbf{A.} I(15; -3), R = 289\quad\quad \textbf{B.} I(15; 3), R = 17\quad\quad \textbf{C.} I(-15; 3), R = 17\quad\quad \textbf{D.} I(15; -3), R = 17\par\vspace{2pt}
\textbf{\color{primary} Câu 15.} Xét đối tượng toán học trong Chương VII: Một vectơ chỉ phương của đường thẳng có phương trình tham số x = 2 + 17t, y = -1 - 4t là. Giá trị hoặc tính chất tương ứng là:\par
\hspace*{1.5em}\textbf{A.} $\vec{u}$ = (17; -4)\quad\quad \textbf{B.} $\vec{u}$ = (4; 17)\quad\quad \textbf{C.} $\vec{u}$ = (2; -1)\quad\quad \textbf{D.} $\vec{u}$ = (17; 4)\par\vspace{2pt}
\textbf{\color{primary} Câu 16.} Hoàn thành phép xác định cơ bản sau đây: Cho đường thẳng $\Delta$ có phương trình tổng quát: 17x - 3y + 7 = 0. Một vectơ pháp tuyến của $\Delta$ là:\par
\hspace*{1.5em}\textbf{A.} $\vec{n}$ = (17; 3)\quad\quad \textbf{B.} $\vec{n}$ = (3; 17)\quad\quad \textbf{C.} $\vec{n}$ = (17; -3)\quad\quad \textbf{D.} $\vec{n}$ = (-3; 17)\par\vspace{2pt}
\textbf{\color{primary} Câu 17.} Trong các khẳng định hoặc giá trị dưới đây, đâu là đáp án đúng cho câu hỏi: Tâm I và bán kính R của đường tròn (C): (x - 18)$^2$ + (y + 3)$^2$ = 400 lần lượt là?\par
\hspace*{1.5em}\textbf{A.} I(18; -3), R = 400\quad\quad \textbf{B.} I(18; 3), R = 20\quad\quad \textbf{C.} I(-18; 3), R = 20\quad\quad \textbf{D.} I(18; -3), R = 20\par\vspace{2pt}
\textbf{\color{primary} Câu 18.} Bài tập khởi động Chương VII: Một vectơ chỉ phương của đường thẳng có phương trình tham số x = 2 + 20t, y = -1 - 4t là:\par
\hspace*{1.5em}\textbf{A.} $\vec{u}$ = (2; -1)\quad\quad \textbf{B.} $\vec{u}$ = (4; 20)\quad\quad \textbf{C.} $\vec{u}$ = (20; 4)\quad\quad \textbf{D.} $\vec{u}$ = (20; -4)\par\vspace{2pt}
\textbf{\color{primary} Câu 19.} Quan sát dữ kiện đề bài: Cho đường thẳng $\Delta$ có phương trình tổng quát: 20x - 3y + 7 = 0. Một vectơ pháp tuyến của $\Delta$ là. Hãy chọn phương án đúng:\par
\hspace*{1.5em}\textbf{A.} $\vec{n}$ = (20; -3)\quad\quad \textbf{B.} $\vec{n}$ = (20; 3)\quad\quad \textbf{C.} $\vec{n}$ = (-3; 20)\quad\quad \textbf{D.} $\vec{n}$ = (3; 20)\par\vspace{2pt}
\textbf{\color{primary} Câu 20.} Sử dụng công thức hoặc tính chất đầu bài của Chương VII để giải quyết yêu cầu: Tâm I và bán kính R của đường tròn (C): (x - 21)$^2$ + (y + 3)$^2$ = 529 lần lượt là:\par
\hspace*{1.5em}\textbf{A.} I(21; -3), R = 23\quad\quad \textbf{B.} I(21; -3), R = 529\quad\quad \textbf{C.} I(-21; 3), R = 23\quad\quad \textbf{D.} I(21; 3), R = 23\par\vspace{2pt}
\subsection*{\color{accent} Phần II. Mức độ Trung bình (Thông hiểu -- Câu 21 đến Câu 40)}
\textbf{\color{primary} Câu 21.} Khoảng cách từ điểm M(1; 2) đến đường thẳng $\Delta$: 3x + 4y + -6 = 0 bằng:\par
\hspace*{1.5em}\textbf{A.} 1\quad\quad \textbf{B.} 0.2\quad\quad \textbf{C.} 2\quad\quad \textbf{D.} 5\par\vspace{2pt}
\textbf{\color{primary} Câu 22.} Vận dụng quy tắc trong Chương VII (Phương pháp tọa độ trong mặt phẳng) để giải bài toán: Số đo góc giữa hai đường thẳng d$_1$: x - 2y + 3 = 0 và d$_2$: 3x - y - 5 = 0 bằng:\par
\hspace*{1.5em}\textbf{A.} 45°\quad\quad \textbf{B.} 60°\quad\quad \textbf{C.} 30°\quad\quad \textbf{D.} 90°\par\vspace{2pt}
\textbf{\color{primary} Câu 23.} Biến đổi và tính toán theo kiến thức Chương VII: Cho elip (E): x$^2$/25 + y$^2$/9 = 1 (với a = 5, b = 3, c = 4). Độ dài tiêu cự 2c và độ dài trục lớn 2a của elip lần lượt là. Kết quả thu được là:\par
\begin{itemize}\setlength{\itemsep}{1pt}\setlength{\parskip}{0pt}
  \item[] \textbf{A.} Tiêu cự bằng 6, trục lớn bằng 10
  \item[] \textbf{B.} Tiêu cự bằng 8, trục lớn bằng 25
  \item[] \textbf{C.} Tiêu cự bằng 4, trục lớn bằng 5
  \item[] \textbf{D.} Tiêu cự bằng 8, trục lớn bằng 10
\end{itemize}\vspace{2pt}
\textbf{\color{primary} Câu 24.} Xét bài toán mức độ thông hiểu sau: Khoảng cách từ điểm M(1; 2) đến đường thẳng $\Delta$: 3x + 4y + 9 = 0 bằng. Đáp án đúng là:\par
\hspace*{1.5em}\textbf{A.} 4\quad\quad \textbf{B.} 20\quad\quad \textbf{C.} 5\quad\quad \textbf{D.} 0.8\par\vspace{2pt}
\textbf{\color{primary} Câu 25.} Để hoàn thành phiếu học tập Chương VII, hãy xác định: Số đo góc giữa hai đường thẳng d$_1$: x - 2y + 3 = 0 và d$_2$: 3x - y - 5 = 0 bằng:\par
\hspace*{1.5em}\textbf{A.} 45°\quad\quad \textbf{B.} 30°\quad\quad \textbf{C.} 90°\quad\quad \textbf{D.} 60°\par\vspace{2pt}
\textbf{\color{primary} Câu 26.} Phân tích dữ kiện bài toán: Cho elip (E): x$^2$/25 + y$^2$/9 = 1 (với a = 5, b = 3, c = 4). Độ dài tiêu cự 2c và độ dài trục lớn 2a của elip lần lượt là. Giá trị cần tìm bằng:\par
\begin{itemize}\setlength{\itemsep}{1pt}\setlength{\parskip}{0pt}
  \item[] \textbf{A.} Tiêu cự bằng 8, trục lớn bằng 10
  \item[] \textbf{B.} Tiêu cự bằng 4, trục lớn bằng 5
  \item[] \textbf{C.} Tiêu cự bằng 8, trục lớn bằng 25
  \item[] \textbf{D.} Tiêu cự bằng 6, trục lớn bằng 10
\end{itemize}\vspace{2pt}
\textbf{\color{primary} Câu 27.} Áp dụng tính chất đã học ở Chương VII vào bài tập sau: Khoảng cách từ điểm M(1; 2) đến đường thẳng $\Delta$: 3x + 4y + 24 = 0 bằng:\par
\hspace*{1.5em}\textbf{A.} 8\quad\quad \textbf{B.} 35\quad\quad \textbf{C.} 7\quad\quad \textbf{D.} 1.4\par\vspace{2pt}
\textbf{\color{primary} Câu 28.} Thực hiện các bước suy luận hợp lý để giải quyết yêu cầu: Số đo góc giữa hai đường thẳng d$_1$: x - 2y + 3 = 0 và d$_2$: 3x - y - 5 = 0 bằng:\par
\hspace*{1.5em}\textbf{A.} 60°\quad\quad \textbf{B.} 90°\quad\quad \textbf{C.} 45°\quad\quad \textbf{D.} 30°\par\vspace{2pt}
\textbf{\color{primary} Câu 29.} Cho bài toán tự luyện (Phương pháp tọa độ trong mặt phẳng): Cho elip (E): x$^2$/25 + y$^2$/9 = 1 (với a = 5, b = 3, c = 4). Độ dài tiêu cự 2c và độ dài trục lớn 2a của elip lần lượt là. Phương án nào sau đây chính xác?\par
\begin{itemize}\setlength{\itemsep}{1pt}\setlength{\parskip}{0pt}
  \item[] \textbf{A.} Tiêu cự bằng 8, trục lớn bằng 10
  \item[] \textbf{B.} Tiêu cự bằng 8, trục lớn bằng 25
  \item[] \textbf{C.} Tiêu cự bằng 6, trục lớn bằng 10
  \item[] \textbf{D.} Tiêu cự bằng 4, trục lớn bằng 5
\end{itemize}\vspace{2pt}
\textbf{\color{primary} Câu 30.} Kết hợp các phép biến đổi trong Chương VII, hãy tìm kết quả của bài toán: Khoảng cách từ điểm M(1; 2) đến đường thẳng $\Delta$: 3x + 4y + 39 = 0 bằng:\par
\hspace*{1.5em}\textbf{A.} 11\quad\quad \textbf{B.} 10\quad\quad \textbf{C.} 2.0\quad\quad \textbf{D.} 50\par\vspace{2pt}
\textbf{\color{primary} Câu 31.} Một bài kiểm tra định kỳ nội dung Chương VII nêu câu hỏi: Số đo góc giữa hai đường thẳng d$_1$: x - 2y + 3 = 0 và d$_2$: 3x - y - 5 = 0 bằng. Hãy chọn đáp số đúng:\par
\hspace*{1.5em}\textbf{A.} 30°\quad\quad \textbf{B.} 60°\quad\quad \textbf{C.} 90°\quad\quad \textbf{D.} 45°\par\vspace{2pt}
\textbf{\color{primary} Câu 32.} Tìm lời giải đúng cho bài toán thông hiểu: Cho elip (E): x$^2$/25 + y$^2$/9 = 1 (với a = 5, b = 3, c = 4). Độ dài tiêu cự 2c và độ dài trục lớn 2a của elip lần lượt là:\par
\begin{itemize}\setlength{\itemsep}{1pt}\setlength{\parskip}{0pt}
  \item[] \textbf{A.} Tiêu cự bằng 8, trục lớn bằng 25
  \item[] \textbf{B.} Tiêu cự bằng 8, trục lớn bằng 10
  \item[] \textbf{C.} Tiêu cự bằng 6, trục lớn bằng 10
  \item[] \textbf{D.} Tiêu cự bằng 4, trục lớn bằng 5
\end{itemize}\vspace{2pt}
\textbf{\color{primary} Câu 33.} Sử dụng mối liên hệ giữa các đại lượng trong Chương VII: Khoảng cách từ điểm M(1; 2) đến đường thẳng $\Delta$: 3x + 4y + 54 = 0 bằng:\par
\hspace*{1.5em}\textbf{A.} 2.6\quad\quad \textbf{B.} 13\quad\quad \textbf{C.} 14\quad\quad \textbf{D.} 65\par\vspace{2pt}
\textbf{\color{primary} Câu 34.} Hãy tính toán và chọn phương án đúng đối với bài tập: Số đo góc giữa hai đường thẳng d$_1$: x - 2y + 3 = 0 và d$_2$: 3x - y - 5 = 0 bằng:\par
\hspace*{1.5em}\textbf{A.} 45°\quad\quad \textbf{B.} 30°\quad\quad \textbf{C.} 60°\quad\quad \textbf{D.} 90°\par\vspace{2pt}
\textbf{\color{primary} Câu 35.} Dựa trên các tính chất trọng tâm của Chương VII, xét bài toán: Cho elip (E): x$^2$/25 + y$^2$/9 = 1 (với a = 5, b = 3, c = 4). Độ dài tiêu cự 2c và độ dài trục lớn 2a của elip lần lượt là:\par
\begin{itemize}\setlength{\itemsep}{1pt}\setlength{\parskip}{0pt}
  \item[] \textbf{A.} Tiêu cự bằng 4, trục lớn bằng 5
  \item[] \textbf{B.} Tiêu cự bằng 8, trục lớn bằng 10
  \item[] \textbf{C.} Tiêu cự bằng 6, trục lớn bằng 10
  \item[] \textbf{D.} Tiêu cự bằng 8, trục lớn bằng 25
\end{itemize}\vspace{2pt}
\textbf{\color{primary} Câu 36.} Thực hành giải toán (Phương pháp tọa độ trong mặt phẳng): Khoảng cách từ điểm M(1; 2) đến đường thẳng $\Delta$: 3x + 4y + 69 = 0 bằng. Đáp án nào dưới đây thỏa mãn?\par
\hspace*{1.5em}\textbf{A.} 17\quad\quad \textbf{B.} 16\quad\quad \textbf{C.} 80\quad\quad \textbf{D.} 3.2\par\vspace{2pt}
\textbf{\color{primary} Câu 37.} Kiểm tra kỹ năng biến đổi và suy luận qua bài toán: Số đo góc giữa hai đường thẳng d$_1$: x - 2y + 3 = 0 và d$_2$: 3x - y - 5 = 0 bằng:\par
\hspace*{1.5em}\textbf{A.} 45°\quad\quad \textbf{B.} 30°\quad\quad \textbf{C.} 60°\quad\quad \textbf{D.} 90°\par\vspace{2pt}
\textbf{\color{primary} Câu 38.} Cho giả thiết thuộc phạm vi Chương VII: Cho elip (E): x$^2$/25 + y$^2$/9 = 1 (với a = 5, b = 3, c = 4). Độ dài tiêu cự 2c và độ dài trục lớn 2a của elip lần lượt là. Kết luận đúng là:\par
\begin{itemize}\setlength{\itemsep}{1pt}\setlength{\parskip}{0pt}
  \item[] \textbf{A.} Tiêu cự bằng 4, trục lớn bằng 5
  \item[] \textbf{B.} Tiêu cự bằng 8, trục lớn bằng 10
  \item[] \textbf{C.} Tiêu cự bằng 8, trục lớn bằng 25
  \item[] \textbf{D.} Tiêu cự bằng 6, trục lớn bằng 10
\end{itemize}\vspace{2pt}
\textbf{\color{primary} Câu 39.} Rút gọn hoặc tính giá trị theo yêu cầu sau của Chương VII: Khoảng cách từ điểm M(1; 2) đến đường thẳng $\Delta$: 3x + 4y + 84 = 0 bằng:\par
\hspace*{1.5em}\textbf{A.} 3.8\quad\quad \textbf{B.} 95\quad\quad \textbf{C.} 20\quad\quad \textbf{D.} 19\par\vspace{2pt}
\textbf{\color{primary} Câu 40.} Xác định phương án đúng sau khi thực hiện phép tính / suy luận: Số đo góc giữa hai đường thẳng d$_1$: x - 2y + 3 = 0 và d$_2$: 3x - y - 5 = 0 bằng:\par
\hspace*{1.5em}\textbf{A.} 30°\quad\quad \textbf{B.} 90°\quad\quad \textbf{C.} 60°\quad\quad \textbf{D.} 45°\par\vspace{2pt}
\subsection*{\color{accent} Phần III. Mức độ Khó (Vận dụng \& Vận dụng cao -- Câu 41 đến Câu 50)}
\textbf{\color{primary} Câu 41.} Phương trình tiếp tuyến của đường tròn (C): (x - 1)$^2$ + (y - 2)$^2$ = 25 tại tiếp điểm M(4; 6) là:\par
\hspace*{1.5em}\textbf{A.} 4x + 3y - 34 = 0\quad\quad \textbf{B.} 3x + 4y - 25 = 0\quad\quad \textbf{C.} 3x - 4y + 12 = 0\quad\quad \textbf{D.} 3x + 4y - 36 = 0\par\vspace{2pt}
\textbf{\color{primary} Câu 42.} Bài toán vận dụng (Chương VII – Phương pháp tọa độ trong mặt phẳng): Cho parabol có phương trình chính tắc (P): y$^2$ = 12x. Tọa độ tiêu điểm F và phương trình đường chuẩn $\Delta$ của (P) là. Kết quả đúng là:\par
\hspace*{1.5em}\textbf{A.} F(6; 0) và $\Delta$: x + 6 = 0\quad\quad \textbf{B.} F(3; 0) và $\Delta$: x + 3 = 0\quad\quad \textbf{C.} F(0; 3) và $\Delta$: y + 3 = 0\quad\quad \textbf{D.} F(3; 0) và $\Delta$: x - 3 = 0\par\vspace{2pt}
\textbf{\color{primary} Câu 43.} Vận dụng kiến thức nâng cao trong phạm vi Chương VII để giải bài toán: Phương trình tiếp tuyến của đường tròn (C): (x - 1)$^2$ + (y - 2)$^2$ = 25 tại tiếp điểm M(4; 6) là:\par
\hspace*{1.5em}\textbf{A.} 3x + 4y - 36 = 0\quad\quad \textbf{B.} 3x + 4y - 25 = 0\quad\quad \textbf{C.} 3x - 4y + 12 = 0\quad\quad \textbf{D.} 4x + 3y - 34 = 0\par\vspace{2pt}
\textbf{\color{primary} Câu 44.} Phân tích và giải bài toán tổng hợp sau của Chương VII: Cho parabol có phương trình chính tắc (P): y$^2$ = 20x. Tọa độ tiêu điểm F và phương trình đường chuẩn $\Delta$ của (P) là:\par
\begin{itemize}\setlength{\itemsep}{1pt}\setlength{\parskip}{0pt}
  \item[] \textbf{A.} F(0; 5) và $\Delta$: y + 5 = 0
  \item[] \textbf{B.} F(5; 0) và $\Delta$: x + 5 = 0
  \item[] \textbf{C.} F(10; 0) và $\Delta$: x + 10 = 0
  \item[] \textbf{D.} F(5; 0) và $\Delta$: x - 5 = 0
\end{itemize}\vspace{2pt}
\textbf{\color{primary} Câu 45.} Thử thách tư duy trong phạm vi kiến thức Chương VII: Phương trình tiếp tuyến của đường tròn (C): (x - 1)$^2$ + (y - 2)$^2$ = 25 tại tiếp điểm M(4; 6) là. Đáp số chính xác là:\par
\hspace*{1.5em}\textbf{A.} 3x - 4y + 12 = 0\quad\quad \textbf{B.} 3x + 4y - 25 = 0\quad\quad \textbf{C.} 4x + 3y - 34 = 0\quad\quad \textbf{D.} 3x + 4y - 36 = 0\par\vspace{2pt}
\textbf{\color{primary} Câu 46.} Sử dụng phương pháp giải toán của Chương VII để tìm đáp án cho bài toán: Cho parabol có phương trình chính tắc (P): y$^2$ = 28x. Tọa độ tiêu điểm F và phương trình đường chuẩn $\Delta$ của (P) là:\par
\begin{itemize}\setlength{\itemsep}{1pt}\setlength{\parskip}{0pt}
  \item[] \textbf{A.} F(7; 0) và $\Delta$: x + 7 = 0
  \item[] \textbf{B.} F(14; 0) và $\Delta$: x + 14 = 0
  \item[] \textbf{C.} F(0; 7) và $\Delta$: y + 7 = 0
  \item[] \textbf{D.} F(7; 0) và $\Delta$: x - 7 = 0
\end{itemize}\vspace{2pt}
\textbf{\color{primary} Câu 47.} Xét bài toán vận dụng thực tiễn / nâng cao thuộc Chương VII: Phương trình tiếp tuyến của đường tròn (C): (x - 1)$^2$ + (y - 2)$^2$ = 25 tại tiếp điểm M(4; 6) là:\par
\hspace*{1.5em}\textbf{A.} 3x - 4y + 12 = 0\quad\quad \textbf{B.} 3x + 4y - 36 = 0\quad\quad \textbf{C.} 4x + 3y - 34 = 0\quad\quad \textbf{D.} 3x + 4y - 25 = 0\par\vspace{2pt}
\textbf{\color{primary} Câu 48.} Tìm kết quả cuối cùng của bài toán phân loại học sinh khá giỏi (Chương VII): Cho parabol có phương trình chính tắc (P): y$^2$ = 36x. Tọa độ tiêu điểm F và phương trình đường chuẩn $\Delta$ của (P) là:\par
\begin{itemize}\setlength{\itemsep}{1pt}\setlength{\parskip}{0pt}
  \item[] \textbf{A.} F(18; 0) và $\Delta$: x + 18 = 0
  \item[] \textbf{B.} F(0; 9) và $\Delta$: y + 9 = 0
  \item[] \textbf{C.} F(9; 0) và $\Delta$: x - 9 = 0
  \item[] \textbf{D.} F(9; 0) và $\Delta$: x + 9 = 0
\end{itemize}\vspace{2pt}
\textbf{\color{primary} Câu 49.} Kết hợp các tính chất đã học trong Chương VII, hãy giải quyết bài toán: Phương trình tiếp tuyến của đường tròn (C): (x - 1)$^2$ + (y - 2)$^2$ = 25 tại tiếp điểm M(4; 6) là:\par
\hspace*{1.5em}\textbf{A.} 3x + 4y - 36 = 0\quad\quad \textbf{B.} 3x + 4y - 25 = 0\quad\quad \textbf{C.} 3x - 4y + 12 = 0\quad\quad \textbf{D.} 4x + 3y - 34 = 0\par\vspace{2pt}
\textbf{\color{primary} Câu 50.} Câu hỏi vận dụng cao cuối Chương VII: Cho parabol có phương trình chính tắc (P): y$^2$ = 44x. Tọa độ tiêu điểm F và phương trình đường chuẩn $\Delta$ của (P) là. Phương án đúng là:\par
\begin{itemize}\setlength{\itemsep}{1pt}\setlength{\parskip}{0pt}
  \item[] \textbf{A.} F(11; 0) và $\Delta$: x + 11 = 0
  \item[] \textbf{B.} F(11; 0) và $\Delta$: x - 11 = 0
  \item[] \textbf{C.} F(22; 0) và $\Delta$: x + 22 = 0
  \item[] \textbf{D.} F(0; 11) và $\Delta$: y + 11 = 0
\end{itemize}\vspace{2pt}
\subsection*{\color{primary} Bảng đáp án Chương VII: Phương pháp tọa độ trong mặt phẳng}
\begin{center}
\renewcommand{\arraystretch}{1.25}
\begin{tabularx}{\textwidth}{*{10}{>{\centering\arraybackslash}X}}
\toprule
\textbf{\color{primary} Cột 1} & \textbf{\color{primary} Cột 2} & \textbf{\color{primary} Cột 3} & \textbf{\color{primary} Cột 4} & \textbf{\color{primary} Cột 5} & \textbf{\color{primary} Cột 6} & \textbf{\color{primary} Cột 7} & \textbf{\color{primary} Cột 8} & \textbf{\color{primary} Cột 9} & \textbf{\color{primary} Cột 10} \\
\midrule
\textbf{1.C} & \textbf{2.D} & \textbf{3.A} & \textbf{4.C} & \textbf{5.C} & \textbf{6.B} & \textbf{7.A} & \textbf{8.D} & \textbf{9.D} & \textbf{10.D} \\
\rowcolor{rowtint} \textbf{11.A} & \textbf{12.B} & \textbf{13.C} & \textbf{14.D} & \textbf{15.A} & \textbf{16.C} & \textbf{17.D} & \textbf{18.D} & \textbf{19.A} & \textbf{20.A} \\
\textbf{21.A} & \textbf{22.A} & \textbf{23.D} & \textbf{24.A} & \textbf{25.A} & \textbf{26.A} & \textbf{27.C} & \textbf{28.C} & \textbf{29.A} & \textbf{30.B} \\
\rowcolor{rowtint} \textbf{31.D} & \textbf{32.B} & \textbf{33.B} & \textbf{34.A} & \textbf{35.B} & \textbf{36.B} & \textbf{37.A} & \textbf{38.B} & \textbf{39.D} & \textbf{40.D} \\
\textbf{41.D} & \textbf{42.B} & \textbf{43.A} & \textbf{44.B} & \textbf{45.D} & \textbf{46.A} & \textbf{47.B} & \textbf{48.D} & \textbf{49.A} & \textbf{50.A} \\
\bottomrule
\end{tabularx}
\end{center}\vspace{10pt}
\section*{\color{primary} CHƯƠNG VIII: ĐẠI SỐ TỔ HỢP}
\textit{Trọng tâm kiến thức SGK: Quy tắc cộng và quy tắc nhân, hoán vị, chỉnh hợp, tổ hợp và nhị thức Newton (khai triển với số mũ n = 4, n = 5).}\par\vspace{4pt}
\subsection*{\color{accent} Phần I. Mức độ Dễ (Nhận biết -- Câu 1 đến Câu 20)}
\textbf{\color{primary} Câu 1.} Một lớp có 18 học sinh nam và 20 học sinh nữ. Số cách chọn ra 1 học sinh bất kỳ làm lớp trưởng là:\par
\hspace*{1.5em}\textbf{A.} 38 cách\quad\quad \textbf{B.} 20 cách\quad\quad \textbf{C.} 18 cách\quad\quad \textbf{D.} 360 cách\par\vspace{2pt}
\textbf{\color{primary} Câu 2.} Thực hiện yêu cầu nhận biết sau: Số các hoán vị của một tập hợp gồm 4 phần tử phân biệt là.\par
\hspace*{1.5em}\textbf{A.} 24\quad\quad \textbf{B.} 16\quad\quad \textbf{C.} 6\quad\quad \textbf{D.} 24 (1)\par\vspace{2pt}
\textbf{\color{primary} Câu 3.} Trong bài học thuộc Chương VIII, xét nội dung: Số cách chọn ra 2 học sinh từ một tổ gồm 8 học sinh (không phân biệt thứ tự) là. Đáp án đúng là:\par
\hspace*{1.5em}\textbf{A.} 16\quad\quad \textbf{B.} 64\quad\quad \textbf{C.} 28\quad\quad \textbf{D.} 56\par\vspace{2pt}
\textbf{\color{primary} Câu 4.} Dựa vào kiến thức chuẩn SGK của Chương VIII (Đại số tổ hợp), một lớp có 21 học sinh nam và 23 học sinh nữ. Số cách chọn ra 1 học sinh bất kỳ làm lớp trưởng là:\par
\hspace*{1.5em}\textbf{A.} 483 cách\quad\quad \textbf{B.} 23 cách\quad\quad \textbf{C.} 21 cách\quad\quad \textbf{D.} 44 cách\par\vspace{2pt}
\textbf{\color{primary} Câu 5.} Chọn kết quả đúng nhất khi xét: Số các hoán vị của một tập hợp gồm 7 phần tử phân biệt là:\par
\hspace*{1.5em}\textbf{A.} 128\quad\quad \textbf{B.} 720\quad\quad \textbf{C.} 49\quad\quad \textbf{D.} 5040\par\vspace{2pt}
\textbf{\color{primary} Câu 6.} Khi làm bài tập cơ bản về đại số tổ hợp, một học sinh gặp câu hỏi: Số cách chọn ra 2 học sinh từ một tổ gồm 11 học sinh (không phân biệt thứ tự) là. Kết quả nào dưới đây đúng?\par
\hspace*{1.5em}\textbf{A.} 121\quad\quad \textbf{B.} 55\quad\quad \textbf{C.} 22\quad\quad \textbf{D.} 110\par\vspace{2pt}
\textbf{\color{primary} Câu 7.} Xác định giá trị hoặc khẳng định đúng cho bài toán sau: Một lớp có 24 học sinh nam và 26 học sinh nữ. Số cách chọn ra 1 học sinh bất kỳ làm lớp trưởng là:\par
\hspace*{1.5em}\textbf{A.} 624 cách\quad\quad \textbf{B.} 26 cách\quad\quad \textbf{C.} 24 cách\quad\quad \textbf{D.} 50 cách\par\vspace{2pt}
\textbf{\color{primary} Câu 8.} Áp dụng trực tiếp định nghĩa và quy tắc trong Chương VIII: Số các hoán vị của một tập hợp gồm 5 phần tử phân biệt là:\par
\hspace*{1.5em}\textbf{A.} 24\quad\quad \textbf{B.} 32\quad\quad \textbf{C.} 120\quad\quad \textbf{D.} 25\par\vspace{2pt}
\textbf{\color{primary} Câu 9.} Kiểm tra nhanh kiến thức nền tảng: Số cách chọn ra 2 học sinh từ một tổ gồm 14 học sinh (không phân biệt thứ tự) là:\par
\hspace*{1.5em}\textbf{A.} 28\quad\quad \textbf{B.} 182\quad\quad \textbf{C.} 91\quad\quad \textbf{D.} 196\par\vspace{2pt}
\textbf{\color{primary} Câu 10.} Cho dữ kiện nhận biết sau đây: Một lớp có 27 học sinh nam và 29 học sinh nữ. Số cách chọn ra 1 học sinh bất kỳ làm lớp trưởng là. Phương án chính xác là:\par
\hspace*{1.5em}\textbf{A.} 29 cách\quad\quad \textbf{B.} 27 cách\quad\quad \textbf{C.} 56 cách\quad\quad \textbf{D.} 783 cách\par\vspace{2pt}
\textbf{\color{primary} Câu 11.} Hãy chỉ ra kết quả đúng trong bốn lựa chọn A, B, C, D đối với yêu cầu: Số các hoán vị của một tập hợp gồm 3 phần tử phân biệt là:\par
\hspace*{1.5em}\textbf{A.} 2\quad\quad \textbf{B.} 9\quad\quad \textbf{C.} 6\quad\quad \textbf{D.} 8\par\vspace{2pt}
\textbf{\color{primary} Câu 12.} Câu hỏi ôn tập lý thuyết và kỹ năng cơ bản (Chương VIII): Số cách chọn ra 2 học sinh từ một tổ gồm 17 học sinh (không phân biệt thứ tự) là:\par
\hspace*{1.5em}\textbf{A.} 289\quad\quad \textbf{B.} 272\quad\quad \textbf{C.} 136\quad\quad \textbf{D.} 34\par\vspace{2pt}
\textbf{\color{primary} Câu 13.} Theo nội dung sách giáo khoa Kết nối tri thức với cuộc sống (Chương VIII), một lớp có 30 học sinh nam và 32 học sinh nữ. Số cách chọn ra 1 học sinh bất kỳ làm lớp trưởng là:\par
\hspace*{1.5em}\textbf{A.} 32 cách\quad\quad \textbf{B.} 30 cách\quad\quad \textbf{C.} 960 cách\quad\quad \textbf{D.} 62 cách\par\vspace{2pt}
\textbf{\color{primary} Câu 14.} Tìm đáp số đúng của bài toán cơ bản sau: Số các hoán vị của một tập hợp gồm 6 phần tử phân biệt là:\par
\hspace*{1.5em}\textbf{A.} 36\quad\quad \textbf{B.} 720\quad\quad \textbf{C.} 120\quad\quad \textbf{D.} 64\par\vspace{2pt}
\textbf{\color{primary} Câu 15.} Xét đối tượng toán học trong Chương VIII: Số cách chọn ra 2 học sinh từ một tổ gồm 20 học sinh (không phân biệt thứ tự) là. Giá trị hoặc tính chất tương ứng là:\par
\hspace*{1.5em}\textbf{A.} 400\quad\quad \textbf{B.} 380\quad\quad \textbf{C.} 190\quad\quad \textbf{D.} 40\par\vspace{2pt}
\textbf{\color{primary} Câu 16.} Hoàn thành phép xác định cơ bản sau đây: Một lớp có 33 học sinh nam và 35 học sinh nữ. Số cách chọn ra 1 học sinh bất kỳ làm lớp trưởng là:\par
\hspace*{1.5em}\textbf{A.} 33 cách\quad\quad \textbf{B.} 68 cách\quad\quad \textbf{C.} 1155 cách\quad\quad \textbf{D.} 35 cách\par\vspace{2pt}
\textbf{\color{primary} Câu 17.} Trong các khẳng định hoặc giá trị dưới đây, đâu là đáp án đúng cho câu hỏi: Số các hoán vị của một tập hợp gồm 4 phần tử phân biệt là?\par
\hspace*{1.5em}\textbf{A.} 6\quad\quad \textbf{B.} 16\quad\quad \textbf{C.} 24\quad\quad \textbf{D.} 24 (1)\par\vspace{2pt}
\textbf{\color{primary} Câu 18.} Bài tập khởi động Chương VIII: Số cách chọn ra 2 học sinh từ một tổ gồm 23 học sinh (không phân biệt thứ tự) là:\par
\hspace*{1.5em}\textbf{A.} 253\quad\quad \textbf{B.} 46\quad\quad \textbf{C.} 506\quad\quad \textbf{D.} 529\par\vspace{2pt}
\textbf{\color{primary} Câu 19.} Quan sát dữ kiện đề bài: Một lớp có 36 học sinh nam và 38 học sinh nữ. Số cách chọn ra 1 học sinh bất kỳ làm lớp trưởng là. Hãy chọn phương án đúng:\par
\hspace*{1.5em}\textbf{A.} 38 cách\quad\quad \textbf{B.} 1368 cách\quad\quad \textbf{C.} 74 cách\quad\quad \textbf{D.} 36 cách\par\vspace{2pt}
\textbf{\color{primary} Câu 20.} Sử dụng công thức hoặc tính chất đầu bài của Chương VIII để giải quyết yêu cầu: Số các hoán vị của một tập hợp gồm 7 phần tử phân biệt là:\par
\hspace*{1.5em}\textbf{A.} 5040\quad\quad \textbf{B.} 128\quad\quad \textbf{C.} 49\quad\quad \textbf{D.} 720\par\vspace{2pt}
\subsection*{\color{accent} Phần II. Mức độ Trung bình (Thông hiểu -- Câu 21 đến Câu 40)}
\textbf{\color{primary} Câu 21.} Từ các chữ số 1, 2, 3, 4, 5, 6 có thể lập được bao nhiêu số tự nhiên có 3 chữ số đôi một khác nhau?\par
\hspace*{1.5em}\textbf{A.} 216 số\quad\quad \textbf{B.} 20 số\quad\quad \textbf{C.} 120 số\quad\quad \textbf{D.} 18 số\par\vspace{2pt}
\textbf{\color{primary} Câu 22.} Vận dụng quy tắc trong Chương VIII (Đại số tổ hợp) để giải bài toán: Trong khai triển nhị thức Newton của (x + 3)$^4$, hệ số của số hạng chứa x$^2$ là:\par
\hspace*{1.5em}\textbf{A.} 36\quad\quad \textbf{B.} 9\quad\quad \textbf{C.} 54\quad\quad \textbf{D.} 18\par\vspace{2pt}
\textbf{\color{primary} Câu 23.} Biến đổi và tính toán theo kiến thức Chương VIII: Trong khai triển nhị thức Newton của (2x - 1)$^5$, tổng tất cả các hệ số của đa thức nhận được (thay x = 1) bằng. Kết quả thu được là:\par
\hspace*{1.5em}\textbf{A.} -1\quad\quad \textbf{B.} 1\quad\quad \textbf{C.} 243\quad\quad \textbf{D.} 32\par\vspace{2pt}
\textbf{\color{primary} Câu 24.} Xét bài toán mức độ thông hiểu sau: Từ các chữ số 1, 2, 3, 4, 5, 6 có thể lập được bao nhiêu số tự nhiên có 3 chữ số đôi một khác nhau. Đáp án đúng là:\par
\hspace*{1.5em}\textbf{A.} 18 số\quad\quad \textbf{B.} 216 số\quad\quad \textbf{C.} 120 số\quad\quad \textbf{D.} 20 số\par\vspace{2pt}
\textbf{\color{primary} Câu 25.} Để hoàn thành phiếu học tập Chương VIII, hãy xác định: Trong khai triển nhị thức Newton của (x + 6)$^4$, hệ số của số hạng chứa x$^2$ là:\par
\hspace*{1.5em}\textbf{A.} 216\quad\quad \textbf{B.} 36\quad\quad \textbf{C.} 144\quad\quad \textbf{D.} 216 (1)\par\vspace{2pt}
\textbf{\color{primary} Câu 26.} Phân tích dữ kiện bài toán: Trong khai triển nhị thức Newton của (2x - 1)$^5$, tổng tất cả các hệ số của đa thức nhận được (thay x = 1) bằng. Giá trị cần tìm bằng:\par
\hspace*{1.5em}\textbf{A.} 243\quad\quad \textbf{B.} 32\quad\quad \textbf{C.} -1\quad\quad \textbf{D.} 1\par\vspace{2pt}
\textbf{\color{primary} Câu 27.} Áp dụng tính chất đã học ở Chương VIII vào bài tập sau: Từ các chữ số 1, 2, 3, 4, 5, 6 có thể lập được bao nhiêu số tự nhiên có 3 chữ số đôi một khác nhau:\par
\hspace*{1.5em}\textbf{A.} 18 số\quad\quad \textbf{B.} 216 số\quad\quad \textbf{C.} 20 số\quad\quad \textbf{D.} 120 số\par\vspace{2pt}
\textbf{\color{primary} Câu 28.} Thực hiện các bước suy luận hợp lý để giải quyết yêu cầu: Trong khai triển nhị thức Newton của (x + 9)$^4$, hệ số của số hạng chứa x$^2$ là:\par
\hspace*{1.5em}\textbf{A.} 486\quad\quad \textbf{B.} 54\quad\quad \textbf{C.} 81\quad\quad \textbf{D.} 324\par\vspace{2pt}
\textbf{\color{primary} Câu 29.} Cho bài toán tự luyện (Đại số tổ hợp): Trong khai triển nhị thức Newton của (2x - 1)$^5$, tổng tất cả các hệ số của đa thức nhận được (thay x = 1) bằng. Phương án nào sau đây chính xác?\par
\hspace*{1.5em}\textbf{A.} -1\quad\quad \textbf{B.} 1\quad\quad \textbf{C.} 32\quad\quad \textbf{D.} 243\par\vspace{2pt}
\textbf{\color{primary} Câu 30.} Kết hợp các phép biến đổi trong Chương VIII, hãy tìm kết quả của bài toán: Từ các chữ số 1, 2, 3, 4, 5, 6 có thể lập được bao nhiêu số tự nhiên có 3 chữ số đôi một khác nhau:\par
\hspace*{1.5em}\textbf{A.} 216 số\quad\quad \textbf{B.} 18 số\quad\quad \textbf{C.} 120 số\quad\quad \textbf{D.} 20 số\par\vspace{2pt}
\textbf{\color{primary} Câu 31.} Một bài kiểm tra định kỳ nội dung Chương VIII nêu câu hỏi: Trong khai triển nhị thức Newton của (x + 12)$^4$, hệ số của số hạng chứa x$^2$ là. Hãy chọn đáp số đúng:\par
\hspace*{1.5em}\textbf{A.} 576\quad\quad \textbf{B.} 72\quad\quad \textbf{C.} 144\quad\quad \textbf{D.} 864\par\vspace{2pt}
\textbf{\color{primary} Câu 32.} Tìm lời giải đúng cho bài toán thông hiểu: Trong khai triển nhị thức Newton của (2x - 1)$^5$, tổng tất cả các hệ số của đa thức nhận được (thay x = 1) bằng:\par
\hspace*{1.5em}\textbf{A.} 32\quad\quad \textbf{B.} 243\quad\quad \textbf{C.} 1\quad\quad \textbf{D.} -1\par\vspace{2pt}
\textbf{\color{primary} Câu 33.} Sử dụng mối liên hệ giữa các đại lượng trong Chương VIII: Từ các chữ số 1, 2, 3, 4, 5, 6 có thể lập được bao nhiêu số tự nhiên có 3 chữ số đôi một khác nhau:\par
\hspace*{1.5em}\textbf{A.} 18 số\quad\quad \textbf{B.} 120 số\quad\quad \textbf{C.} 20 số\quad\quad \textbf{D.} 216 số\par\vspace{2pt}
\textbf{\color{primary} Câu 34.} Hãy tính toán và chọn phương án đúng đối với bài tập: Trong khai triển nhị thức Newton của (x + 15)$^4$, hệ số của số hạng chứa x$^2$ là:\par
\hspace*{1.5em}\textbf{A.} 1350\quad\quad \textbf{B.} 225\quad\quad \textbf{C.} 900\quad\quad \textbf{D.} 90\par\vspace{2pt}
\textbf{\color{primary} Câu 35.} Dựa trên các tính chất trọng tâm của Chương VIII, xét bài toán: Trong khai triển nhị thức Newton của (2x - 1)$^5$, tổng tất cả các hệ số của đa thức nhận được (thay x = 1) bằng:\par
\hspace*{1.5em}\textbf{A.} 1\quad\quad \textbf{B.} 32\quad\quad \textbf{C.} 243\quad\quad \textbf{D.} -1\par\vspace{2pt}
\textbf{\color{primary} Câu 36.} Thực hành giải toán (Đại số tổ hợp): Từ các chữ số 1, 2, 3, 4, 5, 6 có thể lập được bao nhiêu số tự nhiên có 3 chữ số đôi một khác nhau. Đáp án nào dưới đây thỏa mãn?\par
\hspace*{1.5em}\textbf{A.} 20 số\quad\quad \textbf{B.} 18 số\quad\quad \textbf{C.} 120 số\quad\quad \textbf{D.} 216 số\par\vspace{2pt}
\textbf{\color{primary} Câu 37.} Kiểm tra kỹ năng biến đổi và suy luận qua bài toán: Trong khai triển nhị thức Newton của (x + 18)$^4$, hệ số của số hạng chứa x$^2$ là:\par
\hspace*{1.5em}\textbf{A.} 1944\quad\quad \textbf{B.} 108\quad\quad \textbf{C.} 324\quad\quad \textbf{D.} 1296\par\vspace{2pt}
\textbf{\color{primary} Câu 38.} Cho giả thiết thuộc phạm vi Chương VIII: Trong khai triển nhị thức Newton của (2x - 1)$^5$, tổng tất cả các hệ số của đa thức nhận được (thay x = 1) bằng. Kết luận đúng là:\par
\hspace*{1.5em}\textbf{A.} -1\quad\quad \textbf{B.} 1\quad\quad \textbf{C.} 243\quad\quad \textbf{D.} 32\par\vspace{2pt}
\textbf{\color{primary} Câu 39.} Rút gọn hoặc tính giá trị theo yêu cầu sau của Chương VIII: Từ các chữ số 1, 2, 3, 4, 5, 6 có thể lập được bao nhiêu số tự nhiên có 3 chữ số đôi một khác nhau:\par
\hspace*{1.5em}\textbf{A.} 216 số\quad\quad \textbf{B.} 120 số\quad\quad \textbf{C.} 18 số\quad\quad \textbf{D.} 20 số\par\vspace{2pt}
\textbf{\color{primary} Câu 40.} Xác định phương án đúng sau khi thực hiện phép tính / suy luận: Trong khai triển nhị thức Newton của (x + 21)$^4$, hệ số của số hạng chứa x$^2$ là:\par
\hspace*{1.5em}\textbf{A.} 441\quad\quad \textbf{B.} 1764\quad\quad \textbf{C.} 126\quad\quad \textbf{D.} 2646\par\vspace{2pt}
\subsection*{\color{accent} Phần III. Mức độ Khó (Vận dụng \& Vận dụng cao -- Câu 41 đến Câu 50)}
\textbf{\color{primary} Câu 41.} Một tổ có 6 học sinh nam và 4 học sinh nữ. Có bao nhiêu cách chọn ra một nhóm 3 học sinh sao cho trong nhóm có ít nhất 1 học sinh nữ?\par
\hspace*{1.5em}\textbf{A.} 120 cách\quad\quad \textbf{B.} 100 cách\quad\quad \textbf{C.} 96 cách\quad\quad \textbf{D.} 80 cách\par\vspace{2pt}
\textbf{\color{primary} Câu 42.} Bài toán vận dụng (Chương VIII – Đại số tổ hợp): Tìm số tự nhiên n (n $\ge$ 2) thỏa mãn đẳng thức $A_n^2$ + $C_n^2$ = 30 (trong đó $A_n^2$ là chỉnh hợp và $C_n^2$ là tổ hợp chập 2 của n). Kết quả đúng là:\par
\hspace*{1.5em}\textbf{A.} n = 5\quad\quad \textbf{B.} n = 6\quad\quad \textbf{C.} n = 7\quad\quad \textbf{D.} n = 4\par\vspace{2pt}
\textbf{\color{primary} Câu 43.} Vận dụng kiến thức nâng cao trong phạm vi Chương VIII để giải bài toán: Một tổ gồm 7 học sinh nam và 3 học sinh nữ. Số cách chọn ra một nhóm gồm 3 học sinh trong đó có đúng 2 học sinh nam và 1 học sinh nữ là:\par
\hspace*{1.5em}\textbf{A.} 10 cách\quad\quad \textbf{B.} 63 cách\quad\quad \textbf{C.} 21 cách\quad\quad \textbf{D.} 126 cách\par\vspace{2pt}
\textbf{\color{primary} Câu 44.} Phân tích và giải bài toán tổng hợp sau của Chương VIII: Một tổ có 6 học sinh nam và 4 học sinh nữ. Có bao nhiêu cách chọn ra một nhóm 3 học sinh sao cho trong nhóm có ít nhất 1 học sinh nữ:\par
\hspace*{1.5em}\textbf{A.} 100 cách\quad\quad \textbf{B.} 96 cách\quad\quad \textbf{C.} 120 cách\quad\quad \textbf{D.} 80 cách\par\vspace{2pt}
\textbf{\color{primary} Câu 45.} Thử thách tư duy trong phạm vi kiến thức Chương VIII: Tìm số tự nhiên n (n $\ge$ 2) thỏa mãn đẳng thức $A_n^2$ + $C_n^2$ = 30 (trong đó $A_n^2$ là chỉnh hợp và $C_n^2$ là tổ hợp chập 2 của n). Đáp số chính xác là:\par
\hspace*{1.5em}\textbf{A.} n = 4\quad\quad \textbf{B.} n = 5\quad\quad \textbf{C.} n = 7\quad\quad \textbf{D.} n = 6\par\vspace{2pt}
\textbf{\color{primary} Câu 46.} Sử dụng phương pháp giải toán của Chương VIII để tìm đáp án cho bài toán: Một tổ gồm 10 học sinh nam và 3 học sinh nữ. Số cách chọn ra một nhóm gồm 3 học sinh trong đó có đúng 2 học sinh nam và 1 học sinh nữ là:\par
\hspace*{1.5em}\textbf{A.} 270 cách\quad\quad \textbf{B.} 13 cách\quad\quad \textbf{C.} 45 cách\quad\quad \textbf{D.} 135 cách\par\vspace{2pt}
\textbf{\color{primary} Câu 47.} Xét bài toán vận dụng thực tiễn / nâng cao thuộc Chương VIII: Một tổ có 6 học sinh nam và 4 học sinh nữ. Có bao nhiêu cách chọn ra một nhóm 3 học sinh sao cho trong nhóm có ít nhất 1 học sinh nữ:\par
\hspace*{1.5em}\textbf{A.} 120 cách\quad\quad \textbf{B.} 100 cách\quad\quad \textbf{C.} 80 cách\quad\quad \textbf{D.} 96 cách\par\vspace{2pt}
\textbf{\color{primary} Câu 48.} Tìm kết quả cuối cùng của bài toán phân loại học sinh khá giỏi (Chương VIII): Tìm số tự nhiên n (n $\ge$ 2) thỏa mãn đẳng thức $A_n^2$ + $C_n^2$ = 30 (trong đó $A_n^2$ là chỉnh hợp và $C_n^2$ là tổ hợp chập 2 của n):\par
\hspace*{1.5em}\textbf{A.} n = 7\quad\quad \textbf{B.} n = 5\quad\quad \textbf{C.} n = 6\quad\quad \textbf{D.} n = 4\par\vspace{2pt}
\textbf{\color{primary} Câu 49.} Kết hợp các tính chất đã học trong Chương VIII, hãy giải quyết bài toán: Một tổ gồm 13 học sinh nam và 3 học sinh nữ. Số cách chọn ra một nhóm gồm 3 học sinh trong đó có đúng 2 học sinh nam và 1 học sinh nữ là:\par
\hspace*{1.5em}\textbf{A.} 16 cách\quad\quad \textbf{B.} 468 cách\quad\quad \textbf{C.} 234 cách\quad\quad \textbf{D.} 78 cách\par\vspace{2pt}
\textbf{\color{primary} Câu 50.} Câu hỏi vận dụng cao cuối Chương VIII: Một tổ có 6 học sinh nam và 4 học sinh nữ. Có bao nhiêu cách chọn ra một nhóm 3 học sinh sao cho trong nhóm có ít nhất 1 học sinh nữ. Phương án đúng là:\par
\hspace*{1.5em}\textbf{A.} 100 cách\quad\quad \textbf{B.} 120 cách\quad\quad \textbf{C.} 96 cách\quad\quad \textbf{D.} 80 cách\par\vspace{2pt}
\subsection*{\color{primary} Bảng đáp án Chương VIII: Đại số tổ hợp}
\begin{center}
\renewcommand{\arraystretch}{1.25}
\begin{tabularx}{\textwidth}{*{10}{>{\centering\arraybackslash}X}}
\toprule
\textbf{\color{primary} Cột 1} & \textbf{\color{primary} Cột 2} & \textbf{\color{primary} Cột 3} & \textbf{\color{primary} Cột 4} & \textbf{\color{primary} Cột 5} & \textbf{\color{primary} Cột 6} & \textbf{\color{primary} Cột 7} & \textbf{\color{primary} Cột 8} & \textbf{\color{primary} Cột 9} & \textbf{\color{primary} Cột 10} \\
\midrule
\textbf{1.A} & \textbf{2.A} & \textbf{3.C} & \textbf{4.D} & \textbf{5.D} & \textbf{6.B} & \textbf{7.D} & \textbf{8.C} & \textbf{9.C} & \textbf{10.C} \\
\rowcolor{rowtint} \textbf{11.C} & \textbf{12.C} & \textbf{13.D} & \textbf{14.B} & \textbf{15.C} & \textbf{16.B} & \textbf{17.C} & \textbf{18.A} & \textbf{19.C} & \textbf{20.A} \\
\textbf{21.C} & \textbf{22.C} & \textbf{23.B} & \textbf{24.C} & \textbf{25.A} & \textbf{26.D} & \textbf{27.D} & \textbf{28.A} & \textbf{29.B} & \textbf{30.C} \\
\rowcolor{rowtint} \textbf{31.D} & \textbf{32.C} & \textbf{33.B} & \textbf{34.A} & \textbf{35.A} & \textbf{36.C} & \textbf{37.A} & \textbf{38.B} & \textbf{39.B} & \textbf{40.D} \\
\textbf{41.B} & \textbf{42.A} & \textbf{43.B} & \textbf{44.A} & \textbf{45.B} & \textbf{46.D} & \textbf{47.B} & \textbf{48.B} & \textbf{49.C} & \textbf{50.A} \\
\bottomrule
\end{tabularx}
\end{center}\vspace{10pt}
\section*{\color{primary} CHƯƠNG IX: TÍNH XÁC SUẤT THEO ĐỊNH NGHĨA CỔ ĐIỂN}
\textit{Trọng tâm kiến thức SGK: Biến cố và định nghĩa cổ điển của xác suất, biến cố đối và thực hành tính xác suất sử dụng hoán vị, chỉnh hợp, tổ hợp.}\par\vspace{4pt}
\subsection*{\color{accent} Phần I. Mức độ Dễ (Nhận biết -- Câu 1 đến Câu 20)}
\textbf{\color{primary} Câu 1.} Cho biến cố A có xác suất P(A) = 2/7. Xác suất của biến cố đối $\bar{A}$ là:\par
\hspace*{1.5em}\textbf{A.} P($\bar{A}$) = 1/7\quad\quad \textbf{B.} P($\bar{A}$) = 2/7\quad\quad \textbf{C.} P($\bar{A}$) = 5/7\quad\quad \textbf{D.} P($\bar{A}$) = 7/2\par\vspace{2pt}
\textbf{\color{primary} Câu 2.} Thực hiện yêu cầu nhận biết sau: Gieo một đồng xu cân đối 3 lần liên tiếp. Số phần tử của không gian mẫu n($\Omega$) bằng.\par
\hspace*{1.5em}\textbf{A.} 6\quad\quad \textbf{B.} 4\quad\quad \textbf{C.} 9\quad\quad \textbf{D.} 8\par\vspace{2pt}
\textbf{\color{primary} Câu 3.} Trong bài học thuộc Chương IX, xét nội dung: Một hộp có 5 bi đỏ và 5 bi xanh. Lấy ngẫu nhiên 1 viên bi. Xác suất lấy được bi đỏ là. Đáp án đúng là:\par
\hspace*{1.5em}\textbf{A.} 1/5\quad\quad \textbf{B.} 5/10\quad\quad \textbf{C.} 5/5\quad\quad \textbf{D.} 5/10 (1)\par\vspace{2pt}
\textbf{\color{primary} Câu 4.} Dựa vào kiến thức chuẩn SGK của Chương IX (Tính xác suất theo định nghĩa cổ điển), cho biến cố A có xác suất P(A) = 5/13. Xác suất của biến cố đối $\bar{A}$ là:\par
\hspace*{1.5em}\textbf{A.} P($\bar{A}$) = 5/13\quad\quad \textbf{B.} P($\bar{A}$) = 8/13\quad\quad \textbf{C.} P($\bar{A}$) = 1/13\quad\quad \textbf{D.} P($\bar{A}$) = 13/5\par\vspace{2pt}
\textbf{\color{primary} Câu 5.} Chọn kết quả đúng nhất khi xét: Gieo một đồng xu cân đối 3 lần liên tiếp. Số phần tử của không gian mẫu n($\Omega$) bằng:\par
\hspace*{1.5em}\textbf{A.} 6\quad\quad \textbf{B.} 9\quad\quad \textbf{C.} 4\quad\quad \textbf{D.} 8\par\vspace{2pt}
\textbf{\color{primary} Câu 6.} Khi làm bài tập cơ bản về tính xác suất theo định nghĩa cổ điển, một học sinh gặp câu hỏi: Một hộp có 5 bi đỏ và 8 bi xanh. Lấy ngẫu nhiên 1 viên bi. Xác suất lấy được bi đỏ là. Kết quả nào dưới đây đúng?\par
\hspace*{1.5em}\textbf{A.} 1/5\quad\quad \textbf{B.} 5/8\quad\quad \textbf{C.} 5/13\quad\quad \textbf{D.} 8/13\par\vspace{2pt}
\textbf{\color{primary} Câu 7.} Xác định giá trị hoặc khẳng định đúng cho bài toán sau: Cho biến cố A có xác suất P(A) = 8/19. Xác suất của biến cố đối $\bar{A}$ là:\par
\hspace*{1.5em}\textbf{A.} P($\bar{A}$) = 11/19\quad\quad \textbf{B.} P($\bar{A}$) = 1/19\quad\quad \textbf{C.} P($\bar{A}$) = 19/8\quad\quad \textbf{D.} P($\bar{A}$) = 8/19\par\vspace{2pt}
\textbf{\color{primary} Câu 8.} Áp dụng trực tiếp định nghĩa và quy tắc trong Chương IX: Gieo một đồng xu cân đối 3 lần liên tiếp. Số phần tử của không gian mẫu n($\Omega$) bằng:\par
\hspace*{1.5em}\textbf{A.} 6\quad\quad \textbf{B.} 9\quad\quad \textbf{C.} 4\quad\quad \textbf{D.} 8\par\vspace{2pt}
\textbf{\color{primary} Câu 9.} Kiểm tra nhanh kiến thức nền tảng: Một hộp có 5 bi đỏ và 11 bi xanh. Lấy ngẫu nhiên 1 viên bi. Xác suất lấy được bi đỏ là:\par
\hspace*{1.5em}\textbf{A.} 5/16\quad\quad \textbf{B.} 1/5\quad\quad \textbf{C.} 11/16\quad\quad \textbf{D.} 5/11\par\vspace{2pt}
\textbf{\color{primary} Câu 10.} Cho dữ kiện nhận biết sau đây: Cho biến cố A có xác suất P(A) = 11/25. Xác suất của biến cố đối $\bar{A}$ là. Phương án chính xác là:\par
\hspace*{1.5em}\textbf{A.} P($\bar{A}$) = 14/25\quad\quad \textbf{B.} P($\bar{A}$) = 11/25\quad\quad \textbf{C.} P($\bar{A}$) = 25/11\quad\quad \textbf{D.} P($\bar{A}$) = 1/25\par\vspace{2pt}
\textbf{\color{primary} Câu 11.} Hãy chỉ ra kết quả đúng trong bốn lựa chọn A, B, C, D đối với yêu cầu: Gieo một đồng xu cân đối 3 lần liên tiếp. Số phần tử của không gian mẫu n($\Omega$) bằng:\par
\hspace*{1.5em}\textbf{A.} 4\quad\quad \textbf{B.} 6\quad\quad \textbf{C.} 8\quad\quad \textbf{D.} 9\par\vspace{2pt}
\textbf{\color{primary} Câu 12.} Câu hỏi ôn tập lý thuyết và kỹ năng cơ bản (Chương IX): Một hộp có 5 bi đỏ và 14 bi xanh. Lấy ngẫu nhiên 1 viên bi. Xác suất lấy được bi đỏ là:\par
\hspace*{1.5em}\textbf{A.} 5/19\quad\quad \textbf{B.} 5/14\quad\quad \textbf{C.} 1/5\quad\quad \textbf{D.} 14/19\par\vspace{2pt}
\textbf{\color{primary} Câu 13.} Theo nội dung sách giáo khoa Kết nối tri thức với cuộc sống (Chương IX), cho biến cố A có xác suất P(A) = 14/31. Xác suất của biến cố đối $\bar{A}$ là:\par
\hspace*{1.5em}\textbf{A.} P($\bar{A}$) = 17/31\quad\quad \textbf{B.} P($\bar{A}$) = 14/31\quad\quad \textbf{C.} P($\bar{A}$) = 31/14\quad\quad \textbf{D.} P($\bar{A}$) = 1/31\par\vspace{2pt}
\textbf{\color{primary} Câu 14.} Tìm đáp số đúng của bài toán cơ bản sau: Gieo một đồng xu cân đối 3 lần liên tiếp. Số phần tử của không gian mẫu n($\Omega$) bằng:\par
\hspace*{1.5em}\textbf{A.} 8\quad\quad \textbf{B.} 4\quad\quad \textbf{C.} 9\quad\quad \textbf{D.} 6\par\vspace{2pt}
\textbf{\color{primary} Câu 15.} Xét đối tượng toán học trong Chương IX: Một hộp có 5 bi đỏ và 17 bi xanh. Lấy ngẫu nhiên 1 viên bi. Xác suất lấy được bi đỏ là. Giá trị hoặc tính chất tương ứng là:\par
\hspace*{1.5em}\textbf{A.} 1/5\quad\quad \textbf{B.} 5/22\quad\quad \textbf{C.} 17/22\quad\quad \textbf{D.} 5/17\par\vspace{2pt}
\textbf{\color{primary} Câu 16.} Hoàn thành phép xác định cơ bản sau đây: Cho biến cố A có xác suất P(A) = 17/37. Xác suất của biến cố đối $\bar{A}$ là:\par
\hspace*{1.5em}\textbf{A.} P($\bar{A}$) = 1/37\quad\quad \textbf{B.} P($\bar{A}$) = 20/37\quad\quad \textbf{C.} P($\bar{A}$) = 37/17\quad\quad \textbf{D.} P($\bar{A}$) = 17/37\par\vspace{2pt}
\textbf{\color{primary} Câu 17.} Trong các khẳng định hoặc giá trị dưới đây, đâu là đáp án đúng cho câu hỏi: Gieo một đồng xu cân đối 3 lần liên tiếp. Số phần tử của không gian mẫu n($\Omega$) bằng?\par
\hspace*{1.5em}\textbf{A.} 6\quad\quad \textbf{B.} 8\quad\quad \textbf{C.} 4\quad\quad \textbf{D.} 9\par\vspace{2pt}
\textbf{\color{primary} Câu 18.} Bài tập khởi động Chương IX: Một hộp có 5 bi đỏ và 20 bi xanh. Lấy ngẫu nhiên 1 viên bi. Xác suất lấy được bi đỏ là:\par
\hspace*{1.5em}\textbf{A.} 5/25\quad\quad \textbf{B.} 5/20\quad\quad \textbf{C.} 20/25\quad\quad \textbf{D.} 1/5\par\vspace{2pt}
\textbf{\color{primary} Câu 19.} Quan sát dữ kiện đề bài: Cho biến cố A có xác suất P(A) = 20/43. Xác suất của biến cố đối $\bar{A}$ là. Hãy chọn phương án đúng:\par
\hspace*{1.5em}\textbf{A.} P($\bar{A}$) = 23/43\quad\quad \textbf{B.} P($\bar{A}$) = 1/43\quad\quad \textbf{C.} P($\bar{A}$) = 20/43\quad\quad \textbf{D.} P($\bar{A}$) = 43/20\par\vspace{2pt}
\textbf{\color{primary} Câu 20.} Sử dụng công thức hoặc tính chất đầu bài của Chương IX để giải quyết yêu cầu: Gieo một đồng xu cân đối 3 lần liên tiếp. Số phần tử của không gian mẫu n($\Omega$) bằng:\par
\hspace*{1.5em}\textbf{A.} 4\quad\quad \textbf{B.} 8\quad\quad \textbf{C.} 6\quad\quad \textbf{D.} 9\par\vspace{2pt}
\subsection*{\color{accent} Phần II. Mức độ Trung bình (Thông hiểu -- Câu 21 đến Câu 40)}
\textbf{\color{primary} Câu 21.} Một hộp có 4 viên bi đỏ và 5 viên bi xanh. Lấy ngẫu nhiên đồng thời 2 viên bi. Xác suất để lấy được 2 viên bi cùng màu bằng:\par
\hspace*{1.5em}\textbf{A.} 2/9\quad\quad \textbf{B.} 5/9\quad\quad \textbf{C.} 1/3\quad\quad \textbf{D.} 4/9\par\vspace{2pt}
\textbf{\color{primary} Câu 22.} Vận dụng quy tắc trong Chương IX (Tính xác suất theo định nghĩa cổ điển) để giải bài toán: Gieo một đồng xu cân đối 4 lần liên tiếp. Xác suất để có ít nhất 1 lần xuất hiện mặt Ngửa bằng:\par
\hspace*{1.5em}\textbf{A.} 1/2\quad\quad \textbf{B.} 15/16\quad\quad \textbf{C.} 1/16\quad\quad \textbf{D.} 3/4\par\vspace{2pt}
\textbf{\color{primary} Câu 23.} Biến đổi và tính toán theo kiến thức Chương IX: Từ một tổ gồm 5 nam và 4 nữ, chọn ngẫu nhiên 2 người. Xác suất để cả 2 người được chọn đều là nữ bằng. Kết quả thu được là:\par
\hspace*{1.5em}\textbf{A.} 1/3\quad\quad \textbf{B.} 2/9\quad\quad \textbf{C.} 1/6\quad\quad \textbf{D.} 5/18\par\vspace{2pt}
\textbf{\color{primary} Câu 24.} Xét bài toán mức độ thông hiểu sau: Một hộp có 6 viên bi trắng và 2 viên bi đen. Lấy ngẫu nhiên đồng thời 2 viên bi từ hộp. Số phần tử của không gian mẫu n($\Omega$) = C\_\{n\}$^{2}$ và xác suất lấy được 2 viên bi đen lần lượt là. Đáp án đúng là:\par
\hspace*{1.5em}\textbf{A.} n($\Omega$) = 28; P = 1/28\quad\quad \textbf{B.} n($\Omega$) = 8; P = 2/8\quad\quad \textbf{C.} n($\Omega$) = 56; P = 1/56\quad\quad \textbf{D.} n($\Omega$) = 28; P = 2/28\par\vspace{2pt}
\textbf{\color{primary} Câu 25.} Để hoàn thành phiếu học tập Chương IX, hãy xác định: Một hộp có 4 viên bi đỏ và 5 viên bi xanh. Lấy ngẫu nhiên đồng thời 2 viên bi. Xác suất để lấy được 2 viên bi cùng màu bằng:\par
\hspace*{1.5em}\textbf{A.} 1/3\quad\quad \textbf{B.} 2/9\quad\quad \textbf{C.} 5/9\quad\quad \textbf{D.} 4/9\par\vspace{2pt}
\textbf{\color{primary} Câu 26.} Phân tích dữ kiện bài toán: Gieo một đồng xu cân đối 4 lần liên tiếp. Xác suất để có ít nhất 1 lần xuất hiện mặt Ngửa bằng. Giá trị cần tìm bằng:\par
\hspace*{1.5em}\textbf{A.} 1/16\quad\quad \textbf{B.} 1/2\quad\quad \textbf{C.} 3/4\quad\quad \textbf{D.} 15/16\par\vspace{2pt}
\textbf{\color{primary} Câu 27.} Áp dụng tính chất đã học ở Chương IX vào bài tập sau: Từ một tổ gồm 5 nam và 4 nữ, chọn ngẫu nhiên 2 người. Xác suất để cả 2 người được chọn đều là nữ bằng:\par
\hspace*{1.5em}\textbf{A.} 5/18\quad\quad \textbf{B.} 2/9\quad\quad \textbf{C.} 1/3\quad\quad \textbf{D.} 1/6\par\vspace{2pt}
\textbf{\color{primary} Câu 28.} Thực hiện các bước suy luận hợp lý để giải quyết yêu cầu: Một hộp có 10 viên bi trắng và 2 viên bi đen. Lấy ngẫu nhiên đồng thời 2 viên bi từ hộp. Số phần tử của không gian mẫu n($\Omega$) = C\_\{n\}$^{2}$ và xác suất lấy được 2 viên bi đen lần lượt là:\par
\hspace*{1.5em}\textbf{A.} n($\Omega$) = 12; P = 2/12\quad\quad \textbf{B.} n($\Omega$) = 132; P = 1/132\quad\quad \textbf{C.} n($\Omega$) = 66; P = 2/66\quad\quad \textbf{D.} n($\Omega$) = 66; P = 1/66\par\vspace{2pt}
\textbf{\color{primary} Câu 29.} Cho bài toán tự luyện (Tính xác suất theo định nghĩa cổ điển): Một hộp có 4 viên bi đỏ và 5 viên bi xanh. Lấy ngẫu nhiên đồng thời 2 viên bi. Xác suất để lấy được 2 viên bi cùng màu bằng. Phương án nào sau đây chính xác?\par
\hspace*{1.5em}\textbf{A.} 1/3\quad\quad \textbf{B.} 4/9\quad\quad \textbf{C.} 5/9\quad\quad \textbf{D.} 2/9\par\vspace{2pt}
\textbf{\color{primary} Câu 30.} Kết hợp các phép biến đổi trong Chương IX, hãy tìm kết quả của bài toán: Gieo một đồng xu cân đối 4 lần liên tiếp. Xác suất để có ít nhất 1 lần xuất hiện mặt Ngửa bằng:\par
\hspace*{1.5em}\textbf{A.} 3/4\quad\quad \textbf{B.} 15/16\quad\quad \textbf{C.} 1/16\quad\quad \textbf{D.} 1/2\par\vspace{2pt}
\textbf{\color{primary} Câu 31.} Một bài kiểm tra định kỳ nội dung Chương IX nêu câu hỏi: Từ một tổ gồm 5 nam và 4 nữ, chọn ngẫu nhiên 2 người. Xác suất để cả 2 người được chọn đều là nữ bằng. Hãy chọn đáp số đúng:\par
\hspace*{1.5em}\textbf{A.} 5/18\quad\quad \textbf{B.} 1/6\quad\quad \textbf{C.} 1/3\quad\quad \textbf{D.} 2/9\par\vspace{2pt}
\textbf{\color{primary} Câu 32.} Tìm lời giải đúng cho bài toán thông hiểu: Một hộp có 14 viên bi trắng và 2 viên bi đen. Lấy ngẫu nhiên đồng thời 2 viên bi từ hộp. Số phần tử của không gian mẫu n($\Omega$) = C\_\{n\}$^{2}$ và xác suất lấy được 2 viên bi đen lần lượt là:\par
\hspace*{1.5em}\textbf{A.} n($\Omega$) = 120; P = 1/120\quad\quad \textbf{B.} n($\Omega$) = 240; P = 1/240\quad\quad \textbf{C.} n($\Omega$) = 16; P = 2/16\quad\quad \textbf{D.} n($\Omega$) = 120; P = 2/120\par\vspace{2pt}
\textbf{\color{primary} Câu 33.} Sử dụng mối liên hệ giữa các đại lượng trong Chương IX: Một hộp có 4 viên bi đỏ và 5 viên bi xanh. Lấy ngẫu nhiên đồng thời 2 viên bi. Xác suất để lấy được 2 viên bi cùng màu bằng:\par
\hspace*{1.5em}\textbf{A.} 5/9\quad\quad \textbf{B.} 2/9\quad\quad \textbf{C.} 4/9\quad\quad \textbf{D.} 1/3\par\vspace{2pt}
\textbf{\color{primary} Câu 34.} Hãy tính toán và chọn phương án đúng đối với bài tập: Gieo một đồng xu cân đối 4 lần liên tiếp. Xác suất để có ít nhất 1 lần xuất hiện mặt Ngửa bằng:\par
\hspace*{1.5em}\textbf{A.} 3/4\quad\quad \textbf{B.} 1/2\quad\quad \textbf{C.} 1/16\quad\quad \textbf{D.} 15/16\par\vspace{2pt}
\textbf{\color{primary} Câu 35.} Dựa trên các tính chất trọng tâm của Chương IX, xét bài toán: Từ một tổ gồm 5 nam và 4 nữ, chọn ngẫu nhiên 2 người. Xác suất để cả 2 người được chọn đều là nữ bằng:\par
\hspace*{1.5em}\textbf{A.} 1/3\quad\quad \textbf{B.} 2/9\quad\quad \textbf{C.} 1/6\quad\quad \textbf{D.} 5/18\par\vspace{2pt}
\textbf{\color{primary} Câu 36.} Thực hành giải toán (Tính xác suất theo định nghĩa cổ điển): Một hộp có 18 viên bi trắng và 2 viên bi đen. Lấy ngẫu nhiên đồng thời 2 viên bi từ hộp. Số phần tử của không gian mẫu n($\Omega$) = C\_\{n\}$^{2}$ và xác suất lấy được 2 viên bi đen lần lượt là. Đáp án nào dưới đây thỏa mãn?\par
\hspace*{1.5em}\textbf{A.} n($\Omega$) = 190; P = 2/190\quad\quad \textbf{B.} n($\Omega$) = 190; P = 1/190\quad\quad \textbf{C.} n($\Omega$) = 20; P = 2/20\quad\quad \textbf{D.} n($\Omega$) = 380; P = 1/380\par\vspace{2pt}
\textbf{\color{primary} Câu 37.} Kiểm tra kỹ năng biến đổi và suy luận qua bài toán: Một hộp có 4 viên bi đỏ và 5 viên bi xanh. Lấy ngẫu nhiên đồng thời 2 viên bi. Xác suất để lấy được 2 viên bi cùng màu bằng:\par
\hspace*{1.5em}\textbf{A.} 4/9\quad\quad \textbf{B.} 5/9\quad\quad \textbf{C.} 1/3\quad\quad \textbf{D.} 2/9\par\vspace{2pt}
\textbf{\color{primary} Câu 38.} Cho giả thiết thuộc phạm vi Chương IX: Gieo một đồng xu cân đối 4 lần liên tiếp. Xác suất để có ít nhất 1 lần xuất hiện mặt Ngửa bằng. Kết luận đúng là:\par
\hspace*{1.5em}\textbf{A.} 1/16\quad\quad \textbf{B.} 3/4\quad\quad \textbf{C.} 1/2\quad\quad \textbf{D.} 15/16\par\vspace{2pt}
\textbf{\color{primary} Câu 39.} Rút gọn hoặc tính giá trị theo yêu cầu sau của Chương IX: Từ một tổ gồm 5 nam và 4 nữ, chọn ngẫu nhiên 2 người. Xác suất để cả 2 người được chọn đều là nữ bằng:\par
\hspace*{1.5em}\textbf{A.} 1/6\quad\quad \textbf{B.} 2/9\quad\quad \textbf{C.} 1/3\quad\quad \textbf{D.} 5/18\par\vspace{2pt}
\textbf{\color{primary} Câu 40.} Xác định phương án đúng sau khi thực hiện phép tính / suy luận: Một hộp có 22 viên bi trắng và 2 viên bi đen. Lấy ngẫu nhiên đồng thời 2 viên bi từ hộp. Số phần tử của không gian mẫu n($\Omega$) = C\_\{n\}$^{2}$ và xác suất lấy được 2 viên bi đen lần lượt là:\par
\hspace*{1.5em}\textbf{A.} n($\Omega$) = 552; P = 1/552\quad\quad \textbf{B.} n($\Omega$) = 24; P = 2/24\quad\quad \textbf{C.} n($\Omega$) = 276; P = 1/276\quad\quad \textbf{D.} n($\Omega$) = 276; P = 2/276\par\vspace{2pt}
\subsection*{\color{accent} Phần III. Mức độ Khó (Vận dụng \& Vận dụng cao -- Câu 41 đến Câu 50)}
\textbf{\color{primary} Câu 41.} Xếp ngẫu nhiên 5 học sinh gồm A, B, C, D, E thành một hàng ngang. Xác suất để hai bạn A và B luôn đứng cạnh nhau bằng:\par
\hspace*{1.5em}\textbf{A.} 1/5\quad\quad \textbf{B.} 2/5\quad\quad \textbf{C.} 3/5\quad\quad \textbf{D.} 1/2\par\vspace{2pt}
\textbf{\color{primary} Câu 42.} Bài toán vận dụng (Chương IX – Tính xác suất theo định nghĩa cổ điển): Một hộp có 6 quả cầu trắng và 4 quả cầu đen. Lấy ngẫu nhiên đồng thời 3 quả cầu. Xác suất để trong 3 quả cầu lấy ra có ít nhất 1 quả cầu đen bằng. Kết quả đúng là:\par
\hspace*{1.5em}\textbf{A.} 5/6\quad\quad \textbf{B.} 3/4\quad\quad \textbf{C.} 2/3\quad\quad \textbf{D.} 1/6\par\vspace{2pt}
\textbf{\color{primary} Câu 43.} Vận dụng kiến thức nâng cao trong phạm vi Chương IX để giải bài toán: Xếp ngẫu nhiên 6 quyển sách khác nhau (trong đó có 1 quyển Toán và 1 quyển Vật lí) thành một hàng ngang trên giá sách. Xác suất để hai quyển Toán và Vật lí đứng cạnh nhau là:\par
\hspace*{1.5em}\textbf{A.} 2/6\quad\quad \textbf{B.} 1/6\quad\quad \textbf{C.} 4/6\quad\quad \textbf{D.} 2/30\par\vspace{2pt}
\textbf{\color{primary} Câu 44.} Phân tích và giải bài toán tổng hợp sau của Chương IX: Xếp ngẫu nhiên 5 học sinh gồm A, B, C, D, E thành một hàng ngang. Xác suất để hai bạn A và B luôn đứng cạnh nhau bằng:\par
\hspace*{1.5em}\textbf{A.} 2/5\quad\quad \textbf{B.} 1/5\quad\quad \textbf{C.} 3/5\quad\quad \textbf{D.} 1/2\par\vspace{2pt}
\textbf{\color{primary} Câu 45.} Thử thách tư duy trong phạm vi kiến thức Chương IX: Một hộp có 6 quả cầu trắng và 4 quả cầu đen. Lấy ngẫu nhiên đồng thời 3 quả cầu. Xác suất để trong 3 quả cầu lấy ra có ít nhất 1 quả cầu đen bằng. Đáp số chính xác là:\par
\hspace*{1.5em}\textbf{A.} 1/6\quad\quad \textbf{B.} 2/3\quad\quad \textbf{C.} 5/6\quad\quad \textbf{D.} 3/4\par\vspace{2pt}
\textbf{\color{primary} Câu 46.} Sử dụng phương pháp giải toán của Chương IX để tìm đáp án cho bài toán: Xếp ngẫu nhiên 9 quyển sách khác nhau (trong đó có 1 quyển Toán và 1 quyển Vật lí) thành một hàng ngang trên giá sách. Xác suất để hai quyển Toán và Vật lí đứng cạnh nhau là:\par
\hspace*{1.5em}\textbf{A.} 2/9\quad\quad \textbf{B.} 2/72\quad\quad \textbf{C.} 1/9\quad\quad \textbf{D.} 7/9\par\vspace{2pt}
\textbf{\color{primary} Câu 47.} Xét bài toán vận dụng thực tiễn / nâng cao thuộc Chương IX: Xếp ngẫu nhiên 5 học sinh gồm A, B, C, D, E thành một hàng ngang. Xác suất để hai bạn A và B luôn đứng cạnh nhau bằng:\par
\hspace*{1.5em}\textbf{A.} 1/2\quad\quad \textbf{B.} 3/5\quad\quad \textbf{C.} 1/5\quad\quad \textbf{D.} 2/5\par\vspace{2pt}
\textbf{\color{primary} Câu 48.} Tìm kết quả cuối cùng của bài toán phân loại học sinh khá giỏi (Chương IX): Một hộp có 6 quả cầu trắng và 4 quả cầu đen. Lấy ngẫu nhiên đồng thời 3 quả cầu. Xác suất để trong 3 quả cầu lấy ra có ít nhất 1 quả cầu đen bằng:\par
\hspace*{1.5em}\textbf{A.} 2/3\quad\quad \textbf{B.} 3/4\quad\quad \textbf{C.} 5/6\quad\quad \textbf{D.} 1/6\par\vspace{2pt}
\textbf{\color{primary} Câu 49.} Kết hợp các tính chất đã học trong Chương IX, hãy giải quyết bài toán: Xếp ngẫu nhiên 12 quyển sách khác nhau (trong đó có 1 quyển Toán và 1 quyển Vật lí) thành một hàng ngang trên giá sách. Xác suất để hai quyển Toán và Vật lí đứng cạnh nhau là:\par
\hspace*{1.5em}\textbf{A.} 2/12\quad\quad \textbf{B.} 1/12\quad\quad \textbf{C.} 10/12\quad\quad \textbf{D.} 2/132\par\vspace{2pt}
\textbf{\color{primary} Câu 50.} Câu hỏi vận dụng cao cuối Chương IX: Xếp ngẫu nhiên 5 học sinh gồm A, B, C, D, E thành một hàng ngang. Xác suất để hai bạn A và B luôn đứng cạnh nhau bằng. Phương án đúng là:\par
\hspace*{1.5em}\textbf{A.} 1/5\quad\quad \textbf{B.} 1/2\quad\quad \textbf{C.} 3/5\quad\quad \textbf{D.} 2/5\par\vspace{2pt}
\subsection*{\color{primary} Bảng đáp án Chương IX: Tính xác suất theo định nghĩa cổ điển}
\begin{center}
\renewcommand{\arraystretch}{1.25}
\begin{tabularx}{\textwidth}{*{10}{>{\centering\arraybackslash}X}}
\toprule
\textbf{\color{primary} Cột 1} & \textbf{\color{primary} Cột 2} & \textbf{\color{primary} Cột 3} & \textbf{\color{primary} Cột 4} & \textbf{\color{primary} Cột 5} & \textbf{\color{primary} Cột 6} & \textbf{\color{primary} Cột 7} & \textbf{\color{primary} Cột 8} & \textbf{\color{primary} Cột 9} & \textbf{\color{primary} Cột 10} \\
\midrule
\textbf{1.C} & \textbf{2.D} & \textbf{3.B} & \textbf{4.B} & \textbf{5.D} & \textbf{6.C} & \textbf{7.A} & \textbf{8.D} & \textbf{9.A} & \textbf{10.A} \\
\rowcolor{rowtint} \textbf{11.C} & \textbf{12.A} & \textbf{13.A} & \textbf{14.A} & \textbf{15.B} & \textbf{16.B} & \textbf{17.B} & \textbf{18.A} & \textbf{19.A} & \textbf{20.B} \\
\textbf{21.D} & \textbf{22.B} & \textbf{23.C} & \textbf{24.A} & \textbf{25.D} & \textbf{26.D} & \textbf{27.D} & \textbf{28.D} & \textbf{29.B} & \textbf{30.B} \\
\rowcolor{rowtint} \textbf{31.B} & \textbf{32.A} & \textbf{33.C} & \textbf{34.D} & \textbf{35.C} & \textbf{36.B} & \textbf{37.A} & \textbf{38.D} & \textbf{39.A} & \textbf{40.C} \\
\textbf{41.B} & \textbf{42.A} & \textbf{43.A} & \textbf{44.A} & \textbf{45.C} & \textbf{46.A} & \textbf{47.D} & \textbf{48.C} & \textbf{49.A} & \textbf{50.D} \\
\bottomrule
\end{tabularx}
\end{center}\vspace{10pt}
\end{document}
